\section{The Galois link I: admissibility at inert places}\label{sec:pf:linkN-adm}

\providecommand{\linkNadmcol}{\operatorname{col}}
\providecommand{\linkNadmZf}{Z^1_{f}}
\providecommand{\linkNadmHf}{H^1_{f}}
\providecommand{\linkNadmIm}{\operatorname{Im}_f}
\providecommand{\linkNadmlen}{\operatorname{len}}
\providecommand{\linkNadmerr}{\operatorname{err}}
\providecommand{\linkNadmvcc}{v_{cc}}
\providecommand{\linkNadmsP}{s_{\Pi}}
\providecommand{\linkNadmVV}{\mathrm{VV}}
\providecommand{\linkNadmlog}{\log}

\subsection{Introduction}\label{pf:linkN-adm:intro}

This section proves the two inequalities called \emph{mid} and \emph{bb} in the table of local bounds of
the Galois link (\cref{thm:linkN}, step~2 of the outline of its proof, and the local bound
\eqref{eq:families:bd}). Let $D$ be a datum with size data $L$, let $\iota$ be a twisted local index of $D$ whose
prime $\ell=\ell_\iota$ is not one of the two type primes and lies in the \emph{middle range} (it is neither
small nor $3$, nor one of the primes $q,d$ of the host, and it is at most the threshold $\ell^*(L)$ unless it
belongs to the exceptional set $\mathrm{exc}$ of the hub data). Then for every good eigensystem $\pi$ at $\iota$
\begin{align}
\linkNadmvcc(\iota,\pi)&\le \linkNadmsP(\iota,\rho_\iota,\pi)+\sum_{\lambda\mid\ell}\linkNadmlen_0(\iota,\pi,\lambda)
+\linkNadmVV(\ell)+v_\ell(h_E)+\linkNadmerr(\iota,\pi)+\ell^{-2},\label{pf:linkN-adm:eq-mid}\\
\linkNadmerr(\iota,\pi)&\le 7200\,f\,n\,t^2\,n_V\,\Lambda^2,\label{pf:linkN-adm:eq-bb}
\end{align}
where $\linkNadmlen_0$ and $\linkNadmerr$ are the local lengths and the error term recorded in the \emph{hub data}
attached to $D$, $\Lambda$ is the logarithmic scale of $L$, and $f,n,t,n_V$ are ranks of $D$
(\cref{pf:linkN-adm:main}). The hypotheses are the package $\sfname{Hypotheses}_2$ of results from the
literature (\cref{sec:families:galois}), a setup condition on the hub data, a lifting corollary with constant $C_2$
(``Corollary A''), and the \emph{joint hub conditions}, which describe the Hodge--Tate weights of the local
Galois representations at the places above $\ell$ by \emph{words}.

\paragraph{Idea of the proof.}
The inequality \eqref{pf:linkN-adm:eq-mid} is the conclusion of an \emph{integrality-from-metric theorem}
proved elsewhere (\lref{GlobalTheory.T.integralityFromMetric_of_designConditionAll}) once a list of local
conditions, collectively called \emph{admissibility}, holds at $(\iota,\pi)$; the error term $\linkNadmerr$ is
the part of its right-hand side that does not come from the places above $\ell$. At primes split in the host
field the admissibility conditions are collected from the setup. At inert primes the new conditions are the
\emph{row conditions}: at every place $\lambda\mid\ell$ of the CM field $E$, every surjective map $F\colon T_1\to T_2$
between suitable crystalline lattices of the decomposition group has a cokernel on Bloch--Kato finite cohomology
of length at most $\linkNadmlen_0(\iota,\pi,\lambda)$ times the absolute ramification index of the coefficients.
These are local statements in $\ell$-adic Hodge theory. Depending on the word attached to $T_1$ they come from
three local lifting theorems:
\begin{itemize}
\item \emph{long Frobenius orbits} (row~B): Corollary~A, a hypothesis of the main theorem;
\item \emph{a word with a non-thick letter} (row~C): the local lifting theorem~B (\cref{pf:linkN-adm:thmB});
\item \emph{a word all of whose letters are thick} (row~D): the theorem on ramified coefficient chains,
  theorem~C (\cref{pf:linkN-adm:thmC}).
\end{itemize}
Theorems~B and~C are proved from Bloch--Kato's local theory \cite{BlochKato1990}, local Tate duality and the
Euler characteristic formula \cite{Tate1963,Serre1997}, through Wach-type $(\varphi,\Gamma)$-modules
\cite{Wach1996,Berger2004} and Herr cocycles \cite{Herr1998}; their common core is a \emph{counting principle}:
the length of the cokernel is bounded by (number of generators)$\times$(an exponent killing them), and the
exponents are bounded explicitly in terms of the Hodge--Tate data. The inequality \eqref{pf:linkN-adm:eq-bb} is
an explicit estimate of the error term.

\paragraph{Organisation.}
\Cref{pf:linkN-adm:local} fixes the local notions (lattices, finite cocycles, the cokernel length $\linkNadmcol$)
and proves two counting lemmas. \Cref{pf:linkN-adm:global} introduces the global notation, the hub data and the
two inequalities, and states the main theorem; its proof is given in \cref{pf:linkN-adm:proofmain} from the
star inequality (\cref{pf:linkN-adm:star}), the bound for the error term (\cref{pf:linkN-adm:bbmain}), the split step
(\cref{pf:linkN-adm:split}) and the inert step (\cref{pf:linkN-adm:inert}). The inert step uses the row bounds
of \cref{pf:linkN-adm:rows}, which rest on theorems~B and~C, proved in \cref{pf:linkN-adm:thmBsec} and
\cref{pf:linkN-adm:thmCsec}; the admissibility of their local inputs is proved in \cref{pf:linkN-adm:admsec}.
\Cref{pf:linkN-adm:remarks} collects remarks on the formal statements. the part of the formal proof treated here
comprises about $69\,500$ Lean lines \cite{LeanSource2026}. Its deepest parts (the arc lifting theorems for arcs of length
two and three, the relay, splitting and universal statements for Wach modules, and the local lemmas of the ramified
coefficient chains, together about $55\,000$ lines) are reduced to precisely stated results of the formal development,
which are cited and not reproved.

% ---------------------------------------------------------------------------
\subsection{Local lattices and the cokernel length}\label{pf:linkN-adm:local}

\paragraph{Period data.}
Fix a prime $p$ and finite extensions $K_0\subseteq K$ of $\Q_p$, $K_0$ unramified. The development axiomatizes
the period rings: a \emph{period datum} $\mathrm{PD}$ for $G=G_K$ consists of a $K_0$-algebra $B$ with a
$G$-action, an injective Frobenius $\varphi$ semilinear over the Frobenius $\sigma$ of $K_0$ and commuting with
$G$, a filtered $K$-algebra $B_{\mathrm{dR}}$ with a $G$-action and a $G$-equivariant injection $B\to B_{\mathrm{dR}}$,
a character $\chi\colon G\to\Q_p^\times$ and an element $t\in B^\times$ with $g(t)=\chi(g)t$, $\varphi(t)=pt$,
$t\in\mathrm{Fil}^1\setminus\mathrm{Fil}^2$, such that $B^G=K_0$, $B_{\mathrm{dR}}^G=K$, and the fundamental exact
sequence holds in the form $\{b\in B:\varphi(b)=b,\ b\in\mathrm{Fil}^0\}=\Q_p$. These are the properties of
Fontaine's rings $B_{\mathrm{cris}}$, $B_{\mathrm{dR}}$ \cite{Fontaine1982,Fontaine1994}, and the instances used
below are Fontaine's rings (\cref{sec:families:galois}). Let $L_C$ be a finite extension of $\Q_p$ (the
\emph{coefficient field}) with ring of integers $\OO_C$ and $V$ a finite-dimensional continuous $L_C$-linear
representation of $G$. Then $D_{\mathrm{cris}}(V)=(V\otimes_{\Q_p}B)^G$; $V$ is \emph{crystalline} if
$\dim_{L_C}D_{\mathrm{cris}}(V)=[K_0:\Q_p]\dim_{L_C}V$. For an embedding $\tau\colon K\to L_C$ and $a\in\Z$ let
$\mathrm{Fil}^a_\tau(V)$ be the $\tau$-part of $(V\otimes B_{\mathrm{dR}})^G\cap(V\otimes\mathrm{Fil}^a)$; $a$ is a
\emph{jump} (a Hodge--Tate weight) of $V$ at $\tau$ if $\mathrm{Fil}^{a+1}_\tau\ne\mathrm{Fil}^a_\tau$, and
$\mathrm{mult}_\tau(V,a)=\dim\mathrm{Fil}^a_\tau-\dim\mathrm{Fil}^{a+1}_\tau$ is its multiplicity.

\begin{definition}[Lattices and the cokernel length]\label{pf:linkN-adm:def-col}
A \emph{lattice} of $V$ is a $G$-stable finitely generated free $\OO_C$-submodule $T$ spanning $V$; a \emph{lattice
map} $F\colon T_1\to T_2$ is a continuous $G$-equivariant $L_C$-linear map $V_1\to V_2$ with $F(T_1)\subseteq T_2$;
it is \emph{onto} if $F(T_1)=T_2$. Let $Z^1(T)$ be the continuous $1$-cocycles $G\to V$ with values in $T$,
$B^1(T)=\{g\mapsto gx-x:x\in T\}$, and $\linkNadmZf(T)\subseteq Z^1(T)$ the \emph{finite} cocycles: those $c$ for
which there is $y\in V\otimes B$ with $c(g)\otimes1=gy-y$ for all $g$ (the Bloch--Kato condition
\cite{BlochKato1990}). Put $H^1(T)=Z^1(T)/B^1(T)$ and let $\linkNadmHf(T)$ be the image of $\linkNadmZf(T)$. For
a lattice map $F$ put
\[
\linkNadmIm(F)=F_*\linkNadmZf(T_1)+B^1(T_2),\qquad
\linkNadmcol(F)=\linkNadmlen_{\OO_C}\bigl(\linkNadmZf(T_2)/(\linkNadmIm(F)\cap\linkNadmZf(T_2))\bigr)\in\N\cup\{\infty\}.
\]
For $x\in\R$ we write $\linkNadmcol(F)\preceq x$ if $\linkNadmcol(F)$ is finite and
$\linkNadmcol(F)\le x\cdot e_C$, where $e_C=\linkNadmlen_{\OO_C}(\OO_C/p)$ is the absolute ramification index of
$L_C$. For $e\in\N$, $F$ \emph{lifts with exponent} $e$ if $p^e\linkNadmZf(T_2)\subseteq\linkNadmIm(F)$.
\formal{Formalism.WeaklyAdmissibleFiltrations.LatHom.col}
\formal{Formalism.MainNotions.TruncatedSeriesCountingBounds.ColLeR}
\formal{Formalism.WeaklyAdmissibleFiltrations.LatHom.LiftF}
\end{definition}

Thus $\linkNadmcol(F)$ is the length of the cokernel of $F_*\colon\linkNadmHf(T_1)\to\linkNadmHf(T_2)$, computed
on cocycles. If $\linkNadmcol(F)\preceq x$ and $x\le y$ then $\linkNadmcol(F)\preceq y$; we use this monotonicity
without comment. We use the following standard facts on modules over a discrete valuation ring $\OO$ (see e.g.\
\cite{Matsumura1989}): length is additive in short exact sequences; $\linkNadmlen(\OO/a)$ is finite for $a\ne0$; a
module generated by $g$ elements and killed by $a^e$ is a quotient of $(\OO/a^e)^g$ and so has length at most
$g\,e\,\linkNadmlen(\OO/a)$; and a submodule of a module generated by $g$ elements is generated by at most $g$ elements.

\begin{lemma}[Two-step count]\label{pf:linkN-adm:twostep}
Let $\OO$ be a discrete valuation ring, $a\in\OO\setminus\{0\}$, $M$ an $\OO$-module, $N\subseteq M$ a submodule,
$s,Z\subseteq M$ finite sets with $M=\OO s+N$, and $e,e'\in\N$ such that $a^eM\subseteq\OO Z+N$ and
$a^{e'}Z\subseteq N$. Then $\linkNadmlen(M/N)\le(|s|e+|Z|e')\linkNadmlen(\OO/a)$.
\end{lemma}

\begin{proof}
This is \lref{Formalism.TruncatedSeriesCounting.TruncatedSeriesCountingBounds.length_quotient_le_of_two_step}.
The module $M/(\OO Z+N)$ is generated by the images of $s$ and killed by $a^e$, and $(\OO Z+N)/N$ is generated by
the images of $Z$ and killed by $a^{e'}$; add the two lengths.
\end{proof}

\begin{lemma}[Counting principle]\label{pf:linkN-adm:colexp}
Let $\OO_C$ be a discrete valuation ring with $p\ne0$ in $\OO_C$, $F\colon T_1\to T_2$ a lattice map, $g_1,e,e'\in\N$
and $Z\subseteq\linkNadmZf(T_2)$ a finite set. Assume that $\linkNadmHf(T_2)$ is generated by at most $g_1$
elements, that $p^e\linkNadmZf(T_2)\subseteq\linkNadmIm(F)+\OO_CZ$, and that $F$
lifts with exponent $e'$. Then $\linkNadmcol(F)\preceq g_1e+|Z|e'$.
\formallabel{Formalism.MainNotions.N.Ln17.colLeR_of_exponent}
\end{lemma}

\begin{proof}
Choose a set $s$ of at most $g_1$ finite cocycles whose classes generate $\linkNadmHf(T_2)$, and put
$M=\linkNadmZf(T_2)$, $N=\linkNadmIm(F)\cap M$. Since $B^1(T_2)\subseteq\linkNadmIm(F)$, every $c\in M$ differs
from an $\OO_C$-combination of $s$ by a coboundary lying in $N$; hence $M=\OO_Cs+N$. If $c\in M$, write
$p^ec=i+z$ with $i\in\linkNadmIm(F)$ and $z\in\OO_CZ\subseteq M$; then $i\in N$, so $p^eM\subseteq\OO_CZ+N$.
Finally $p^{e'}Z\subseteq p^{e'}M\subseteq N$ because $F$ lifts with exponent $e'$. \Cref{pf:linkN-adm:twostep}
with $a=p$ gives $\linkNadmcol(F)\le(g_1e+|Z|e')\,e_C$, which is finite since $p\ne0$.
\end{proof}

\begin{lemma}[Filtration count]\label{pf:linkN-adm:filtration}
Let $\OO$ be a discrete valuation ring, $a\in\OO\setminus\{0\}$, $k_0=\linkNadmlen(\OO/a)$, $M$ a free $\OO$-module of
finite rank, $P\subseteq M$ a submodule, $M=\mathrm{Fl}_0\supseteq\mathrm{Fl}_1\supseteq\dots\supseteq\mathrm{Fl}_k$ a
chain of submodules and $\mathrm{gen}\subseteq\{0,\dots,k-1\}$. Let $m,e,e',L\in\N$ and assume:
(i) for $i\in\mathrm{gen}$, $\mathrm{Fl}_i$ is spanned modulo $\mathrm{Fl}_{i+1}$ by at most $m$ elements of
$\mathrm{Fl}_i$; (ii) $a^{e'}M\subseteq P$; (iii) for $i\notin\mathrm{gen}$,
$\linkNadmlen(\mathrm{Fl}_i/(\mathrm{Fl}_{i+1}+P)\cap\mathrm{Fl}_i)\le L$; (iv) $a^e\mathrm{Fl}_k\subseteq P$. Then
\[
\linkNadmlen(M/P)\le \rank(M)\,e\,k_0+|\mathrm{gen}|\,m\,e'k_0+(k-|\mathrm{gen}|)\,L .
\]
\formallabel{Formalism.MainNotions.RamifiedCoefficientChains.length_le_of_filtration}
\end{lemma}

\begin{proof}
By additivity of length along the chain $M=\mathrm{Fl}_0+P\supseteq\mathrm{Fl}_1+P\supseteq\dots\supseteq
\mathrm{Fl}_k+P\supseteq P$,
\[
\linkNadmlen(M/P)\le\sum_{i<k}\linkNadmlen\bigl(\mathrm{Fl}_i/(\mathrm{Fl}_{i+1}+P)\cap\mathrm{Fl}_i\bigr)
+\linkNadmlen\bigl(\mathrm{Fl}_k/(P\cap\mathrm{Fl}_k)\bigr),
\]
because $(\mathrm{Fl}_i+P)/(\mathrm{Fl}_{i+1}+P)$ is a quotient of $\mathrm{Fl}_i/(\mathrm{Fl}_{i+1}+P)\cap
\mathrm{Fl}_i$ and $(\mathrm{Fl}_k+P)/P\cong\mathrm{Fl}_k/(P\cap\mathrm{Fl}_k)$. The last term is at most
$\rank(M)\,e\,k_0$, since $\mathrm{Fl}_k$ is a submodule of a free module of rank $\rank(M)$ over a principal ideal
domain, hence generated by $\rank(M)$ elements, and is killed by $a^e$ modulo $P$. For $i\in\mathrm{gen}$ the $i$-th
term is at most $m\,e'k_0$, since by (i) the quotient is generated by $m$ elements and by (ii) it is killed by $a^{e'}$;
for $i\notin\mathrm{gen}$ it is at most $L$ by (iii). Summing gives the claim.
\end{proof}

% ---------------------------------------------------------------------------
\subsection{The global setting and the main theorem}\label{pf:linkN-adm:global}

\paragraph{The datum.}
We use the notation of \cref{sec:families:galois} for the Galois side of a datum $D$. Its ranks include
positive integers $f,n,m,k,r,t$ with $m\le k\le r$, $k\ge5$ and $r+5\le t$, and $n_V=r+t$. The host provides a CM field $E$ with
$[E:\Q]=2f$ and class number $h_E$, a finite set $\mathrm{small}$ of small primes containing $2$, primes $q,d$,
two \emph{type primes} $g_0,g_1$ with $\mathcal G=\{g_0,g_1\}$, and a \emph{split bound}; every odd prime $p$ at most
the split bound satisfies $\bigl(\frac{-d}{p}\bigr)=1$. Each prime $\ell\notin\mathrm{small}\cup\{q,d\}$ is either
\emph{split} or \emph{inert} (\lref{Formalism.MainNotions.TwistedPlaces.split_or_inert}); an inert prime
satisfies $\bigl(\frac{-d}{\ell}\bigr)\ne1$. A local index $\iota$ carries a prime $\ell=\ell_\iota$, a discrete
valuation ring $\OO_\iota$ finite over $\Z_\ell$ with fraction field $K_\iota$, finite sets
$Z^{\mathrm{good}}_\iota\subseteq Z_\iota\subseteq X_\iota$ of eigensystems, the congruence lengths
$\linkNadmvcc(\iota,\pi)\in\Q_{\ge0}$, finite sets $S_\iota,Q_\iota$ of places of $E$ with elements
$\zeta_\iota(v)\in\OO_\iota$ ($v\in Q_\iota$), a local level $s_0(\iota)$ and
$\delta_\iota=\sum_{0<k<t}(s_0(\iota)+v_\ell(k))$. It is \emph{twisted} if its place $w_\iota$ is twisted. Lengths of
$\OO_\iota$-modules are normalized: $\linkNadmlen(M)=\linkNadmlen_{\OO_\iota}(M)/\linkNadmlen_{\OO_\iota}(\OO_\iota/\ell)$.
Galois families are maps $r_\iota\colon X_\iota\to\Hom(G_E,\GL_{n_V}(\OO_\iota))$ and $\rho_\iota\colon
X_\iota\to\Hom(G_E,\GL_r(\OO_\iota))$; $\linkNadmsP(\iota,\rho_\iota,\pi)$ is the normalized length of the Selmer
group of $\rho_{\iota,\pi}$ of \cref{sec:families:galois}. For a prime $\ell$, $\mathrm{sh}(\ell,e)$ is the
$\Z_\ell$-length of $\Sha^1(G_{E,S_\ell},\Q_\ell/\Z_\ell(-e))$ ($0$ if it is infinite) and $\linkNadmVV(\ell)=\sum_{0<e<t}
\mathrm{sh}(\ell,e)$; condition (Sha) means that all these groups with $0<e<t$ are finite, for every prime $\ell$.

For a place $\lambda\mid\ell$ of $E$ let $n_\lambda=e_\lambda f_\lambda$ be its local degree, so
$\sum_{\lambda\mid\ell}n_\lambda=2f$ (\lref{Formalism.MainNotions.TwistedPlaces.sum_nLam_le}), let $K_\lambda$
be the completion and $G_\lambda$ the decomposition group. A \emph{coefficient extension} at $\lambda$ is a finite
extension $L_C$ of $\Q_\ell$ containing the local coefficient field of the datum at $\lambda$ (which contains all
conjugates of $K_\lambda$); it is \emph{good} if it contains the roots of unity of order dividing $\lcm(\ell,N)$, $N$ the
conductor of the host. If $\ell\ne2$ and $K_\lambda/\Q_\ell$ is unramified, a good $L_C$ contains all conjugates of
$K_\lambda(\mu_\ell)$; we call such $L_C$ \emph{admissible}. The
localizations at $\lambda$, with coefficients in $L_C$, of the twisted duals
$\rho^\vee_{\iota,\pi}(\chi_{\iota,1}\varepsilon_\ell)$ and $r^\vee_{\iota,x}(\chi_{\iota,1}\varepsilon_\ell)$
($\varepsilon_\ell$ the cyclotomic character, $\chi_{\iota,1}$ a character of the datum) are denoted $T^\vee_\pi$ and
$r^\vee_x$; the localization of the twist $r_{\iota,x}\chi_{\iota,1}^{-1}$ is denoted $r^{\mathrm{tw}}_x$. The datum
provides period data $\mathrm{PD}_\lambda$ for $G_\lambda$ and, for every admissible $L_C$, a \emph{Wach structure}
$W_{\lambda,C}$: the $(\varphi,\Gamma)$-module theory with coefficients in $\OO_C$ in the sense of
\cite{Wach1996,Berger2004}, axiomatized (\cref{pf:linkN-adm:wach}). Its \emph{letters} are the
$n_\lambda=[K_\lambda:\Q_\ell]$ embeddings of $K_\lambda$ into $L_C$, indexed by $\Z/n_\lambda$ compatibly with
Frobenius; an embedding $\tau$ \emph{lies over} the letter $a$ if it is the $a$-th one.

\paragraph{Size data.}
Size data $L$ of $D$ provide a scale $x>1$, $b_L=e^\Lambda$ with $\Lambda\ge1$
(\lref{Formalism.MainNotions.TwistedPlaces.L_ge_one}), and positive reals $S_c,\eta_{\mathrm{sz}}$. Put
\[
h_B=b_L+S_c+\eta_{\mathrm{sz}}+t+2,\qquad \eta_3=\eta_{\mathrm{sz}}+t+2,\qquad G_P=w_P+S_c+\eta_3,
\]
where $w_P=w_P(L)$ is an explicit width, of the form $w_B+(t+1)/2$ with $w_B\ge0$ a maximum of non-negative
quantities, so $w_P\ge0$; let $\ell_A(L)$, $\ell^*(L)$ be the thresholds of $L$ and $j_A(\ell)=\max(1,\lceil
\log_{\ell-1}\ell_A\rceil)$. The \emph{parameter bound} (PB) is the inequality $S_c+\eta_{\mathrm{sz}}+t+2\le b_L$; it
gives $h_B\le2b_L=2e^\Lambda$.

\paragraph{The middle range.}
For a finite set $\mathrm{exc}$ of primes, $\ell$ is in the \emph{middle range} if
$\ell\notin\mathrm{small}\cup\{3,q,d\}$ and ($\ell\le\ell^*(L)$ or $\ell\in\mathrm{exc}$). Then $\ell\ge5$ is odd and
$s_0(\iota)=1$ (\lref{UpperTheory.Small.localLevel_eq_one_of_not_mem_smallPrimes}), so
$\delta_\iota\le t+v_\ell(t!)\le2t$.

\begin{definition}[Hub data and rows]\label{pf:linkN-adm:def-hub}
\emph{Hub data} $X$ for $D$ consist of families $r=(r_\iota)$, $\rho=(\rho_\iota)$, local lengths
$\linkNadmlen_0(\iota,\pi,\lambda)\in\R$, error terms $\linkNadmerr(\iota,\pi)\in\R$, for each $\iota$ a
\emph{declaration} (a partition of the places above $\ell_\iota$ into \emph{rows} and \emph{columns}), row lengths
$\linkNadmlen_{\mathrm{row}}(\iota,\lambda)\in\Q$, exponents $e'(\iota,\lambda),c_B(\iota),e_S(\iota),c_0(\iota,\lambda)\in\N$,
a predicate $\mathrm{band}(\iota,\lambda)$ and a finite set $\mathrm{exc}$ of primes. For a local index $\iota$ with
declaration $(\mathrm{row},\mathrm{col})$ put
\begin{align*}
\linkNadmerr_{\mathrm{dec}}(\iota)&=t\sum_{w\in\mathrm{row}}(n_w+2)\,c_B(\iota)+4n_V(|\mathrm{col}|+1)\,\delta_\iota
+t\,|\mathrm{col}|\,h_0(\iota,\mathrm{col}),\\
\linkNadmerr_{\mathrm{pl}}(\iota,\pi)&=\sum_{u\in\mathrm{col}}\bigl(n_u m\,\delta_\iota+2\delta_\iota+h^W_0(\iota,u,\pi)
+r\,c_0(\iota,u)\bigr)+\linkNadmerr_{\mathrm{dec}}(\iota),
\end{align*}
where $h_0(\iota,\mathrm{col})$ is the maximum over $u\in\mathrm{col}$ and $0<e<t$ of the normalized length of the
$G_u$-invariants of $K_\iota/\OO_\iota$ twisted by $\varepsilon_\ell^{-e}$, and $h^W_0(\iota,u,\pi)$ is the normalized
length of the $G_u$-invariants of $(K_\iota/\OO_\iota)^r$ with the action of $\rho_{\iota,\pi}\chi_{\iota,1}^{-1}$. A prime
$\ell$ is of \emph{row~B} if $\ell<\ell_A(L)$ and the order of the image of $\ell$ in $(\Z/2^{e+2}\Z)^\times/\{\pm1\}$
exceeds $4L_w$ ($e$, $L_w$ data of $D$); otherwise of \emph{rows~C/D}. With
\[
\mathrm{core}_B(\ell,\lambda)=(n_\lambda+1)\,r\,\Bigl(\frac{h_B}{(\ell-1)^{j_A(\ell)+1}}+(G_P+1)(j_A(\ell)+1)+\Lambda\Bigr)
\]
let $\linkNadmlen_B(\ell,\lambda)$ be a rational number strictly between $3C_2\,\mathrm{core}_B$ and
$(3C_2+1)\mathrm{core}_B$ when $3C_2\,\mathrm{core}_B<(3C_2+1)\mathrm{core}_B$, and $0$ otherwise. For a proposition
$\beta$ put
$[\beta]\in\{0,1\}$, $\mathrm{core}_E(\ell)=1+h_B/\ell+\log_\ell h_B$ and
\begin{multline}\label{pf:linkN-adm:eq-bdCD}
\mathrm{bd}_{CD}(\beta,\ell,\lambda;K_1,K_2,K_3)=K_1\frac{(n_\lambda+1)rh_B}{\ell^2}\\
+K_2(n_\lambda+1)r\Bigl(\frac{w_P+\eta_3}{\ell}+S_c+\Lambda+\eta_3[\beta]\Bigr)+K_3\,n_\lambda m\,\mathrm{core}_E(\ell);
\end{multline}
$\linkNadmlen_{CD}(\ell,\lambda)$ is a rational number strictly between $\mathrm{bd}_{CD}(\beta;K_1,K_2,K_3)$ and
$\mathrm{bd}_{CD}(\beta;K_1,K_2+1,K_3)$ with $\beta=\mathrm{band}(\iota,\lambda)$ (chosen once and for all; $0$ if the
interval is empty).
The hub data are \emph{adapted} to $L$ and constants $(C_2,C_B,K_1,K_2,K_3,C_A)$ if for every $\iota$:
\begin{itemize}
\item if $\ell_\iota$ is inert: there are no columns; $\linkNadmlen_{\mathrm{row}}(\iota,\cdot)$ is $\linkNadmlen_B$ in
  row~B and $\linkNadmlen_{CD}$ otherwise; $\linkNadmlen_0(\iota,\pi,\lambda)=\linkNadmlen_{\mathrm{row}}(\iota,\lambda)$
  and $\linkNadmerr(\iota,\pi)=\linkNadmerr_{\mathrm{dec}}(\iota)$;
\item if $\ell_\iota$ is split: $e_S(\iota)=e_A(\ell)=\lceil C_A(1+2\lfloor h_B\rfloor/\ell+1+\log_\ell(2\lfloor
  h_B\rfloor+1))\rceil$, $\linkNadmlen_0(\iota,\pi,\lambda)=(2n_\lambda m+1)e_S(\iota)$ on columns and $0$ on rows,
  and $\linkNadmerr(\iota,\pi)=\linkNadmerr_{\mathrm{pl}}(\iota,\pi)$;
\item $\mathrm{band}(\iota,\lambda)$ is the band predicate of \cref{pf:linkN-adm:def-words}.
\end{itemize}
The constants used are the \emph{row constants} $\mathcal K(C_2)=(C_2,C_B,C_C,C_A)=(C_2,320,25,23)$, for which
$K_1=C_B+48=368$, $K_2=6C_B+12+6C_C=2082$, $K_3=3C_AC_B+12C_A+9=22365$.
\formal{Formalism.MainNotions.N.KillingConditions.HubData}
\formal{Formalism.MainNotions.N.KillingConditions.DataAD}
\formal{GlobalTheory.NLink.rowC}
\end{definition}

\begin{definition}[The two inequalities]\label{pf:linkN-adm:def-star}
Let $X$ be hub data. $X$ satisfies the \emph{star inequality} if for every twisted local index $\iota$ with
$\ell=\ell_\iota\notin\mathcal G$ in the middle range for $\mathrm{exc}=X.\mathrm{exc}$, and every $\pi\in
Z^{\mathrm{good}}_\iota$, inequality \eqref{pf:linkN-adm:eq-mid} holds, with $\rho_\iota=X.\rho_\iota$,
$\linkNadmlen_0=X.\linkNadmlen_0$ and $\linkNadmerr=X.\linkNadmerr$. $X$ satisfies the \emph{error bound with constant
$C_5$} if for the same $\iota,\pi$, $\linkNadmerr(\iota,\pi)\le C_5fnt^2n_V\Lambda^2$.
\formal{Formalism.MainNotions.N.KillingConditions.StarAD}
\formal{Formalism.MainNotions.N.KillingConditions.BBLe}
\end{definition}

\begin{definition}[Words]\label{pf:linkN-adm:def-words}
Let $\lambda\mid\ell$ be a place with $K_\lambda/\Q_\ell$ unramified of degree $n_\lambda$, and $L_C$ an admissible
coefficient extension. A \emph{word} on the letters $\Z/n_\lambda$ consists of integers $b,w,S,\eta,m\ge0$, a \emph{type}
$\mathrm{up}\colon\Z/n_\lambda\to\{\mathrm{true},\mathrm{false}\}$ and functions
$\mathrm{sz},M^-,M^+\colon\Z/n_\lambda\to\N$. The \emph{window} of the letter $a$ is
\[
\mathrm{win}(a)=\begin{cases}[\mathrm{sz}_a,\mathrm{sz}_a+S]\cup[-\eta,\eta]&\text{if }\mathrm{up}(a),\\
[-\mathrm{sz}_a-S,-\mathrm{sz}_a]\cup[-\eta,\eta]&\text{otherwise;}\end{cases}
\]
$a$ is \emph{thick} if $b\le\mathrm{sz}_a\le b+\eta$. A lattice $T$ has its jumps \emph{in the windows} if for every
letter $a$ and every $\tau$ over $a$ all jumps of $T$ at $\tau$ lie in $\mathrm{win}(a)$, and \emph{within $(M^-,M^+)$} if
they lie in $[-M^-_a,M^+_a]$. The low count of a lattice $T$ at $a$ is $\sum_{\tau\text{ over }a}\sum_{|x|\le\eta}
\mathrm{mult}_\tau(T,x)$. The word is \emph{based} at $(\iota,\pi,\lambda,C)$ if
$K_\lambda$ is unramified of degree $n_\lambda$, $L_C$ contains all conjugates of $K_\lambda$, $\OO_C$ is a discrete
valuation ring with fraction field $L_C$, and:
(a) every lattice $T_1$ of a subrepresentation of a finite direct sum of copies of the $r^\vee_x$ ($x\in X_\iota$) is
crystalline with jumps in the windows and within $(M^-,M^+)$; (b) every lattice $T_2$ of $T^\vee_\pi$ is crystalline with
low count at most $m$ at every letter; (c) every letter is thick or has $\mathrm{sz}_a\le w$; (d) $b+\eta+S\le h_B$,
$w\le w_P$, $S\le S_c$, $\eta\le\eta_3$ and $m$ equals the rank $m$ of $D$. Further conditions on a word:
\begin{itemize}
\item[(B$_{C_2}$)] $j_A(\ell)+1\le n_\lambda$, $C_2(j_A(\ell)+1)(w+S+\eta+1)\le b$, and the thick letters
  form \emph{clusters} of span $j=j_A(\ell)$: for each thick $t_0$, either none of $t_0-1,\dots,t_0-j$ is thick, or there
  is $d\in[1,j]$ with $t_0-d$ thick and no other thick letter among $t_0-1,\dots,t_0-d-j$;
\item[(Box)] $M^-_a+M^+_a\le2h_B$; for a non-thick letter, $|\mathrm{sz}_a-s|\le(t+1)/2$ for some $s$ in
  an explicit interval $[w''(L),w_B(L)]$; and $b_L-(t+1)/2\le b\le b_L-(t+1)/2+1$;
\item[(NT)] (a word with a non-thick letter) $\eta<\mathrm{sz}_a$ for all $a$, $20w\le b$, $20\eta\le b$, $20S\le b$,
  some letter is not thick, and the type alternates ($\mathrm{up}(a+1)\ne\mathrm{up}(a)$) with no three consecutive
  thick letters;
\item[(AT$_d$)] (an all-thick word of parity $d$) $n_\lambda$ is even, $\mathrm{up}(a)$ holds iff $a+d$ is odd, all
  letters are thick, $2a_R(\ell,H)+\eta+1\le\mathrm{sz}_a$ and $H_0/2\le\mathrm{sz}_a$ for all $a$, where
  $H_0=\max_a\mathrm{sz}_a+S$, $H=H_0+\eta$ and $a_R$ is defined in \eqref{pf:linkN-adm:eq-aR};
\item[(Mult)] $L_C$ contains a primitive $\ell$-th root of unity, and every lattice $T_2$ of $T^\vee_\pi$
  satisfies $\sum_{|x|\le\eta}\mathrm{mult}_\tau(T_2,x)\le m$ for every embedding $\tau$.
\end{itemize}
The \emph{band} condition of a word with constant $C_B$ is: some non-thick letter $a$ has a thick neighbour and
$|\ell\,\mathrm{sz}_a-b|\le C_B(\ell\eta+\ell S+b/\ell+\mathrm{sz}_a+\ell\log_\ell(2h))$, where $h=b+\eta+S$ if some
letter is thick and $h=w+S$ otherwise. The \emph{band predicate} $\mathrm{band}(\iota,\lambda)$ holds if $\ell\ge11$ and
some $\pi\in Z^{\mathrm{good}}_\iota$, admissible $L_C$ and word based at $(\iota,\pi,\lambda,C)$ satisfying (Box) and
(NT) satisfy the band condition with $C_B=320$.
\formal{Formalism.MainNotions.N.WordBase}
\formal{Formalism.MainNotions.N.WordsB}
\formal{Formalism.MainNotions.N.WordBox}
\formal{Formalism.MainNotions.N.WordsNT}
\formal{Formalism.MainNotions.N.WordsAT}
\formal{Formalism.MainNotions.N.WordBaseC}
\formal{Formalism.MainNotions.N.bandOf}
\end{definition}

\begin{definition}[Joint hub conditions; Corollary A]\label{pf:linkN-adm:def-joint}
The \emph{joint hub conditions} for the row constants $\mathcal K(C_2)$ state: for every
twisted $\iota$ with $\ell=\ell_\iota$ inert, every $\pi\in Z^{\mathrm{good}}_\iota$ and every $\lambda\mid\ell$, $K_\lambda$
is unramified over $\Q_\ell$ and, for every admissible $L_C$,
\begin{itemize}
\item in row~B: there is a word based at $(\iota,\pi,\lambda,C)$ satisfying (B$_{C_2}$);
\item in rows~C/D: there is a word based at $(\iota,\pi,\lambda,C)$ satisfying (Box) and either (NT), or (Mult) and
  (AT$_d$) for some $d$.
\end{itemize}
\emph{Corollary A with constant $C$} states: for every prime $\ell\ge11$, unramified $K$ with $n$ letters, coefficients
$L_C$ satisfying the standard conditions (they split $K$, $\OO_C$ is a discrete valuation ring with fraction field $L_C$),
crystalline lattices $T_1,T_2$ and an onto lattice map $F\colon T_1\to T_2$, and every word $(b,w,S,\eta,\mathrm{up},
\mathrm{sz})$ with $1\le j$, $j+1\le n$, jumps of $T_1$ in the windows, every letter thick or of size at most $w$,
$C(j+1)(w+S+\eta+1)\le b$, and the thick letters in clusters of span $j$, one has
\begin{equation}\label{pf:linkN-adm:eq-corA}
\linkNadmcol(F)\preceq C(n+1)\,\rank T_2\,\Bigl(\frac{h_{\mathrm{cor}}}{(\ell-1)^{j+1}}+(w+S+\eta+1)(j+1)
+\log_\ell h_{\mathrm{cor}}\Bigr),
\end{equation}
with $h_{\mathrm{cor}}=b+\eta+S$ if some letter is thick and $h_{\mathrm{cor}}=w+S+\eta$ otherwise.
\formal{Formalism.MainNotions.N.KillingConditions.JointParametersHubConditions}
\formal{Formalism.MainNotions.N.LocalRepresentationBoundCondition}
\formal{Formalism.MainNotions.N.LocalRepresentationCongruenceCondition}
\formal{Formalism.MainNotions.WachCoordinates.Conclusion2}
\end{definition}

\begin{definition}[The setup conditions]\label{pf:linkN-adm:def-setup}
Let $\mathrm{nch}$ be a norm choice for $D$ and $C_2\in\R$. The hub data $X$ satisfy the \emph{setup conditions} if: the
norm-choice facts hold for $\mathrm{nch}$; $X$ is the \emph{canonical} hub data for the row constants $\mathcal K(C_2)$ (its
families are good and continuous, it is adapted in the sense of \cref{pf:linkN-adm:def-hub}, and its declarations at
split primes and its exponents $c_B,c_0$ are the canonical ones); the parameters are admissible; (PB) holds;
$\Lambda\ge8$; the split bound is at least $60$; the \emph{killing condition} holds; and (Sha) holds. The killing
condition states: for twisted $\iota$ and $\pi\in Z^{\mathrm{good}}_\iota$ the Selmer group of $X.\rho_{\iota,\pi}$ is
finite with $\linkNadmsP(\iota,X.\rho_\iota,\pi)\le\nu_\ell+\ell^{-2}$, and
\begin{equation}\label{pf:linkN-adm:eq-kill}
t\sum_{v\in Q_\iota}\linkNadmlen\bigl(\OO_\iota/(\zeta_\iota(v)-1)\bigr)\le\ell^{-2}.
\end{equation}
\formal{GlobalTheory.NLink.SetupAlt}
\formal{Formalism.MainNotions.N.KillingConditions.KillAt}
\end{definition}

\begin{theorem}[The middle-range bounds]\label{pf:linkN-adm:main}
Assume $\sfname{Hypotheses}_2$. Let $D$ be a datum with size data $L$, parameter choice, norm choice
$\mathrm{nch}$, local Hodge--Tate, Fontaine and comparison data, hub data $X$ and $C_2\in\R$ such that the setup
conditions hold, $C_2\ge1$, Corollary A holds with constant $C_2$, and the joint hub conditions hold for the row constants
$\mathcal K(C_2)$. Then $X$ satisfies the star inequality and the error bound with constant $7200$: for every twisted local
index $\iota$ with $\ell_\iota\notin\mathcal G$ in the middle range and every $\pi\in Z^{\mathrm{good}}_\iota$,
\eqref{pf:linkN-adm:eq-mid} and \eqref{pf:linkN-adm:eq-bb} hold.
\formallabel{GlobalTheory.T.star_ad_and_bb_le_of_joint_cond}
\formallabel{GlobalTheory.NLink.rowsAD_in_setting_alt}
\end{theorem}

The proof is given in \cref{pf:linkN-adm:proofmain}. The theorem supplies the fields \emph{mid} and \emph{bb} of the
conclusion of the table of local bounds in the proof of \cref{thm:linkN}; there $C_2$ is the constant of
``Theorem~E and Corollary~A'' and the joint hub conditions come from Fontaine's theory.

% ---------------------------------------------------------------------------
\subsection{Admissibility and the star inequality}\label{pf:linkN-adm:starsec}

\paragraph{Row and lift bounds.}
Let $\iota$ be a local index, $\lambda\mid\ell_\iota$ and $L_C$ a coefficient extension at $\lambda$. For $x\in\R$, the
\emph{row bound} $\mathrm{Row}_{\lambda,C}(x)$ for $(r_\iota,\rho_\iota,\pi)$ is the statement: for every lattice $T_2$ of
$T^\vee_\pi$, every lattice $T_1$ of a subrepresentation of a finite direct sum of copies of the $r^\vee_x$ ($x\in
X_\iota$), and every onto lattice map $F\colon T_1\to T_2$ of $G_\lambda$-representations, $\linkNadmcol(F)\preceq x$
(\cref{pf:linkN-adm:def-col}, with $\mathrm{PD}=\mathrm{PD}_\lambda$). Note $\rank T_2=r$. The \emph{lift bound}
$\mathrm{Lift}_{\lambda,C}(e')$ is the statement: every onto lattice map between lattices of subquotients of finite direct
sums of copies of the $r^{\mathrm{tw}}_x$ lifts with exponent $e'$.

\begin{definition}[Admissibility at an inert point]\label{pf:linkN-adm:def-adm}
Let $\iota$ be a local index, $r_\iota,\rho_\iota$ families, $\pi\in X_\iota$, $(\mathrm{row},\mathrm{col})$ a declaration
at $\iota$, $\linkNadmlen_{\mathrm{row}}\colon\{\lambda\mid\ell\}\to\Q$, $e'\colon\{\lambda\mid\ell\}\to\N$ and
$c_B\in\N$. The point is \emph{admissible} if twenty conditions hold. Seventeen of
them concern the global and local Galois representations only: (G1) the coefficients of the characteristic
polynomials of the $r_{\iota,x}$ lie in the Hecke algebra; (G2) the \emph{gain condition}: a gain bound for the twisted
localizations at the places above $\ell$, and their crystallinity; (G3) absolute irreducibility and the residual shape
of $r_{\iota,x}$ relative to $\rho_{\iota,\pi}$; unramifiedness outside $\ell$, $S_\iota$ and $Q_\iota$; the tame shape at
$S_\iota\cup Q_\iota$; the tame generator conditions (in particular $\zeta_\iota(v)\ne1$ for $v\in Q_\iota$); an
eigenvector criterion for membership in $Z_\iota$; local Tate duality for $H^2$, the Bloch--Kato dimension formula and
Bloch--Kato duality at good coefficients; continuity; semisimplicity; the decomposition $r_{\iota,\pi}\otimes K_\iota\cong
(\rho_{\iota,\pi}\oplus\bigoplus_{c=1}^t\chi_{\iota,c})\otimes K_\iota$; a cyclotomic inertia image; compatibility of local
coefficients; that each $K_\lambda$ lies in $\Q_\ell(\mu_{\lcm(\ell,N)})$; Fontaine's properties of crystalline
representations at the places above $\ell$; an integral $H^2$-duality; and the bound $\linkNadmlen\le v_\ell(h_E)$ for
the module of everywhere unramified homomorphisms $G_E\to K_\iota/\OO_\iota$. The other three are:
\begin{itemize}
\item[(TB)] the \emph{Bloch--Kato condition with exponent $c_B$}: at every $\lambda\mid\ell$ and good
  $L_C$, for the crystalline lattices with jumps among those of $r^{\mathrm{tw}}_\pi$, multiplication by $p^{c_B}$ makes the
  relevant extension cocycles with values in $L_C(\varepsilon^{-e})/\OO_C(\varepsilon^{-e})$, $0\le e<t$, finite, and
  the finite divisible cocycles of these twists are coboundaries ($e>0$), resp.\ unramified up to coboundaries ($e=0$),
  with colength at most $([K_\lambda:\Q_\ell]+1)c_Be_C$ for those that become finite after multiplication by
  $p^{c_B}$;
\item[(Row)] for every row place $\lambda$ and good $L_C$: $\mathrm{Row}_{\lambda,C}(\linkNadmlen_{\mathrm{row}}(\lambda))$;
\item[(Lift)] for every column place $\lambda$ and good $L_C$: $\mathrm{Lift}_{\lambda,C}(e'(\lambda))$.
\end{itemize}
\formal{Formalism.MainNotions.TruncatedSeriesCountingIdentities.admitted}
\formal{Formalism.MainNotions.TruncatedSeriesCountingIdentities.admTB}
\formal{Formalism.MainNotions.TruncatedSeriesCountingIdentities.LengthBoundAtPlaces}
\formal{Formalism.MainNotions.N.RowAt}
\end{definition}

We use two results proved elsewhere. The first is the \emph{integrality-from-metric theorem}: under
$\sfname{Hypotheses}_2$ and a design condition contained in the norm-choice facts,
\lref{GlobalTheory.T.integralityFromMetric_of_designConditionAll} proves that for every local index $\iota$, good
families $r_\iota$, $\rho_\iota$, $\pi$, declaration and admissible data as in \cref{pf:linkN-adm:def-adm}: if
$\ell\notin\mathcal G$, $\pi\in Z_\iota$, the Selmer group of $\rho_{\iota,\pi}$ is finite and the groups
$\Sha^1(G_{E,S_\ell},\Q_\ell/\Z_\ell(-e))$, $0<e<t$, are finite, then
\begin{multline}\label{pf:linkN-adm:eq-ifm}
\linkNadmvcc(\iota,\pi)\le\sum_{\lambda\mid\ell}\linkNadmlen_W(\lambda)+\linkNadmsP(\iota,\rho_\iota,\pi)
+\sum_{0<e<t}\mathrm{sh}(\ell,e)+v_\ell(h_E)\\
+t\sum_{w\in\mathrm{row}}(n_w+2)c_B+4n_V(|\mathrm{col}|+1)\delta_\iota+t|\mathrm{col}|\,h_0(\iota,\mathrm{col})
+t\sum_{v\in Q_\iota}\linkNadmlen\bigl(\OO_\iota/(\zeta_\iota(v)-1)\bigr),
\end{multline}
where $\linkNadmlen_W(\lambda)=\linkNadmlen_{\mathrm{row}}(\lambda)$ on rows and is an explicit column length on columns.
The second is the \emph{split-case bound} \lref{GlobalTheory.T.splitCase_bound_at_point_of_admitted_items_alt}:
under global Artin reciprocity, local Tate duality and the Euler characteristic formula, let a point satisfy the
thirty-five \emph{split admissibility conditions} (\lref{Formalism.MainNotions.EmbeddingTypes.admitted}: they contain
the conditions above except (Row), in a split form, and add column conditions with an exponent $e'$ and constants
$c_0(\lambda)$) and, if $\ell\notin\mathcal G$ and $\pi\in Z_\iota$, the gain condition (G2). If $\ell$ is split,
all places of $E$ above $\ell$ are unramified over $\Q$, $\ell\notin\mathcal G$, the declaration is valid for the chosen
embedding types, $\pi\in Z_\iota$, the Selmer group of $\rho_{\iota,\pi}$ and the groups $\Sha^1(G_{E,S_\ell},\Q_\ell/\Z_\ell(-e))$
($0<e<t$) are finite and $\zeta_\iota(v)\ne1$ for $v\in Q_\iota$, then
\begin{multline}\label{pf:linkN-adm:eq-split}
\linkNadmvcc(\iota,\pi)\le\sum_{u\in\mathrm{col}}\bigl(e'+n_um(e'+\delta_\iota)+2\delta_\iota+h^W_0(\iota,u,\pi)
+n_ume'+r\,c_0(u)\bigr)\\
+\linkNadmsP(\iota,\rho_\iota,\pi)+\sum_{0<e<t}\mathrm{sh}(\ell,e)+v_\ell(h_E)
+t\sum_{w\in\mathrm{row}}(n_w+2)c_B+4n_V(|\mathrm{col}|+1)\delta_\iota\\
+t|\mathrm{col}|\,h_0(\iota,\mathrm{col})+t\sum_{v\in Q_\iota}\linkNadmlen\bigl(\OO_\iota/(\zeta_\iota(v)-1)\bigr).
\end{multline}

For hub data $X$ and a local index $\iota$, the \emph{point of $X$ at $(\iota,\pi)$} consists of $\iota$, $\pi$ and the
data of $X$ at $\iota$: the families $X.r_\iota,X.\rho_\iota$, the declaration, the row lengths
$X.\linkNadmlen_{\mathrm{row}}(\iota,\cdot)$ and the exponents $X.e'(\iota,\cdot)$, $X.c_B(\iota)$. The \emph{step
conditions} are: (a) for every twisted $\iota$ with $\ell_\iota\notin\mathcal G$ in the middle range and inert, and every
$\pi\in Z^{\mathrm{good}}_\iota$, the point of $X$ at $(\iota,\pi)$ is admissible; (b) for every such $\iota$ with
$\ell_\iota$ split there is a choice of embedding types for which the declaration is valid, all places above $\ell_\iota$
are unramified in $E/\Q$, and every $\pi\in Z^{\mathrm{good}}_\iota$ satisfies the split admissibility conditions with
$e'=X.e_S(\iota)$ and $c_0=X.c_0(\iota,\cdot)$; (c) the families $X.r_\iota$ are good.

\begin{proposition}[The star inequality]\label{pf:linkN-adm:star}
Let $X$ be hub data adapted to $L$ and some constants, satisfying the step conditions. Assume the conclusion of the
integrality-from-metric theorem \eqref{pf:linkN-adm:eq-ifm} for $D$, the split-case bound \eqref{pf:linkN-adm:eq-split}
for $D$, the gain condition (G2) at every $\pi\in Z^{\mathrm{good}}_\iota$ with $\ell_\iota\notin\mathcal G$,
the killing condition for $X.\rho$ and (Sha). Then $X$ satisfies the star inequality.
\formallabel{GlobalTheory.T.starInequality_rowsAD_of_global_theorems}
\formallabel{GlobalTheory.NLink.starAD_of_global}
\end{proposition}

\begin{proof}
Fix a twisted $\iota$ with $\ell=\ell_\iota\notin\mathcal G$ in the middle range and $\pi\in Z^{\mathrm{good}}_\iota$. Then
$\ell\notin\mathrm{small}\cup\{3,q,d\}$, so $\ell$ is odd, and $\ell\ne g_0,g_1$; $\pi\in Z_\iota$ since
$Z^{\mathrm{good}}_\iota\subseteq Z_\iota$. By the killing condition the Selmer group of $X.\rho_{\iota,\pi}$ is finite and
\eqref{pf:linkN-adm:eq-kill} holds; by (Sha) the groups $\Sha^1(G_{E,S_\ell},\Q_\ell/\Z_\ell(-e))$,
$0<e<t$, are finite, and $\linkNadmVV(\ell)=\sum_{0<e<t}\mathrm{sh}(\ell,e)$. The prime $\ell$ is split or inert.

\emph{Split case.} By (b) of the step conditions choose embedding types with a valid declaration and unramified places;
the point satisfies the split admissibility conditions, and the tame generator conditions among them give
$\zeta_\iota(v)\ne1$ for $v\in Q_\iota$. The split-case bound \eqref{pf:linkN-adm:eq-split} applies with $e'=e_S(\iota)$.
Since the rows and columns partition the places above $\ell$ and $X$ is adapted,
$\sum_{\lambda\mid\ell}\linkNadmlen_0(\iota,\pi,\lambda)=\sum_{u\in\mathrm{col}}(2n_um+1)e_S(\iota)$. Splitting the column sum
in \eqref{pf:linkN-adm:eq-split} as $\sum_{u}(2n_um+1)e_S+\sum_u(n_um\delta_\iota+2\delta_\iota+h^W_0+rc_0(u))$ and using
$\linkNadmerr(\iota,\pi)=\linkNadmerr_{\mathrm{pl}}(\iota,\pi)$, the right-hand side of \eqref{pf:linkN-adm:eq-split} equals
\[
\sum_{\lambda\mid\ell}\linkNadmlen_0(\iota,\pi,\lambda)+\linkNadmsP+\linkNadmVV(\ell)+v_\ell(h_E)+\linkNadmerr(\iota,\pi)
+t\sum_{v\in Q_\iota}\linkNadmlen(\OO_\iota/(\zeta_\iota(v)-1)),
\]
and \eqref{pf:linkN-adm:eq-kill} bounds the last term by $\ell^{-2}$.

\emph{Inert case.} By (a) of the step conditions the point is admissible, and the families are good by (c); so
\eqref{pf:linkN-adm:eq-ifm} applies. As $X$ is adapted, there are no columns, so every place above $\ell$ is a row,
$\linkNadmlen_W(\lambda)=\linkNadmlen_{\mathrm{row}}(\iota,\lambda)=\linkNadmlen_0(\iota,\pi,\lambda)$, and the terms with
$|\mathrm{col}|$ in \eqref{pf:linkN-adm:eq-ifm} reduce to $4n_V\delta_\iota$; hence the second line of
\eqref{pf:linkN-adm:eq-ifm} is $\linkNadmerr_{\mathrm{dec}}(\iota)=\linkNadmerr(\iota,\pi)$ plus the $Q_\iota$-term, which is at
most $\ell^{-2}$ by \eqref{pf:linkN-adm:eq-kill}. This is \eqref{pf:linkN-adm:eq-mid}.
\end{proof}

% ---------------------------------------------------------------------------
\subsection{The size of the error term}\label{pf:linkN-adm:bbsec}

\begin{definition}[Invariant bounds]\label{pf:linkN-adm:def-inv}
For $c\in\R$, hub data $X$ satisfy the \emph{invariant bounds with constant $c$} if for every twisted $\iota$ with
$\ell_\iota\notin\mathcal G$ in the middle range: $c_B(\iota)\le c\,tnn_V\Lambda$, $h_0(\iota,\mathrm{col})\le ct$, and, if
$\ell_\iota$ is split, for every column $u$: $c_0(\iota,u)\le c\,nn_V\Lambda$ and $h^W_0(\iota,u,\pi)\le c\,rnn_V\Lambda$ for
$\pi\in Z^{\mathrm{good}}_\iota$.
\formal{Formalism.MainNotions.N.KillingConditions.InvariantsBoundsWithCoarseColumnConstant}
\end{definition}

\begin{lemma}\label{pf:linkN-adm:blk}
Let $\iota$ be a local index with $\ell_\iota\notin\mathrm{small}$, $(\mathrm{row},\mathrm{col})$ a declaration at $\iota$,
$c\ge1$, $c_B\le c\,tnn_V\Lambda$ and $h_0(\iota,\mathrm{col})\le ct$. Then $\linkNadmerr_{\mathrm{dec}}(\iota)\le
32\,c\,fnt^2n_V\Lambda$. If moreover $c_0(u)\le c\,nn_V\Lambda$ and $h^W_0(\iota,u,\pi)\le c\,rnn_V\Lambda$ for every column
$u$, then $\linkNadmerr_{\mathrm{pl}}(\iota,\pi)\le48\,c\,fnt^2n_V\Lambda$.
\formallabel{GlobalTheory.NLink.BBSize.blockBoundDeclaration_le_middleRange}
\formallabel{GlobalTheory.NLink.BBSize.blockBoundAtZero_le_middleRange}
\end{lemma}

\begin{proof}
We use $f,t,n,n_V,\Lambda\ge1$, $\delta_\iota\le2t$, $m\le r\le t$, $\sum_{w\mid\ell}n_w\le2f$, and $|S|\le\sum_{w\in S}n_w$
for every set $S$ of places above $\ell$ (as $n_w\ge1$). Put $A=\sum_{w\in\mathrm{row}}(n_w+2)\le2f+2\cdot2f=6f$ and
$k=|\mathrm{col}|\le2f$. Then
\begin{align*}
tAc_B&\le t\cdot6f\cdot ctnn_V\Lambda=6\,cfnt^2n_V\Lambda,\\
4n_V(k+1)\delta_\iota&\le4n_V\cdot3f\cdot2t\le24\,cfnt^2n_V\Lambda,\\
tk\,h_0(\iota,\mathrm{col})&\le t\cdot2f\cdot ct\le2\,cfnt^2n_V\Lambda,
\end{align*}
and the sum is $32\,cfnt^2n_V\Lambda$. For the second bound, $\sum_{u\in\mathrm{col}}(n_um\delta_\iota+2\delta_\iota+h^W_0+rc_0(u))
\le A'm\delta_\iota+k(2\delta_\iota+2crnn_V\Lambda)$ with $A'=\sum_{u\in\mathrm{col}}n_u\le2f$; and $A'm\delta_\iota\le
2f\cdot t\cdot2t$, $2k\delta_\iota\le2f\cdot4t$, $2kcrnn_V\Lambda\le4cfnt\,n_V\Lambda$, which add up to at most
$16\,cfnt^2n_V\Lambda$. Adding the first bound gives $48$.
\end{proof}

\begin{proposition}\label{pf:linkN-adm:bbmain}
Let $c\ge1$ and let $X$ be adapted hub data satisfying the invariant bounds with constant $c$. Then $X$ satisfies the error
bound with constant $60c$.
\formallabel{GlobalTheory.T.nonPlacePart_localBound_le_middleRange}
\formallabel{GlobalTheory.NLink.BBSize.main}
\end{proposition}

\begin{proof}
Let $\iota,\pi$ be as in the definition of the error bound. As $\ell\notin\mathrm{small}\cup\{q,d\}$, $\ell$ is split or
inert. In the
inert case $\linkNadmerr(\iota,\pi)=\linkNadmerr_{\mathrm{dec}}(\iota)\le32cfnt^2n_V\Lambda$, in the split case
$\linkNadmerr(\iota,\pi)=\linkNadmerr_{\mathrm{pl}}(\iota,\pi)\le48cfnt^2n_V\Lambda$ (\cref{pf:linkN-adm:blk}). Since
$\Lambda\ge1$, both are at most $60\,cfnt^2n_V\Lambda^2$.
\end{proof}

\begin{lemma}\label{pf:linkN-adm:h0}
Let $\iota$ be a local index with $\ell=\ell_\iota\notin\mathrm{small}\cup\{q,d\}$ and $\mathrm{col}$ a set of places above
$\ell$. Then $h_0(\iota,\mathrm{col})\le2t$.
\formallabel{GlobalTheory.T.maxLength_localInvariants_twisted_fractionModInteger_le}
\formallabel{GlobalTheory.InvariantsKO.main}
\end{lemma}

\begin{proof}
Let $u\in\mathrm{col}$ and $0<e<t$. Since $s_0(\iota)=1$, the inertia group at $u$ contains an element $g$ with
$\varepsilon_\ell(g)=1+\ell$ (\lref{GlobalTheory.PtCore.EpsAll.exists_generator}). As $\ell$ is odd,
$a=\varepsilon_\ell(g)^e-1$ has valuation $(1+v_\ell(e))v(\ell)$ in $\OO_\iota$, and $a\ne0$. The invariants of
$(K_\iota/\OO_\iota)(\varepsilon_\ell^{-e})$ are killed by $\varepsilon_\ell(g)^{-e}-1$, a unit multiple of $a$, so they lie in
$a^{-1}\OO_\iota/\OO_\iota$, of normalized length $1+v_\ell(e)\le e+1\le t\le2t$. Taking the maximum gives the claim (and
$h_0=0$ if $\mathrm{col}=\emptyset$).
\end{proof}

\begin{lemma}\label{pf:linkN-adm:h0W}
Let $\iota$ be a local index, $u\mid\ell_\iota$, and let $T$ be the standard lattice $\OO^r$ of the localization at $u$ of
$\rho_{\iota,\pi}\chi_{\iota,1}^{-1}$, with the coefficients of the datum at $u$. If $\ell^c$ kills the $G_u$-invariants
of $V/T$, then $h^W_0(\iota,u,\pi)\le rc$.
\formallabel{GlobalTheory.T.lenUnit_invariants_twistedLattice_le_of_killingExponent}
\formallabel{GlobalTheory.LengthH0Bound.main}
\end{lemma}

\begin{proof}
Transporting along the inclusion of $\OO_\iota$ into the local coefficient ring, the $G_u$-invariants of
$(K_\iota/\OO_\iota)^r$ are killed by $\ell^c$; hence they lie in $(\ell^{-c}\OO_\iota/\OO_\iota)^r$, whose normalized length is
$rc$.
\end{proof}

\begin{proposition}\label{pf:linkN-adm:sizes}
Assume $\sfname{Hypotheses}_2$ and the setup conditions. Then $X$ satisfies the invariant bounds with constant $120$.
\formallabel{GlobalTheory.T.invariantsBoundNonsplit_of_setup}
\formallabel{GlobalTheory.NSizes.main_alt_alt}
\end{proposition}

\begin{proof}
The bounds for canonical hub data \lref{GlobalTheory.T.split_bound_of_params_alt}, applied with constant $120$ (their
hypotheses are fields of the norm-choice facts, Fontaine's $B_{\mathrm{cris}}$, local finiteness and Chebotarev's theorem),
give $c_B(\iota)\le120\,tnn_V\Lambda$ and, at split $\ell$, $c_0(\iota,u)\le120\,nn_V\Lambda$ for every column $u$; they
also give, for every column $u$ and $\pi\in Z^{\mathrm{good}}_\iota$, Hodge--Tate bounds $(M_0,h_0)$ and a non-jump
condition for the lattice $T$ of \cref{pf:linkN-adm:h0W}, for which the crystalline Wach module exists
(\lref{GlobalTheory.T.crystallineLatticeWachModuleExists_proof}), and an exponent $c'\le c_0(\iota,u)$ attached to
them. By the torsion exponent theorem \lref{GlobalTheory.T.torsion_exponent_kills_invariants_of_good} (from
Bloch--Kato, local duality and the Euler characteristic formula; it states that $\ell^{c'}$ kills the invariants of $V/T$
under these conditions), $\ell^{c_0(\iota,u)}$ kills the $G_u$-invariants of $V/T$, and \cref{pf:linkN-adm:h0W} gives
$h^W_0\le r\,c_0(\iota,u)\le120\,rnn_V\Lambda$. Finally $h_0(\iota,\mathrm{col})\le2t\le120t$ by \cref{pf:linkN-adm:h0}.
\end{proof}

% ---------------------------------------------------------------------------
\subsection{The split step}\label{pf:linkN-adm:splitsec}

\begin{proposition}[The split step]\label{pf:linkN-adm:split}
Assume $\sfname{Hypotheses}_2$ and the setup conditions. Let $\iota$ be twisted with
$\ell=\ell_\iota\notin\mathcal G$ in the middle range and split. Then there is a choice of embedding types for which the
declaration of $X$ at $\iota$ is valid, every place of $E$ above $\ell$ is unramified over $\Q$, and every $\pi\in
Z^{\mathrm{good}}_\iota$ satisfies the thirty-five split admissibility conditions with $e'=X.e_S(\iota)$, $c_0=X.c_0(\iota,
\cdot)$ and $c_B=X.c_B(\iota)$.
\formallabel{GlobalTheory.T.unramified_places_declared'}
\formallabel{GlobalTheory.NLink.splitPrime_hyps_alt}
\end{proposition}

\begin{proof}
The proof collects the conditions. The jump statement \lref{GlobalTheory.JumpStatement.holds_all_alt} provides
the embedding types (it bounds the Hodge--Tate weights of the local representations of the datum); the bounds for
canonical hub data \lref{GlobalTheory.T.split_bound_of_params_alt} provide the valid declaration, the unramifiedness
of the places above $\ell$, the type condition, the column point data (unramified and dual-place conditions, the twisted
dual, the Wach and filtration conditions and finiteness of $H^1_f$) and the block condition;
\lref{GlobalTheory.T.set_pt_of_setup_not_twisted} provides the seventeen conditions of the first kind of
\cref{pf:linkN-adm:def-adm}; \lref{GlobalTheory.T.admG2_of_setup_jumps} and \lref{GlobalTheory.T.admG5_of_setup}
provide the split forms of the gain and tame-shape conditions; the Bloch--Kato annihilator and Tate finiteness conditions
come from \lref{GlobalTheory.T.blochKatoAnnihilator_all_of_tateDuality_eulerChar} and
\lref{GlobalTheory.T.tateFinitenessAll_of_localTateDuality}; the condition (TB) is proved as in the inert step
(\cref{pf:linkN-adm:TB}). Three column conditions are proved in the region of this section:
\begin{itemize}
\item the lift exponent conditions at columns and at their conjugate places
  (\lref{GlobalTheory.T.liftExponentColumn_and_dualPlaces}): the jumps of the localized $r^{\mathrm{tw}}_x$ and
  $r^\vee_x$ lie in $[-\lfloor h_B\rfloor,\lfloor h_B\rfloor]$ by the jump statement, and the lifting theorem in case~A
  with constant $C_A=23$ (part of \lref{GlobalTheory.T.localConstants_from_admitted_local_Galois_statements})
  bounds the lifting exponent by $e_A(\ell)=e_S(\iota)$;
\item the column invariant bound (\lref{GlobalTheory.T.h0_column_length_bound_of_setup}): the invariants of the
  twisted dual lattice are killed by $\ell^{c_0(\iota,u)}$, by the torsion exponent theorem as in
  \cref{pf:linkN-adm:sizes};
\item the ambient conditions (\lref{GlobalTheory.T.ambSub_ambCont_h0_of_setup}): subrepresentations of sums of the
  $r^{\mathrm{tw}}_x$ are crystalline with jumps among those of $r^{\mathrm{tw}}_\pi$ (the twists have equal Hodge--Tate
  multiplicities by the gain condition, and Fontaine's properties apply), the localized families are continuous, and
  the column $H^0$ vanishes.
\end{itemize}
The remaining conditions (characteristic polynomial relations) come from
\lref{GlobalTheory.charpoly_span_and_hom_relations_of_admissible}, which derives them from the first condition (G1).
\end{proof}

% ---------------------------------------------------------------------------
\subsection{Row bounds at inert primes}\label{pf:linkN-adm:rows}

Throughout this subsection $\iota$ is a local index with $\ell=\ell_\iota\ge11$, $\lambda\mid\ell$ a place of $E$ with
$K_\lambda/\Q_\ell$ unramified, $L_C$ an admissible coefficient extension and $W=W_{\lambda,C}$; the number of letters is
$n_\lambda$ (as $K_\lambda$ is unramified, $[K_0:\Q_\ell]=[K_\lambda:\Q_\ell]=n_\lambda$), and $W$ is in \emph{case~A}:
$K_0=K_\lambda$ because $K_\lambda$ is unramified, the level of the Wach structure attached to admissible coefficients is $1$,
and $\ell\ne2$. The row bounds $\mathrm{Row}_{\lambda,C}(x)$ are those of \cref{pf:linkN-adm:starsec}.

\begin{lemma}\label{pf:linkN-adm:rowat}
If $\mathrm{Row}_{\lambda,C}(x)$ holds and $x\le y\in\Q$, then $\mathrm{Row}_{\lambda,C}(y)$ holds.
\formallabel{GlobalTheory.T.rowBound_at_place_of_latticeColengthBound}
\formallabel{GlobalTheory.RowAtOfColLen.main}
\end{lemma}

\begin{proof}
Monotonicity of $\preceq$. (In the formal development the hypothesis is formulated for the lattices of the local
situation at $\lambda$ and the conclusion for arbitrary representations given by the same matrices; the two
formulations quantify over the same lattices and maps.)
\end{proof}

\paragraph{Elementary inequalities.} For $\ell\ge11$ we have $\log\ell\ge1$, hence $\log_\ell y\le\log y\le y-1<y$ for
$y\ge1$ and $\log_\ell y\le y$ for all $y>0$. Consequently, for $j\in\N$ and $H\ge0$,
\begin{equation}\label{pf:linkN-adm:eq-logsplit}
\log_\ell H\le\frac{H}{(\ell-1)^{j+1}}+(j+1),
\end{equation}
since $\log_\ell H=\log_\ell\bigl(H/(\ell-1)^{j+1}\bigr)+(j+1)\log_\ell(\ell-1)$ for $H>0$, and $\log_\ell(\ell-1)\le1$
(for $H=0$ the left side is $0$ by convention).

\begin{lemma}\label{pf:linkN-adm:bracket}
Let $\ell\ge11$, $j\in\N$, and $H,G_0,G,h,\Lambda'\ge0$ with $G_0\le G$ and ($H\le h$ or $H\le G_0$). Then
\[
\frac{H}{(\ell-1)^{j+1}}+(G_0+1)(j+1)+\log_\ell H\le3\Bigl(\frac{h}{(\ell-1)^{j+1}}+(G+1)(j+1)+\Lambda'\Bigr).
\]
\formallabel{GlobalTheory.RowBLong.bracket_le}
\end{lemma}

\begin{proof}
Write $q=(\ell-1)^{j+1}\ge1$. By \eqref{pf:linkN-adm:eq-logsplit} the left side is at most $2H/q+(G_0+1)(j+1)+(j+1)$. If
$H\le h$ this is at most $2h/q+2(G+1)(j+1)$. If $H\le G_0$, then $H/q\le G_0$ and the left side is at most
$2G_0+(j+1)+(G_0+1)(j+1)\le3(G_0+1)(j+1)\le3(G+1)(j+1)$. In both cases the claim follows.
\end{proof}

\begin{proposition}[Row~B]\label{pf:linkN-adm:rowB}
Let $C_2\ge1$ and assume Corollary A with constant $C_2$, (PB), $\ell\ge11$, and
the row~B condition of the joint hub conditions at $(\iota,\pi)$. Then for every declaration at $\iota$, every row place
$\lambda$ and every good $L_C$, $\mathrm{Row}_{\lambda,C}(\linkNadmlen_B(\ell,\lambda))$ holds; that is, condition (Row)
holds with $\linkNadmlen_{\mathrm{row}}=\linkNadmlen_B(\ell,\cdot)$.
\formallabel{GlobalTheory.T.longOrbits_rowColength_inertLocalisation}
\formallabel{GlobalTheory.RowBLong.main}
\formallabel{GlobalTheory.Survey.InertPlacesGoodSet.rowL_of_wordsB}
\end{proposition}

\begin{proof}
By the joint condition $K_\lambda$ is unramified; a good $L_C$ is admissible, and there is a word based at
$(\iota,\pi,\lambda,C)$ satisfying (B$_{C_2}$); put $j=j_A(\ell)$. Let $T_1,T_2,F$ be as in the definition of
$\mathrm{Row}_{\lambda,C}$. By the word, $T_1$ and $T_2$ are crystalline, the jumps of $T_1$ lie in the windows, every letter is
thick or has size at most $w$, and the hypotheses $1\le j$, $j+1\le n_\lambda$, $C_2(j+1)(w+S+\eta+1)\le b$ and the
cluster condition hold. Corollary A \eqref{pf:linkN-adm:eq-corA} gives
\[
\linkNadmcol(F)\preceq C_2(n_\lambda+1)\,r\Bigl(\frac{h_{\mathrm{cor}}}{(\ell-1)^{j+1}}+(w+S+\eta+1)(j+1)
+\log_\ell h_{\mathrm{cor}}\Bigr).
\]
Apply \cref{pf:linkN-adm:bracket} with $H=h_{\mathrm{cor}}$, $G_0=w+S+\eta\le w_P+S_c+\eta_3=G_P$, $h=h_B$ and
$\Lambda'=\Lambda$: if some letter is thick, $h_{\mathrm{cor}}=b+\eta+S\le h_B$ by the word; otherwise
$h_{\mathrm{cor}}=G_0$. Hence $\linkNadmcol(F)\preceq3C_2\,\mathrm{core}_B(\ell,\lambda)$, that is $\mathrm{Row}_{\lambda,C}(3C_2
\mathrm{core}_B)$. If $\mathrm{core}_B>0$ then $3C_2\mathrm{core}_B<\linkNadmlen_B(\ell,\lambda)$; otherwise
$3C_2\mathrm{core}_B\le0=\linkNadmlen_B(\ell,\lambda)$. \Cref{pf:linkN-adm:rowat} concludes.
\end{proof}

\paragraph{Two local inputs.}
The rows~C/D use theorems~B and~C below with the constants $C=320$ and $C=25$, and the following consequence of the
local constants \lref{GlobalTheory.T.localConstants_from_admitted_local_Galois_statements} (which, under
Bloch--Kato's theory, local Tate duality and the Euler characteristic formula, proves the local lifting theory with
constants $(10,23,32,32)$): \emph{in case~A, if $T_1$ is crystalline with jumps within $(M^-,M^+)$ at every letter, then}
\begin{equation}\label{pf:linkN-adm:eq-ePA}
e_{PA}\bigl(\ell,M_{\max},H_{\max},J\bigr)\le23\Bigl(1+\frac{H_{\max}}{\ell}+1+\log_\ell(H_{\max}+1)\Bigr),
\end{equation}
where $M_{\max}=\max_aM^-_a$, $H_{\max}=\max_a(M^-_a+M^+_a)$, $J$ is the set of jumps of $T_1$ in $[-M_{\max},\max_aM^+_a]$, and
$e_{PA}(\ell,M,H,J)=2\lfloor H\ell/(\ell-1)^2+3+\log_\ell(H+1)\rfloor+\sum_{j\in J'}v_\ell(j)+a+v_\ell(a!)$ with
$a=\lfloor M/(\ell-1)\rfloor+1$ and $J'=\{j+k:j\in J,\ 1\le k\le M,\ j+k\ne0\}\cup[1,M]$. The number $e_{PA}$ is the
exponent of the lifting theorem in case~A: an onto map $F\colon T_1\to T_2$ of crystalline lattices lifts with exponent
$e_{PA}$. For a word we write $e_{\mathrm{lift}}=e_{PA}(\ell,M_{\max},H_{\max},J)$.

\begin{lemma}\label{pf:linkN-adm:numNT}
Let $\ell\ge11$ and let $n,r,m,h,w,\eta,S,\mathrm{lg},\beta_0,\beta,e',h_B,w_P,\eta_3,S_c,\Lambda,c_E$ be reals with
$n,r,m,w,\eta,S,h\ge0$, $h\le h_B$, $w\le w_P$, $\eta\le\eta_3$, $S\le S_c$, $\mathrm{lg}\le3+\Lambda$, $\Lambda\ge1$,
$\beta_0\le\beta$, $\beta\ge0$, $e'\le69c_E$ and $c_E\ge0$. Then
\begin{multline*}
320(n+1)r\Bigl(\frac{h}{\ell^2}+\frac{w+\eta}{\ell}+S+1+\mathrm{lg}+\eta\beta_0\Bigr)+320\,nm\,e'\\
\le368(n+1)r\frac{h_B}{\ell^2}+2082(n+1)r\Bigl(\frac{w_P+\eta_3}{\ell}+S_c+\Lambda+\eta_3\beta\Bigr)+22365\,nm\,c_E.
\end{multline*}
\formallabel{GlobalTheory.NLink.RowNT.bound_real}
\end{lemma}

\begin{proof}
We have $h/\ell^2\le h_B/\ell^2$, $(w+\eta)/\ell\le(w_P+\eta_3)/\ell$, $\eta\beta_0\le\eta_3\beta$ (if $\beta_0<0$ the
left side is $\le0$), and $S+1+\mathrm{lg}\le S_c+4+\Lambda\le S_c+5\Lambda$. Hence the bracket on the left is at most
$h_B/\ell^2+6\bigl((w_P+\eta_3)/\ell+S_c+\Lambda+\eta_3\beta\bigr)$. Multiply by $320(n+1)r\ge0$ and use $320\le368$,
$6\cdot320=1920\le2082$, and $320\cdot69=22080\le22365$ for the last term.
\end{proof}

\begin{proposition}[Row~C: words with a non-thick letter]\label{pf:linkN-adm:rowNT}
Assume theorem~B with constant $320$ (\cref{pf:linkN-adm:thmB}), the bound \eqref{pf:linkN-adm:eq-ePA},
(PB), $\ell\ge11$ and $\pi\in Z^{\mathrm{good}}_\iota$. If a word based at
$(\iota,\pi,\lambda,C)$ satisfies (Box) and (NT), then
$\mathrm{Row}_{\lambda,C}\bigl(\mathrm{bd}_{CD}(\mathrm{band}(\iota,\lambda),\ell,\lambda;368,2082,22365)\bigr)$ holds.
\formallabel{GlobalTheory.T.rowBounds_of_nonThick_letter}
\formallabel{GlobalTheory.NLink.RowNT.main}
\formallabel{GlobalTheory.Survey.InertPlacesGoodSet.setting_of_wordsNT}
\formallabel{GlobalTheory.NLink.RowNT.lift_shape_le}
\formallabel{GlobalTheory.NLink.RowNT.logb_le_of_le_four_exp}
\end{proposition}

\begin{proof}
Let $T_1,T_2,F$ be as in the definition of $\mathrm{Row}_{\lambda,C}$. The word provides the \emph{setting for theorem~B} of \cref{pf:linkN-adm:thmB} for
$(W,T_1,T_2,F)$ and the word's parameters: case~A, crystallinity of $T_1,T_2$ and ontoness of $F$, $\eta<\mathrm{sz}_a$,
the jumps of $T_1$ in the windows and within $(M^-,M^+)$, the low counts of $T_2$ at most $m$, and $\mathrm{sz}_a\le w$ at
non-thick letters. Together with $20w,20\eta,20S\le b$, a non-thick letter, and the alternation of the type, part~(a)
of theorem~B gives
\[
\linkNadmcol(F)\preceq320(n_\lambda+1)r\Bigl(\frac{h}{\ell^2}+\frac{w+\eta}{\ell}+S+1+\mathrm{lg}+\eta\beta_0\Bigr)
+320\,n_\lambda m\,e_{\mathrm{lift}},
\]
with $h$ the height of the word, $\mathrm{lg}=\log_\ell(2h)$ and $\beta_0\in\{0,1\}$ the band indicator of the word.
We check the hypotheses of \cref{pf:linkN-adm:numNT} with $\beta=[\mathrm{band}(\iota,\lambda)]$ and
$c_E=\mathrm{core}_E(\ell)=1+h_B/\ell+\log_\ell h_B$.
(i) $h\le h_B$: if a letter is thick, $h=b+\eta+S\le h_B$; otherwise $h=w+S\le b+\eta+S\le h_B$ as $20w\le b$.
(ii) $\mathrm{lg}\le3+\Lambda$: $2h\le2h_B\le4e^\Lambda$, so $\log_\ell(2h)\le\log(4e^\Lambda)\le3+\Lambda$.
(iii) $\beta_0\le\beta$: if the band condition of the word holds, the word itself witnesses $\mathrm{band}(\iota,\lambda)$
(as $\ell\ge11$ and $\pi\in Z^{\mathrm{good}}_\iota$).
(iv) $e_{\mathrm{lift}}\le69\,c_E$: by \eqref{pf:linkN-adm:eq-ePA} and $H_{\max}\le2h_B$ (from (Box)),
$e_{\mathrm{lift}}\le23(1+2h_B/\ell+1+\log_\ell(3h_B))\le23(3+2h_B/\ell+\log_\ell h_B)\le69c_E$, using $h_B\ge1$ and
$\log_\ell3\le1$.
The remaining hypotheses are part of the word ($w\le w_P$, $\eta\le\eta_3$, $S\le S_c$) and $\Lambda\ge1$. The right side
of \cref{pf:linkN-adm:numNT} is $\mathrm{bd}_{CD}(\beta,\ell,\lambda;368,2082,22365)$ by \eqref{pf:linkN-adm:eq-bdCD}, with
$n=n_\lambda$ and the rank $r$ of $T_2$ and the rank $m$ of $D$.
\end{proof}

\begin{lemma}\label{pf:linkN-adm:numAT}
Let $\ell\ge11$ and let $n,r,m,S,\eta,H,H',e,e',h_B,\eta_3,w_P,S_c,\Lambda,\mathrm{flt}$ be reals with $n,r,m,S,\eta,H,H'\ge0$,
$S\le S_c$, $\eta\le\eta_3\le h_B$, $H\le h_B+\eta_3$, $1\le h_B\le2e^\Lambda$, $\Lambda\ge1$, $w_P\ge0$,
$(w_P+\eta_3)/\ell+S_c+\Lambda\le\mathrm{flt}$, $H'\le H$,
\[
e\le\frac{32H}{\ell^2}+4S+\frac{10\eta}{\ell-1}+25\bigl(\log_\ell(2H)+1\bigr),\qquad
e'\le23\Bigl(1+\frac{H'}{\ell}+1+\log_\ell(H'+1)\Bigr).
\]
Then, with $c_D=(H+1)/(\ell-1)+1$ and $c_E=1+h_B/\ell+\log_\ell h_B$,
\[
(n+1)re+4nme'+\tfrac32n\bigl(rS+2mc_D\bigr)\le368(n+1)r\frac{h_B}{\ell^2}+2082(n+1)r\,\mathrm{flt}+22365\,nm\,c_E.
\]
\formallabel{GlobalTheory.RowsThick.colBd_le_row}
\end{lemma}

\begin{proof}
Since $2H\le2(h_B+\eta_3)\le4h_B\le8e^\Lambda$ and $\log8\le7$, $\log_\ell(2H)\le7+\Lambda$. Since $H'+1\le H+1\le3h_B$ and
$3\le\ell$, $\log_\ell(H'+1)\le1+\log_\ell h_B$. Next $32H/\ell^2\le32h_B/\ell^2+3\eta_3/\ell$ (as $\eta_3/\ell^2\le
\eta_3/(11\ell)$) and $10\eta/(\ell-1)\le11\eta_3/\ell$ (as $\ell\ge11$). Hence, with $E_1=h_B/\ell^2$,
$e\le32E_1+14\eta_3/\ell+4S_c+25(8+\Lambda)\le32E_1+225\,\mathrm{flt}$, using $200\le200\Lambda$. Also $H'/\ell\le2h_B/\ell$
gives $e'\le23(3+2h_B/\ell+\log_\ell h_B)\le69c_E$, and $(H+1)/(\ell-1)\le3h_B/(\ell-1)\le4h_B/\ell$ gives $c_D\le4c_E$.
Finally $nrS\le(n+1)r\,\mathrm{flt}$. Therefore the left side is at most
\[
(n+1)r(32E_1+225\,\mathrm{flt})+276\,nmc_E+\tfrac32(n+1)r\,\mathrm{flt}+12\,nmc_E,
\]
which is at most the right side.
\end{proof}

\begin{proposition}[Row~D: words all of whose letters are thick]\label{pf:linkN-adm:rowAT}
Assume theorem~C with constant $25$ (\cref{pf:linkN-adm:thmC}), the bound \eqref{pf:linkN-adm:eq-ePA},
(PB) and $\ell\ge11$. If a word based at $(\iota,\pi,\lambda,C)$ satisfies
(Box), (Mult) and (AT$_d$), then for every proposition $\beta$,
$\mathrm{Row}_{\lambda,C}\bigl(\mathrm{bd}_{CD}(\beta,\ell,\lambda;368,2082,22365)\bigr)$ holds.
\formallabel{GlobalTheory.T.rowBounds_of_all_letters_thick}
\formallabel{GlobalTheory.RowsThick.rowL_of_allThick}
\formallabel{GlobalTheory.Survey.InertPlacesGoodSet.setCS_of_wordsAT}
\end{proposition}

\begin{proof}
Let $T_1,T_2,F$ be as in the definition of $\mathrm{Row}_{\lambda,C}$. The word provides a setting of parity $d$ for theorem~C (\cref{pf:linkN-adm:def-setC}) for
$(W,d,T_1,T_2,F,\mathrm{sz},S,\eta,m)$: case~A, $\ell\ge11$, $n_\lambda$ even, $\OO_C$ a discrete valuation ring with fraction
field $L_C$, a primitive $\ell$-th root of unity in $L_C$, crystallinity and ontoness, the multiplicity bound
$\sum_{|x|\le\eta}\mathrm{mult}_\tau(T_2,x)\le m$, and the two size conditions; and the jumps of $T_1$ lie in the windows
of theorem~C, because by (AT$_d$) the type is $\mathrm{up}(a)\Leftrightarrow a+d$ odd, so the word's window at $a$
is $[-\mathrm{sz}_a-S,-\mathrm{sz}_a]\cup[-\eta,\eta]$ if $a+d$ is even and $[\mathrm{sz}_a,\mathrm{sz}_a+S]\cup[-\eta,\eta]$
otherwise. Theorem~C gives $e\in\N$ with $e\le e_{\mathrm{Bd}}$ and
\[
\linkNadmcol(F)\preceq(n_\lambda+1)re+4n_\lambda m\,e''+\tfrac32n_\lambda\bigl(rS+2mc_D(\ell,H)\bigr),
\]
where $H=H_0+\eta$, $H_0=\max_a\mathrm{sz}_a+S$, $e''$ is the lifting exponent $e_{PA}$ for the windows of theorem~C, and
$e_{\mathrm{Bd}}$ is the bound of \cref{pf:linkN-adm:numAT}. Since the jumps of $T_1$ lie within the windows,
\eqref{pf:linkN-adm:eq-ePA} bounds $e''$ with $H'=H_{\max}\le H_0+\eta=H$. As all letters are thick,
$\max_a\mathrm{sz}_a\le b+\eta$, so $H\le b+\eta+S+\eta\le h_B+\eta_3$ by the word; moreover $\eta_3\le h_B$, $h_B\ge1$,
$h_B\le2e^\Lambda$, $\Lambda\ge1$, $w_P\ge0$, and $\mathrm{flt}=(w_P+\eta_3)/\ell+S_c+\Lambda+\eta_3[\beta]$ satisfies the
hypothesis of \cref{pf:linkN-adm:numAT}. That lemma, with $n=n_\lambda$ and $r,m$ the ranks, gives the claim.
\end{proof}

\begin{proposition}[Rows~C and~D]\label{pf:linkN-adm:rowCD}
Assume theorem~B with constant $320$, theorem~C with constant $25$, the local constants (so that
\eqref{pf:linkN-adm:eq-ePA} holds), (PB), $\ell\ge11$, $\pi\in
Z^{\mathrm{good}}_\iota$, and the rows~C/D condition of the joint hub conditions at $(\iota,\pi)$. Then for every declaration
at $\iota$, every row place $\lambda$ and every good $L_C$, $\mathrm{Row}_{\lambda,C}(\linkNadmlen_{CD}(\ell,\lambda))$
holds, where
$\linkNadmlen_{CD}$ is formed with the band predicate and $(K_1,K_2,K_3)=(368,2082,22365)$.
\formallabel{GlobalTheory.T.rowsCD_rowCondition_localBounds}
\formallabel{GlobalTheory.NLink.lengthBoundAtPlaces_of_jointCond}
\formallabel{GlobalTheory.NLink.RowsCD.bdCD_lt}
\end{proposition}

\begin{proof}
By the joint condition $K_\lambda$ is unramified, a good $L_C$ is admissible, and there is a word based at
$(\iota,\pi,\lambda,C)$ satisfying (Box) and either (NT) or (Mult) and
(AT$_d$). By \cref{pf:linkN-adm:rowNT}, resp.\ \cref{pf:linkN-adm:rowAT} with
$\beta=\mathrm{band}(\iota,\lambda)$, $\mathrm{Row}_{\lambda,C}(\mathrm{bd}_{CD}(\beta;368,2082,22365))$ holds. The coefficient of
$K_2$ in \eqref{pf:linkN-adm:eq-bdCD} is $(n_\lambda+1)r(\dots)>0$, since $r\ge k\ge5$ and $(w_P+\eta_3)/\ell+S_c+\Lambda>0$;
hence $\mathrm{bd}_{CD}(\beta;K_1,K_2,K_3)<\mathrm{bd}_{CD}(\beta;K_1,K_2+1,K_3)$ and $\linkNadmlen_{CD}(\ell,\lambda)$ lies
strictly between them. \Cref{pf:linkN-adm:rowat} concludes.
\end{proof}

% ---------------------------------------------------------------------------
\subsection{The inert step}\label{pf:linkN-adm:inertsec}

\begin{lemma}\label{pf:linkN-adm:eleven}
If $\ell$ is a prime, inert for $D$, and the split bound is at least $10$, then $\ell\ge11$.
\formallabel{GlobalTheory.Survey.InertPlacesGoodSet.eleven_le}
\end{lemma}

\begin{proof}
An inert prime is not small, so $\ell\ne2$, and $\bigl(\frac{-d}{\ell}\bigr)\ne1$. If $\ell\le10$, then $\ell$ is an odd
prime at most the split bound, so $\bigl(\frac{-d}{\ell}\bigr)=1$ by the host axioms; a contradiction.
\end{proof}

\begin{lemma}[The Bloch--Kato condition (TB)]\label{pf:linkN-adm:TB}
Let $\iota$, $r_\iota$, $\pi$, $M_0,h_0,c_B\in\N$. Assume: the local theorems ``TB'' and ``TBe'' for all primes, unramified
fields and admissible coefficients; the existence of Wach modules for crystalline
lattices with jumps in $[-M_0,h_0]$; the condition (CrysOn) that crystalline lattices are crystalline on the subgroup of $G_\lambda$ on which the cyclotomic character is $\equiv1\bmod\ell^s$ ($s$ the level of the Wach structure); Fontaine's properties at
the places above $\ell_\iota$; and the \emph{joint tame condition} with $(M_0,h_0,c_B)$: $\ell\ne2$, and at every
$\lambda\mid\ell$ the inertia subgroup of the datum is the inertia group of $K_\lambda$, $K_\lambda\subseteq
K_0(\mu_\ell)$, and for every admissible $L_C$ and every crystalline lattice $Y$ whose jumps are among those of
$r^{\mathrm{tw}}_\pi$, the jump set of $Y$ lies in $[-M_0,h_0]$ and the constants $c_{TB}(Y,M_0,h_0,e)$, $e<t$, are at
most $c_B$. Then condition (TB) holds at $(\iota,r_\iota,\pi,c_B)$.
\formallabel{GlobalTheory.T.localBlochKatoCondition_of_wachModules}
\formallabel{GlobalTheory.TBLocal.main}
\end{lemma}

\begin{proof}
Let $\lambda\mid\ell$ and $L_C$ good; then $L_C$ is admissible. For the first part of (TB), let $Y$, the invariant
functional $p$ with $p(e_0)=1$, $0\le e<t$, the lattice of $L_C(\varepsilon^{-e})$ and the map $\varphi$ be given. The kernel
$Y'=Y\cap\ker p$ is a lattice of a subrepresentation, crystalline with jumps among those of $Y$ by Fontaine's
properties; the joint tame condition applies to $Y'$, the Wach hypothesis provides its Wach module, and
(CrysOn) applies to $Y$. Then \lref{GlobalTheory.TBLocal.tb_part1} (the local theorem TB applied to
these data) shows that $p^{c_B}$ times the cocycle $g\mapsto\varphi(ge_0-e_0)$ is finite.
For the second part, \lref{GlobalTheory.TBLocal.tbe_clauses} derives the three clauses for each $e<t$ from the local
theorem TBe and the equality of inertia groups.
\end{proof}

\begin{theorem}[The inert step]\label{pf:linkN-adm:inert}
Under the hypotheses of \cref{pf:linkN-adm:main}, let $\iota$ be twisted with $\ell=\ell_\iota\notin\mathcal G$ in the middle
range and inert, and $\pi\in Z^{\mathrm{good}}_\iota$. Then the point of $X$ at $(\iota,\pi)$ is admissible.
\formallabel{GlobalTheory.T.adm_at_inert_of_range}
\formallabel{GlobalTheory.NLink.InertStep.main_alt_alt}
\end{theorem}

\begin{proof}
\emph{The seventeen conditions of the first kind} are given by
\lref{GlobalTheory.T.set_pt_of_setup_not_twisted}, which proves them from $\sfname{Hypotheses}_2$ and the setup for
every $\iota$ with $\ell_\iota\notin\mathcal G$ and every $\pi\in Z^{\mathrm{good}}_\iota$.

\emph{Condition (TB).} The closing statement for the canonical hub data,
\lref{GlobalTheory.closingFamilyX_alt_alt} (from Fontaine's $B_{\mathrm{cris}}$, local finiteness, the jump statement
\lref{GlobalTheory.JumpStatement.holds_all_alt}, the admissible parameters and the facts of the setup, and
$c_B(\iota)$ being at least the canonical value), gives $M_0,h_0$ for which the joint tame condition holds with
$(M_0,h_0,c_B(\iota))$ at $\iota$ in the middle range. The local theorems TB and TBe are
\lref{GlobalTheory.T.theorem_tb_of_local_blochKato_duality_euler} (from Bloch--Kato, local duality and the Euler
characteristic), the existence of Wach modules is \lref{GlobalTheory.T.crystallineLatticeWachModuleExists_proof},
(CrysOn) is \lref{GlobalTheory.crysOnCond_all}, and Fontaine's properties are among the conditions of
the first kind. \Cref{pf:linkN-adm:TB} gives (TB).

\emph{Condition (Row).} The split bound is at least $60$, so $\ell\ge11$ by \cref{pf:linkN-adm:eleven}. Since $X$ is
adapted, $X.\linkNadmlen_{\mathrm{row}}(\iota,\cdot)$ is $\linkNadmlen_B(\ell,\cdot)$ if $\ell$ is of row~B and
$\linkNadmlen_{CD}(\ell,\cdot)$ otherwise, the latter formed with the band predicate $X.\mathrm{band}(\iota,\cdot)$ and
$\mathcal K(C_2)$. In row~B, \cref{pf:linkN-adm:rowB} applies (with Corollary A, $C_2\ge1$, the parameter bound of the
setup and the row~B joint condition). In rows~C/D, \cref{pf:linkN-adm:rowCD} applies, with theorem~B for $C=320$
(\cref{pf:linkN-adm:thmB}), theorem~C (\cref{pf:linkN-adm:thmC}), the local constants
\lref{GlobalTheory.T.localConstants_from_admitted_local_Galois_statements}, and the rows~C/D joint condition.

\emph{Condition (Lift)} is vacuous: since $X$ is adapted and $\ell$ is inert, the declaration has no columns.
\end{proof}

% ---------------------------------------------------------------------------
\subsection{Proof of the main theorem}\label{pf:linkN-adm:proofmain}

\begin{proof}[Proof of \cref{pf:linkN-adm:main}]
Let the data be as in the theorem. By the setup, $X$ is adapted to $L$ and $\mathcal K(C_2)$, its families are good, and
the killing condition and (Sha) hold. The step conditions hold: part (a) is the inert step
(\cref{pf:linkN-adm:inert}), part (b) the split step (\cref{pf:linkN-adm:split}), part (c) a field of the setup. The
integrality-from-metric theorem \eqref{pf:linkN-adm:eq-ifm} for $D$ is
\lref{GlobalTheory.T.integralityFromMetric_of_designConditionAll}, applied with global Artin reciprocity, Bloch--Kato's
local theory, local Tate duality and the Euler characteristic formula (fields of $\sfname{Hypotheses}_2$) and the design
condition among the norm-choice facts. The split-case bound \eqref{pf:linkN-adm:eq-split} is
\lref{GlobalTheory.T.splitCase_bound_at_point_of_admitted_items_alt}, applied with Artin reciprocity, local duality and
the Euler characteristic formula. The gain condition (G2) at the points $(\iota,\pi)$ with $\ell_\iota\notin\mathcal G$ and
$\pi\in Z^{\mathrm{good}}_\iota$ is one of the conditions of \lref{GlobalTheory.T.set_pt_of_setup_not_twisted}. Hence
\cref{pf:linkN-adm:star} gives the star inequality for $X$. By \cref{pf:linkN-adm:sizes} the hub data satisfy the invariant bounds
with constant $120$, so \cref{pf:linkN-adm:bbmain} gives the error bound with constant $60\cdot120=7200$.
\end{proof}

% ---------------------------------------------------------------------------
\subsection{The local lifting theorem B}\label{pf:linkN-adm:thmBsec}

\begin{remark}[Wach structures]\label{pf:linkN-adm:wach}
For a prime $p$, a finite extension $K/\Q_p$ whose maximal unramified subfield $K_0$ has degree $n$, a coefficient
field $L_C$ and a level $s\ge1$, the local Hodge--Tate data of the datum provide a \emph{Wach structure} $W$. It consists
of: the ring $A=\OO_C[[X]]^{\Z/n}$ (one factor for each letter) with the Frobenius $\varphi$ acting by $X\mapsto
Xq_0(X)$, $Xq_0(X)=(1+X)^p-1$, on each factor and shifting the letters; a quotient $\Gamma$ of $G_K$ acting on $A$
compatibly with $\varphi$, with $\gamma(X)\equiv\chi(\gamma)X\bmod X^2$; an $A$-algebra $\mathcal E$ in which $X$ is
invertible, $p$-adically complete and dense; for every lattice $T$ an \'etale $(\varphi,\Gamma)$-module $D(T)$ over
$\mathcal E$, free of rank $\rank T$, functorial in $T$, and a finitely generated $\Gamma$-stable $A$-submodule $N(T)$ with
$(\gamma-1)N(T)\subseteq XN(T)$, stable under $\varphi$ up to powers of $q=\varphi^{s-1}(q_0)$, which is a basis of $D(T)$
when $T$ is crystalline (the \emph{Wach module}, \cite{Wach1996,Berger2004}); maps from Herr cocycles of $D(T)$ to lattice
cohomology \cite{Herr1998}; a comparison map $N(T)\to D_{\mathrm{cris}}(V)$; and level maps to $V\otimes B_{\mathrm{dR}}$
attached to $p$-power roots of unity. All these are data of the main theorem with stated axioms; theorems~B and~C use
them only through these axioms and the admissibility inputs of \cref{pf:linkN-adm:admsec}. \emph{Case~A} means
$K=K_0$, $s=1$ and $p\ne2$.
\end{remark}

For a word $(b,w,S,\eta,m,\mathrm{up},\mathrm{sz},M^-,M^+)$ on the letters of $W$ (\cref{pf:linkN-adm:def-words}), its
\emph{height} is $h=b+\eta+S$ if some letter is thick and $h=w+S$ otherwise, and we put
\begin{equation}\label{pf:linkN-adm:eq-Sg}
\mathrm{lg}=\log_\ell(2h),\qquad \mathrm{Sg}=\frac h{\ell^2}+\frac{w+S+\eta}{\ell}+\mathrm{lg}+1 .
\end{equation}
Lattices $T_1,T_2$, a lattice map $F$ and a word form a \emph{setting for theorem~B} if: $W$ is in case~A; $T_1,T_2$ are crystalline
and $F$ is onto; $\eta<\mathrm{sz}_a$ for all letters; the jumps of $T_1$ lie in the windows and within $(M^-,M^+)$; the low
count of $T_2$ at every letter is at most $m$; and $\mathrm{sz}_a\le w$ at every non-thick letter.

\begin{theorem}[Theorem~B]\label{pf:linkN-adm:thmB}
Assume Bloch--Kato's local theory, local Tate duality and the local Euler characteristic formula. Let $D$ be a datum with
local Hodge--Tate, Fontaine and comparison data, and $C\ge320$. Let $\ell\ge11$ be a prime, $K/\Q_\ell$ finite and
unramified with $n=[K:\Q_\ell]$ letters, $L_C$ admissible for $K$, $W$ the Wach structure, and let $T_1,T_2$, $F$ and a word
form a setting for theorem~B with $20w\le b$, $20\eta\le b$, $20S\le b$ and some letter not thick. Let $\beta\in\{0,1\}$ be the band
indicator of the word with constant $C$, $r=\rank T_2$ and $e'=e_{\mathrm{lift}}$ the lifting exponent of the word. Then:
\begin{enumerate}[label=(\alph*)]
\item if the type alternates and no three consecutive letters are thick,
\[
\linkNadmcol(F)\preceq C(n+1)r\Bigl(\frac h{\ell^2}+\frac{w+\eta}{\ell}+S+1+\mathrm{lg}+\eta\beta\Bigr)+C\,nm\,e';
\]
\item for every $R\ge1$, if the type is constant and no $R$ consecutive letters are thick,
\[
\linkNadmcol(F)\preceq C(n+1)r\Bigl(\frac h{\ell^2}+\frac{w+\eta}{\ell}+R^2(S+1+\mathrm{lg})+\eta\beta\Bigr)+C\,nm\,R^2e'.
\]
\end{enumerate}
\formallabel{GlobalTheory.T.localLiftingTheoremB_of_blochKato_tateDuality_eulerChar}
\formallabel{GlobalTheory.theoremB_closing}
\end{theorem}

\paragraph{Reduction to the admissible inputs.}
Let $\mathcal A_B$ denote the admissibility package of theorem~B with the constants $(10,23,32,32,8)$; its components are
listed in \cref{pf:linkN-adm:def-AB}. The \emph{jump inputs} $\mathcal J$ are the one-place Fontaine inequality (Fontaine's
inequality $\dim_{L_C}D_{\mathrm{cris}}(V)\le[K_0:\Q_\ell]\dim_{L_C}V$ at every place), the exactness of jumps (the statement
(cj) of \cref{pf:linkN-adm:def-cores} at every place) and the global letter condition (\cref{pf:linkN-adm:def-letter}).
Among the components of $\mathcal A_B$ used below are:
\begin{itemize}
\item[(gen)] the \emph{generator bound}: if $K$ is unramified and $T_2$ is crystalline, $\linkNadmHf(T_2)$ is generated by
  at most $([K:\Q_\ell]+1)\rank T_2$ classes (an integral analogue of the bound $\dim\linkNadmHf(V)\le([K:\Q_\ell]+1)\dim V$
  given by the dimension formula of \cite{BlochKato1990});
\item[(lift)] \emph{the lifting theorem in case~A}: if $T_1,T_2$ are crystalline, $F$ is onto and the jumps of $T_1$ lie
  within $(M^-,M^+)$, then $F$ lifts with exponent $e_{PA}(\ell,M_{\max},H_{\max},J)$;
\item[(cores)] the seventeen statements of the bounds package $\mathcal B$, which together with the jump inputs give the
  core local statements at every place in case~A with $\ell\ge11$ (\cref{pf:linkN-adm:cores}).
\end{itemize}

\begin{proposition}[Reduction of theorem~B]\label{pf:linkN-adm:thmBred}
Let $C\ge320$ and assume the jump inputs. Then $\mathcal A_B$ implies the conclusion of \cref{pf:linkN-adm:thmB} with
constant $C$, for all primes $\ell$, unramified $K$ and admissible $L_C$.
\formallabel{GlobalTheory.T.theoremB_of_localFilteredHypotheses}
\formallabel{GlobalTheory.theoremB_of_admitted_inputs}
\formallabel{Formalism.MainNotions.N.Ln17.frobenius_statement_of_fontaine}
\formallabel{Formalism.MainNotions.N.Ln17.stmt_holds}
\formallabel{Formalism.MainNotions.N.Ln17.frobenius_statement_of_inputs}
\end{proposition}

\begin{proof}
Extract from $\mathcal A_B$ the seventeen statements of the bounds package $\mathcal B$ (\cref{pf:linkN-adm:cores}) used in
\cref{pf:linkN-adm:thmBat} and apply that proposition,
with the jump inputs, at every prime $\ell$, unramified $K$ and admissible $L_C$.
\end{proof}

\begin{proof}[Proof of \cref{pf:linkN-adm:thmB}]
The package $\mathcal A_B$ is proved from Bloch--Kato's local theory, local duality and the Euler characteristic formula in
\cref{pf:linkN-adm:admB}. The jump inputs are
\lref{GlobalTheory.FontaineIneq.one_place_admissible_of_gap_constant_le} (the one-place Fontaine inequality for all
fields and coefficients), \lref{GlobalTheory.jumps_exact_adm} (exactness of the Hodge--Tate jumps in short exact
sequences of crystalline lattices) and the global letter condition, which \cref{pf:linkN-adm:letter} derives from the norm
bound for $\Delta$-invariant lattices proved in \lref{GlobalTheory.T.wachLattice_laterClauses_hold}. Apply
\cref{pf:linkN-adm:thmBred}.
\end{proof}

The rest of this subsection proves the one-place statement. It bounds $\linkNadmcol(F)$ by the counting principle
(\cref{pf:linkN-adm:colexp}): the number of generators of $\linkNadmHf(T_2)$ is at most $(n+1)r$, and the exponents come from
a decomposition of the letters into \emph{arcs}.

\begin{proposition}[Theorem~B at one place]\label{pf:linkN-adm:thmBat}
Assume the bounds package $\mathcal B$ and the jump inputs. Then the conclusion of
\cref{pf:linkN-adm:thmB} holds for every $C\ge320$ at every prime $\ell$, unramified $K$ and admissible $L_C$.
\formallabel{Formalism.MainNotions.N.Ln17.thmB_at}
\end{proposition}

\begin{proof}
As $\OO_C$ is a discrete valuation ring and $\ell\ne0$ in $\OO_C$, the counting principle applies. By (gen) applied to $T_2$,
$\linkNadmHf(T_2)$ is generated by at most $(n+1)r$ classes; by (lift), $F$ lifts with exponent $e'$. The height satisfies
$\mathrm{lg}\ge0$, and \eqref{pf:linkN-adm:eq-Sg} gives
\begin{equation}\label{pf:linkN-adm:eq-Sgle}
\mathrm{Sg}\le X_0:=\frac h{\ell^2}+\frac{w+\eta}{\ell}+S+1+\mathrm{lg},
\end{equation}
since $S/\ell\le S$.

(a) Assume the type alternates and no three consecutive letters are thick.
\emph{First case: some letter is thick and $b/4<40\,\mathrm{Sg}$.} By \cref{pf:linkN-adm:thick}, $F$ lifts with an exponent
$e_0\le160\,\mathrm{Sg}$. The counting principle with $Z=\emptyset$ and $e=e_0$ gives $\linkNadmcol(F)\preceq(n+1)re_0$,
and $e_0\le160X_0\le C(X_0+\eta\beta)$ as $C\ge320$; the term $Cnme'$ is non-negative.
\emph{Second case: no letter is thick, or $40\,\mathrm{Sg}\le b/4$.} By \cref{pf:linkN-adm:sumalt} there are $e\in\N$ and a
finite set $Z$ of finite cocycles with $e\le55\,\mathrm{Sg}+13S+\eta\beta$, $|Z|\le4nm$ and
$\ell^e\linkNadmZf(T_2)\subseteq\linkNadmIm(F)+\OO_CZ$. The counting principle gives $\linkNadmcol(F)\preceq(n+1)re+|Z|e'$.
Now $e\le55X_0+13S+\eta\beta\le69(X_0+\eta\beta)\le C(X_0+\eta\beta)$ and $|Z|\le4nm\le Cnm$.

(b) Assume the type is constant and no $R$ consecutive letters are thick. By \cref{pf:linkN-adm:sumalt} there are $e$ and
$Z$ with $e\le18\,\mathrm{Sg}+3R^2(S+1+\mathrm{lg})+1$, $|Z|\le(R+1)nm$ and the same exponent property. Put
$Y=h/\ell^2+(w+\eta)/\ell+R^2(S+1+\mathrm{lg})+\eta\beta\ge1$. Since $\mathrm{Sg}\le X_0\le Y$, $e\le21Y+1\le22Y\le CY$; and
$R+1\le2R^2\le CR^2$. The counting principle gives $\linkNadmcol(F)\preceq(n+1)re+|Z|e'\le C(n+1)rY+CnmR^2e'$.
\end{proof}

\begin{definition}[Arcs]\label{pf:linkN-adm:def-arcs}
An \emph{arc datum} $\mathcal A$ (for $n$ letters) consists of $G_0,S,\eta,h,j\in\N$, a \emph{type} for every letter (thin,
$\mathrm{up}(b)$ or $\mathrm{down}(b)$ with $b\in\N$), a set of \emph{cuts} and a \emph{length} $\mathrm{len}(c)$ for every
cut. The \emph{window} of a type is $[-G_0,G_0]$ for thin, $[b,b+S]\cup[-\eta,\eta]$ for $\mathrm{up}(b)$ and
$[-b-S,-b]\cup[-\eta,\eta]$ for $\mathrm{down}(b)$. The datum is \emph{valid} if there is a cut, $j\ge1$, every cut $c$
satisfies $j\le\mathrm{len}(c)\le n$, $c-\mathrm{len}(c)$ is a cut and no $c-d$ with $1\le d<\mathrm{len}(c)$ is a cut,
$1\le h$, $G_0\le h$, and conditions relating the sizes of the types to $S,\eta,h,G_0$ and to the letters following the
cuts hold. Attached to $\mathcal A$ and the prime $\ell$ are $M=\lfloor2h/(\ell-1)\rfloor+1$ and the \emph{shift}
$\sigma(\mathcal A)=\lceil(2G_0+M/\ell^{j-1})\ell/(\ell-1)^2+1\rceil$. Attached to a cut $c$ of length $L$ are: the exponents
$H_s$ ($s<L$), equal to $G_0$ if the letter $c-s$ is thin or $L-s\le j$, and to $h$ otherwise; the finite set $J_c\subseteq\Z$,
the union of $[-2G_0,G_0]$ (if some letter $c-s$, $s<L$, is thin) and of the intervals $[-b-S,-b)$, resp.\ $[b,b+S)$, for
the letters $c-s$, $s<L$, of type $\mathrm{up}(b)$, resp.\ $\mathrm{down}(b)$; the number
$\kappa_c=L+\sum_{i\in J_c\setminus\{0\}}(L+1+v_\ell(i))$; and the number $\mathrm{ro}(c)\le L$ of $s<L$ for which the
letter $c-s$ is not thin.

Let a word with some non-thick letter be given. A letter $c$ is a \emph{cut} if $c+1$ is not thick. The \emph{length}
$\mathrm{len}(c)$ of a cut is the least $d\ge1$ such that $c-d+1$ is not thick; the \emph{arc} of $c$ consists of the letters
$c-\mathrm{len}(c)+1,\dots,c$ (a non-thick letter followed by $\mathrm{len}(c)-1$ thick letters), and the arcs partition the
letters. The \emph{arc datum of the word} has these cuts and lengths, $G_0=w+S$, the parameters $S,\eta$ of the word, $h$ its
height, $j=1$, and at the letter $a$ the type $\mathrm{up}(\mathrm{sz}_a)$ if $\mathrm{up}(a)$ holds and
$\mathrm{down}(\mathrm{sz}_a)$ otherwise.
\formal{Formalism.MainNotions.N.Ln17.arcOf}
\formal{Formalism.MainNotions.ArcLiftingSetup.ArcP.OK}
\formal{Formalism.MainNotions.ArcLiftingSetup.ArcP.σ}
\formal{Formalism.MainNotions.ArcLiftingSetup.ArcP.JS}
\formal{Formalism.MainNotions.N.Ln17.kapC}
\end{definition}

\paragraph{Wach coordinates.}
The arc results are statements about elements of the modules $D(T)$ of a Wach structure, expressed in coordinates
attached to the arcs. We fix the notation.

\begin{definition}[Wach coordinates at a cut]\label{pf:linkN-adm:def-wc}
Let $W$ be a Wach structure in case~A (\cref{pf:linkN-adm:wach}), with $p\ge5$, ring $A=\OO_C[[X]]^{\Z/n}$, idempotents
$e_i\in A$ of the letters, generator $\gamma_0$ of $\Gamma$ acting as $\gamma$ on every $D(T)$, and ring $E_W$ (the ring
$\mathcal E$ of $W$, on which $\gamma_0$ acts by $\gamma_E$). Let $A^\Delta\subseteq A$ be the subring of elements fixed by
the torsion part $\Delta$ of $\Gamma$, $\pi_0\in A^\Delta$ the uniformiser of $W$, $q=\pi_0+p$, $q_s=\varphi^s(q)$ for
$s\in\N$, and $\mathcal P\subseteq A$ the multiplicative monoid generated by the $q_s$. For a lattice $T$ let $N^\Delta(T)$
be the $\Delta$-invariant part of the Wach module $N(T)$ and, for a letter $c$, $N_c(T)=\{y\in N^\Delta(T):e_cy=y\}$; for
$f\in A$ put $fN_c(T)=\{fy:y\in N_c(T)\}$.
\begin{enumerate}[label=(\alph*)]
\item $I_c(T)$ is the set of $w\in D(T)$ with $e_cw=w$, $w$ fixed by $\Delta$, and $(\gamma-1)w\in\pi_0N_c(T)$.
\item For a letter $c'$ and $e\in\N$, $P_{c',e}(T)\subseteq D(T)$ is the additive group generated by the elements
  $a\varphi^e(x)$ with $a\in A^\Delta$ and $x\in N_{c'}(T)$.
\item A \emph{horizontal family} $(b,\mathrm{sol})$ at $c'$ consists of a \emph{$\Delta$-frame} $b=(b_1,\dots,b_d)$ in
  $N_{c'}(T)$ (every element of $N_{c'}(T)$ is uniquely $\sum_ma_mb_m$ with $a_m\in e_{c'}A^\Delta$) and sequences
  $\mathrm{sol}_k=(\mathrm{sol}_{k,r})_{r\ge0}$ in $L_C^d$, $\mathrm{sol}_{k,0}$ the $k$-th unit vector, such that for
  all $R,\kappa\in\N$ for which the vectors $p^\kappa\mathrm{sol}_{k,r}$ ($r<R$) are integral, the element
  $H=\sum_{r<R}\pi_0^r\sum_mp^\kappa\mathrm{sol}_{k,r,m}b_m$ satisfies $\gamma H-H\in\pi_0^RN_{c'}(T)$.
\item Let $\mathcal A$ be a valid arc datum (\cref{pf:linkN-adm:def-arcs}) and $c$ a cut of length $L$; put $c'=c-L$,
  $\Lambda_c=\prod_{s<L}q_s^{H_s}\in\mathcal P$ with the exponents $H_s\in\{G_0,h\}$ fixed by the types,
  $M=\lfloor2h/(\ell-1)\rfloor+1$ and $\sigma=\sigma(\mathcal A)$. For a horizontal family at $c'$, an element $x\in D(T)$
  \emph{has coefficients} $f=(f_k)\in E_W^d$ if there are $\kappa,\kappa'\in\N$, elements $a_k\in A^\Delta$ with
  $\Lambda_cf_k=a_k$ in $E_W$, and truncations $H_k$ as in (c) of the horizontal sections (with $R=M$ and the multiplier
  $p^\kappa$), such that $p^{\kappa'}\bigl(p^\kappa\Lambda_cx-\sum_ka_k\varphi^L(H_k)\bigr)\in\pi_0\Lambda_c^2P_{c',L}(T)$.
\item The \emph{class relation} at $c$: $w\in D(T)$ and $x\in H^1(T)$ are related, written $w\mapsto_cx$, if there is
  a \emph{transport term} $\tau$ for $w$ (an element of $\pi_0P_{c,j}(T)$ built from a solution of the $\gamma$-equation for
  $(\gamma-1)w$ in a $\Delta$-frame at $c$, as in the formal definition) such that $x$ is the image under the class map of
  $W$ of a $\Delta$-Herr cocycle $(p^\sigma\varphi^j(w)-\tau,z)$ with $z\in\pi_0^MN^\Delta(T)$.
\item Let $F\colon T_1\to T_2$ be a lattice map and $(b,\mathrm{sol})$ a horizontal family of $T_2$ at $c'$. A vector
  $g\in E_W^d$ \emph{lifts with exponent} $e$ if there are $w_1\in I_c(T_1)$ with $\Lambda'w_1\in N_c(T_1)$ for some
  $\Lambda'\in\mathcal P$, and $z\in\pi_0N_c(T_2)$, such that $\Lambda_cF(w_1)\in P_{c',L}(T_2)$ and $F(w_1)-z$ has the
  coefficients $(p^eg_k)_k$. $F$ has the \emph{arc lifting property} for $(b,\mathrm{sol},k_0,f)$ ($f\in E_W$) at $c$ with
  exponent $e$ if the vector $(\delta_{k,k_0}f)_k$ lifts with exponent $e$.
\item $F$ \emph{kills at the letter} $i$ with $a\in A$ if $ay\in F(N_i'(T_1))$ for every $\Delta$-invariant $y$ in the
  letter-$i$ part $N_i'(T_2)$ of $N(T_2)$.
\item A lattice $T$ satisfies the \emph{window condition} for $\mathcal A$ if it is crystalline and its jumps at the
  embeddings over every letter $i$ lie in the window of the type of $i$.
\end{enumerate}
The formal development also defines a \emph{filtration condition} on an index $k_0$ of a horizontal family (the
vector $\varphi^L\mathrm{cmp}(b_{k_0})$ lies in the $\varphi^s$-images of suitable steps of the Hodge filtration on
$D_{\mathrm{cris}}$ at the letters $c-s$ of type up), a condition (Id) on an element $f\in E_W$ (membership in fractional
ideals of $E_W$ determined by the types of the arc), and a \emph{pole condition} on an element $w$ relative to a
horizontal family (its coefficients are combinations of elements of such ideals with poles bounded by the types). We use
these only as hypotheses that are passed on.
\formal{Formalism.MainNotions.ArcLiftingSetup.InI}
\formal{Formalism.MainNotions.ArcLiftingSetup.Pp}
\formal{Formalism.MainNotions.ArcLiftingSetup.IsHor}
\formal{Formalism.MainNotions.ArcLiftingSetup.RepVV}
\formal{Formalism.MainNotions.ArcLiftingSetup.IsRho}
\formal{Formalism.MainNotions.N.Ln17.PCutBoundRho.LV}
\formal{Formalism.MainNotions.N.FrobeniusPowerSeries.LiftsAt}
\formal{Formalism.MainNotions.LiftingTheory.KillsΔ}
\formal{Formalism.MainNotions.ArcLiftingSetup.Win}
\end{definition}

\begin{definition}[The core local statements]\label{pf:linkN-adm:def-cores}
Let $W$ be as in \cref{pf:linkN-adm:def-wc}, $C_R\in\R$, and $\mathrm{Pt}$ a predicate on triples (letter, element of
$L_C$, embedding $K\to L_C$), the \emph{pole points}. The \emph{core local statements} at $W$ are the following. Except
for (C0), (r2), (ext) and the jump statements, each one concerns $W$ in case~A and is quantified over valid arc data
$\mathcal A$, cuts $c$ and lattices with the window condition; we describe them briefly.
\begin{itemize}
\item[(C0)] $p\ge5$, and $\OO_C\to L_C$ is injective with image $\{x:v(x)\ge0\}$.
\item[(wr)] The \emph{reduction statement} with constant $C_R$: let $F\colon T_1\to T_2$ be an onto map of crystalline
  lattices such that the jumps of $T_1$ at the embeddings over every letter $i$ lie in $[-M_i,h_i]$; put $H_k=M_k+h_k$,
  $H_{\max}=\max_kH_k$ and $c_H=\lceil1+\log_\ell(H_{\max}+1)\rceil$. Then for every letter $i$ there are
  $s_i,r_i\in\N$ with
  \[
  s_i\le\sum_{a\ge0}\bigl(H_{i+a+1}+c_H+3\bigr)(\ell-1)^{-a},\qquad r_i\le\frac{s_i\ell}{(\ell-1)^2}+c_H,
  \]
  such that $F$ kills at $i$ with $p^{r_i}$; the statement contains further bounds of the same kind.
\item[(H), (Tr)] Two statements of the local lifting theory: (H) describes $\gamma$ on $\pi_0^rN_i(T)$ through a unit
  $u_1\in\OO_C$ with $v(u_1^r-1)=1+v_p(r)$ and bounds the coefficients of solutions of the $\gamma$-equation in a
  $\Delta$-frame; (Tr) gives decompositions of the elements $\varphi^k(q)$ and of the powers of $\varphi^j(\pi_0)$.
\item[(dc)] $D_{\mathrm{cris}}$ is exact and strict on short exact sequences of crystalline representations.
\item[(relay)] The \emph{relay statement}: for every crystalline $T$ with the window condition, the elements of
  $P_{c,j}(T)$ are mapped into $\pi_0^MN_{c+j}(T)$ by $\mathrm{JJ}=\pi_0^M\prod_{k<j}q_k^{2G_0}$ and into $N_{c+j}(T)$ by
  $\prod_{s<j}q_s^{G_0}$; every $w\in I_c(T)$ has a transport term, unique modulo $\mathrm{JJ}\cdot P_{c,j}(T)$; the class
  relation $\mapsto_c$ is the graph of an additive map on $I_c(T)$ which vanishes on $I_c(T)\cap\pi_0N_c(T)$ and takes
  values in $\linkNadmHf(T)$ on the elements $w$ with $\Lambda'w\in N_c(T)$ for some $\Lambda'\in\mathcal P$; and it is
  natural: if $w\mapsto_cx$ then $F(w)\mapsto_cF_*x$ for every lattice map $F$ between lattices with the window condition.
\item[(arc)] The \emph{arc statement}: for $j\le e\le L$, with $a=c'+e$ and $\Lambda=\Lambda_{a,e}$ the analogue of
  $\Lambda_c$ for $(a,e)$: $\varphi^{e-j}(\mathrm{JJ})\in\pi_0\Lambda^2A^\Delta$; $\Lambda P_{c',e}(T)\subseteq N_a(T)$ and
  $\Lambda N_a(T)\subseteq P_{c',e}(T)$; and for every horizontal family at $c'$: the vectors $\varphi^e\mathrm{cmp}(b_k)$
  form a basis of the part of $D_{\mathrm{cris}}$ at the embedding of the letter $a$; every $x$ with $\Lambda x\in
  P_{c',e}(T)$ is such that $p^\sigma x$ has coefficients in $\Lambda^{-1}e_aA^\Delta$, and conversely every such coefficient
  vector $f$ is the vector of coefficients of $p^\sigma x$ for some $x$ with $\Lambda x\in P_{c',e}(T)$; the analogous
  statements for $N_a(T)$ with a \emph{lattice condition} (Lat) on the coefficients; and the \emph{frame property}: if
  $\Lambda x\in P_{c',e}(T)$ and $x$ has coefficients $0$, then $x\in\pi_0\Lambda P_{c',e}(T)$.
\item[(arc$'$)] If $F$ is onto between crystalline lattices and kills at a letter $c'$ with $p^r$, then for every $e$ and
  $x_2\in N_{c'}(T_2)$ there is $x_1\in N_{c'}(T_1)$ with $F(\varphi^e\mathrm{cmp}\,x_1)=p^r\varphi^e\mathrm{cmp}\,x_2$
  in $V_2\otimes B$, where $\mathrm{cmp}$ is the comparison map of $W$.
\item[(split)] The \emph{splitting statement}: if $Y$ is a crystalline extension of the unit lattice $\mathbf 1$ by $T_2$
  attached to a cocycle $\xi$ of $T_2$ and the projection $Y\to\mathbf 1$ kills at every cut with $p^r$, there are
  $w_c\in I_c(T_2)$ with $\Lambda_cw_c\in P_{c',L}(T_2)$ satisfying the pole condition, and classes $x_c$ with
  $w_c\mapsto_cx_c$, such that $p^{r+2\sigma}[\xi]=\pm\sum_cx_c$.
\item[(univ), (nn)] The \emph{universal statements}: (univ) provides, for $T_2$ with the window condition and low counts
  at most $m$ at the letters of type up or down, and for a cut $c$, a horizontal family of $T_2$ at $c'$ with the
  properties (U0)--(U3) listed in the proof of \cref{pf:linkN-adm:arclen}; (nn) is a related statement on constant
  coefficient vectors.
\item[(sat), (sat$'$)] Saturation: if $f,f'$ are products of the $q_s$, $f'x\in N^\Delta(T)$ (resp.\ $f'x\in P_{c,e}(T)$)
  and $p^kx\in fN^\Delta(T)$ (resp.\ $p^kx\in fP_{c,e}(T)$) for some $k$, then $x\in fN^\Delta(T)$ (resp.\
  $x\in fP_{c,e}(T)$).
\item[(kill1), (kill2)] The \emph{killer statements} for the polynomials
  \[
  P_{L,J}(X)=\prod_{i\in J\setminus\{0\}}\bigl(X^{p^L}-\chi(\gamma_0)^{ip^L}\bigr)\cdot\sum_{k<p^L}X^k\in\OO_C[X]
  \qquad(L\in\N,\ J\subseteq\Z\text{ finite}):
  \]
  (kill1) $v(P_{L,J}(1))=L+\sum_{i\in J\setminus\{0\}}(L+1+v_p(i))$, together with the action of $P_{L,J}(\gamma_E)$ on
  fractional ideals of $E_W$; (kill2) $P_{L,J}(\gamma)w-P_{L,J}(1)w\in\pi_0N_c(T)$ whenever $T$ is crystalline, $e_cw=w$
  and $(\gamma-1)w\in\pi_0N_c(T)$.
\item[(r2)] The \emph{content divisibility statement}, a statement on power series over discrete valuation rings.
\item[(P)] The clauses (P2), (P3), (P4) (with the pole points $\mathrm{Pt}$), (P5) and (P9) of the theory of
  $\Delta$-invariant Wach lattices; for instance (P2) states that $N(T)$ is $\Gamma$-stable with $(\gamma-1)N(T)\subseteq
  XN(T)$ for crystalline $T$, and (P9) gives the ring relations of $\pi_0$ and the structure of $N(T)$ over the letters.
\item[(pt)] For every letter $i$ and embedding $\tau$ over $i$ there is a primitive $p$-th root of unity $\zeta$ with
  $\mathrm{Pt}(i,\zeta,\tau)$; and $\gamma(\pi_0)\in\pi_0A^\Delta$.
\item[(ext)] Every finite cocycle of a crystalline lattice $T$ is attached to a crystalline extension of $\mathbf 1$ by $T$.
\item[(cj), (jq), (jy)] In a short exact sequence of crystalline lattices the jumps of the middle term are those of the
  two others; if $F\colon T_1\to T_2$ is onto between crystalline lattices, every jump of $T_2$ is a jump of $T_1$; every
  jump of a crystalline extension of $\mathbf 1$ by a crystalline $T_2$ is $0$ or a jump of $T_2$.
\item[(c4)] The letter condition at $W$ (\cref{pf:linkN-adm:def-letter}).
\item[(sec), (gl), (hf), (fl)] Statements on coefficients relative to horizontal families: compatibility with lattice
  maps (sec), lifting of coefficient vectors along $\gamma$ (gl), existence of a horizontal family at every letter for
  every lattice with the window condition (hf), and low steps of the filtration (fl).
\end{itemize}
\formal{Formalism.MainNotions.N.Ln17.CoresW}
\formal{Formalism.MainNotions.LiftingTheory.coreWR}
\formal{Formalism.MainNotions.ArcLiftingSetup.RelayCore}
\formal{Formalism.MainNotions.ArcLiftingSetup.ArcCore}
\formal{Formalism.MainNotions.ArcLiftingSetup.ArcLCore}
\formal{Formalism.MainNotions.ArcLiftingSetup.SplitCore}
\formal{Formalism.MainNotions.ArcLiftingSetup.UnivCore}
\formal{Formalism.MainNotions.ArcLiftingSetup.KillerLevCore}
\formal{Formalism.MainNotions.ArcLiftingSetup.KillerModCore}
\end{definition}

\paragraph{The killing exponent and the arc setup.}

\begin{lemma}[Killing exponent at a cut]\label{pf:linkN-adm:killcut}
Let $W$ be in case~A with $\ell\ge11$ and (wr). Let $\Phi\colon T\to T'$ be an onto map of crystalline lattices such that
the jumps of $T$ over every letter $i$ lie in $[-M_i,h_i]$, and let $h,w,S,\eta\in\N$ with $\eta+1\le h$ and
$M_i+h_i\le h+\eta$ for all $i$. Put
\[
r_B=\frac{2(w+S+\eta)}\ell+\frac{2h}{\ell^2}+2\log_\ell(2h)+3 .
\]
If $c$ is a letter with $M_{c+1}+h_{c+1}\le w+S+\eta$, then $\Phi$ kills at $c$ with $p^{\lfloor r_B\rfloor}$. Moreover
$r_B\le3\bigl(h/\ell^2+(w+S+\eta)/\ell+\log_\ell(2h)+1\bigr)$.
\formallabel{Formalism.MainNotions.N.FrobeniusPowerSeries.SetupA.kill_cut}
\formallabel{Formalism.MainNotions.N.FrobeniusPowerSeries.SetupA.rr_le}
\formallabel{Formalism.MainNotions.N.FrobeniusPowerSeries.SetupA.tsum_three}
\formallabel{Formalism.MainNotions.N.FrobeniusPowerSeries.SetupA.cW_le}
\formallabel{Formalism.MainNotions.N.FrobeniusPowerSeries.SetupA.killsΔ_mono}
\end{lemma}

\begin{proof}
Let $s=s_c$, $r=r_c$ be given by (wr), $H_1=H_{c+1}\le w+S+\eta$, $H=H_{\max}\le h+\eta$ and $c_H$ as in (wr). Put
$x_0=H_1+c_H+3$, $X=H+c_H+3$ and $u=(\ell-1)^{-1}<1$. The terms of the series in (wr) are nonnegative, the term for
$a=0$ is at most $x_0$ and the others at most $X$, so comparison with a geometric series gives
\[
s\le x_0+Xu+\frac{Xu^2}{1-u}=x_0+\frac X{\ell-1}+\frac X{(\ell-1)(\ell-2)}=x_0+\frac X{\ell-2}.
\]
Since $H+1\le h+\eta+1\le2h$ and $\lceil y\rceil<y+1$, we have $c_H\le2+\lambda$ with $\lambda=\log_\ell(2h)\ge0$. Put
$P=\ell/(\ell-1)^2$ and $Q=\ell/((\ell-1)^2(\ell-2))$. For $\ell\ge11$ one has $P+Q=\ell/((\ell-1)(\ell-2))\le\min(2/\ell,1/8)$
and $Q\le2/\ell^2$ (direct verifications). From $r\le sP+c_H$,
\[
r\le(w+S+\eta)P+(h+\eta)Q+(c_H+3)(P+Q)+c_H\le\frac{2(w+S+\eta)}\ell+\frac{2h}{\ell^2}+\frac{9c_H+3}8,
\]
using $\eta Q\le(w+S+\eta)Q$; and $(9c_H+3)/8\le(21+9\lambda)/8\le2\lambda+3$. Hence $r\le r_B$, and as $r\in\N$,
$r\le\lfloor r_B\rfloor$. If $\Phi$ kills at $c$ with $p^r$ it kills with $p^{r'}$ for every $r'\ge r$, since the image
$\Phi(N_c'(T))$ is an $\OO_C$-module. The last inequality holds termwise, as $\lambda\ge0$.
\end{proof}

\begin{proposition}[Arc setup]\label{pf:linkN-adm:arcsetup}
Let $W$ be in case~A, let $T_1,T_2$, $F$ and a word $(b,w,S,\eta,m,\mathrm{up},\mathrm{sz},M^-,M^+)$ form a setting for
theorem~B with $\ell\ge11$, $20w,20\eta,20S\le b$ and some letter not thick, and let $\mathcal A$ be the arc datum of the
word. Assume (wr) for some $C_R$, (dc), (cj), and that every jump of $T_2$ is a jump of $T_1$. Then: $\mathcal A$ is valid;
$T_1$ and $T_2$ satisfy the window condition; $\sum_c\mathrm{len}(c)=n$; $\mathrm{len}(c)\le3$ for all cuts if the type
alternates and no three consecutive letters are thick; $\mathrm{len}(c)\le R$ for all cuts if $R\ge1$, the type is
constant and no $R$ consecutive letters are thick; $\sigma(\mathcal A)\le3\,\mathrm{Sg}$; and there is $r_r\in\N$ with
$r_r\le3\,\mathrm{Sg}$ such that $F$ kills at every cut with $p^{r_r}$ and, for every crystalline extension $Y$ of the unit
lattice by $T_2$ attached to a cocycle of $T_2$, $Y$ satisfies the window condition and the projection $Y\to\mathbf 1$
kills at every cut with $p^{r_r}$.
\formallabel{Formalism.MainNotions.N.FrobeniusPowerSeries.arc_setup_theorem}
\formallabel{Formalism.MainNotions.N.FrobeniusPowerSeries.SetupA.arcSetup_of}
\end{proposition}

\begin{proof}
The validity of $\mathcal A$, the sum of the lengths, the two length bounds and the bound for $\sigma(\mathcal A)$ are
\lref{Formalism.MainNotions.N.FrobeniusPowerSeries.SetupA.arcBasic_core}, a combinatorial analysis of the arcs of a
word together with a numerical estimate of $\sigma$ (it uses only the setting). The hypothesis (dc) is not used.

\emph{Windows.} The type of the letter $i$ is $\mathrm{up}(\mathrm{sz}_i)$ or $\mathrm{down}(\mathrm{sz}_i)$, with window
$[\mathrm{sz}_i,\mathrm{sz}_i+S]\cup[-\eta,\eta]$, resp.\ $[-\mathrm{sz}_i-S,-\mathrm{sz}_i]\cup[-\eta,\eta]$. The setting
says that the jumps of $T_1$ lie in these windows and that $T_1,T_2$ are crystalline; the jumps of $T_2$ are jumps of
$T_1$. So $T_1$ and $T_2$ satisfy the window condition.

\emph{Killing.} Let $h$ be the height of the word. Put $M'_i=\eta$, $h'_i=\mathrm{sz}_i+S$ if $\mathrm{up}(i)$, and
$M'_i=\mathrm{sz}_i+S$, $h'_i=\eta$ otherwise. As $\eta<\mathrm{sz}_i$, every window is contained in $[-M'_i,h'_i]$, so the
jumps of every lattice with the window condition over the letter $i$ lie in $[-M'_i,h'_i]$. Next
$M'_i+h'_i=\mathrm{sz}_i+S+\eta\le h+\eta$: for a thick letter $\mathrm{sz}_i\le b+\eta$ and $h=b+\eta+S$, and for a
non-thick letter $\mathrm{sz}_i\le w$ and $h\ge w+S$. A non-thick letter $i$ gives $\eta+1\le\mathrm{sz}_i\le w\le h$. If
$c$ is a cut, $c+1$ is not thick, so $M'_{c+1}+h'_{c+1}\le w+S+\eta$. Let $r_r=\lfloor r_B\rfloor$ with $r_B$ as in
\cref{pf:linkN-adm:killcut} for $(h,w,S,\eta)$. Since $r_B\ge0$ and the right-hand side of the last inequality of that
lemma is $3\,\mathrm{Sg}$, we get $r_r\le3\,\mathrm{Sg}$, and the lemma applied to $F$ shows that $F$ kills at every cut
with $p^{r_r}$.

\emph{Extensions.} Let $Y$ be a crystalline extension of $\mathbf 1$ by $T_2$, with projection $\pi\colon Y\to\mathbf 1$.
The unit lattice is crystalline with the single jump $0$ (\lref{Formalism.MainNotions.LiftingTheory.lone_HOne0}, which
states this for the unit lattice of $W$). By (cj) every jump of $Y$ is a jump of $T_2$ or $0$, and $0$ lies in every window;
so $Y$ satisfies the window condition. Since $\pi$ is onto between crystalline lattices and the jumps of $Y$ lie in
$[-M'_i,h'_i]$, \cref{pf:linkN-adm:killcut} applies to $\pi$ with the same data and shows that $\pi$ kills at every cut
with $p^{r_r}$.
\end{proof}

\paragraph{Arc lifting.}
In the next two propositions $W$ is in case~A, $T_1,T_2$, $F$ and a word form a setting for theorem~B with $\ell\ge11$,
$20w,20\eta,20S\le b$ and some letter not thick, $\mathcal A$ is the arc datum of the word, $T_1$ and $T_2$ satisfy the
window condition, and $r_r\in\N$ satisfies $r_r\le3\,\mathrm{Sg}$ and $F$ kills at every cut with $p^{r_r}$. We assume the
core local statements at $W$ (only some of them are used in each proof). For a cut $c$ we consider a horizontal family
$(b,\mathrm{sol})$ of $T_2$ at $c'=c-\mathrm{len}(c)$, an index $k_0$ satisfying the filtration condition, and $f\in E_W$
satisfying (Id).

\begin{proposition}[Arcs of a single type]\label{pf:linkN-adm:single}
Let $c$ be a cut such that $\mathrm{up}(c-s)=\mathrm{up}(c)$ for all $s<\mathrm{len}(c)$. Then $F$ has the arc lifting
property for $(b,\mathrm{sol},k_0,f)$ at $c$ with some exponent $e_1\le6\,\mathrm{Sg}+1$.
\formallabel{Formalism.MainNotions.N.FrobeniusPowerSeries.arc_single_type_bound}
\end{proposition}

\begin{proof}
Take $e_1=r_r+\sigma(\mathcal A)\le3\,\mathrm{Sg}+3\,\mathrm{Sg}$ (\cref{pf:linkN-adm:arcsetup}). Every letter of the arc
datum of a word has type up or down. Let $L=\mathrm{len}(c)$.

If $\mathrm{up}(c)$ holds, all letters of the arc have type up. The lifting lemma for arcs of type up
(\lref{Formalism.MainNotions.N.FrobeniusPowerSeries.liftsAt_thick_letter_type}: under the core local statements, if
$F$ is onto and kills at every cut with $p^{r}$, then for a horizontal family at $c'$, an index with the filtration
condition and $f$ satisfying the ideal conditions of (Id), the arc lifting property holds with exponent
$r+\sigma(\mathcal A)$) gives the claim.

Otherwise all letters of the arc have type down, and the denominator of the ideal condition in (Id) is $1$. The letter
$c'$ is a cut, since $\mathcal A$ is valid. By (hf) there is a horizontal family $(b^1,\mathrm{sol}^1)$ of $T_1$ at $c'$. By
(arc$'$), applied to the killing at the cut $c'$ and to $b_{k_0}\in N_{c'}(T_2)$, there is $x_1\in N_{c'}(T_1)$ with
$F(\varphi^L\mathrm{cmp}\,x_1)=p^{r_r}\varphi^L\mathrm{cmp}\,b_{k_0}$. Write $x_1=\sum_ma_mb^1_m$ with $a_m\in e_{c'}A^\Delta$;
the comparison map is additive and satisfies $\mathrm{cmp}(ay)=a(0)\,\mathrm{cmp}(y)$ for $y$ at the letter $c'$, where
$a(0)\in\OO_C$ is the constant coefficient of the component of $a$ at $c'$. With $c_1=(a_m(0))_m\in\OO_C^{d_1}$ we obtain
$F\bigl(\sum_mc_{1,m}\varphi^L\mathrm{cmp}\,b^1_m\bigr)=p^{r_r}\varphi^L\mathrm{cmp}\,b_{k_0}$. The lifting lemma for arcs
with a constant vector (\lref{Formalism.MainNotions.N.FrobeniusPowerSeries.liftsAt_thin_letter_type}: under the core
local statements, given horizontal families of $T_1$ and $T_2$ at $c'$, a set $S_i$ of up letters of the arc, a vector
$c_1$ satisfying the filtration conditions at $S_i$ with $F(\sum_mc_{1,m}\varphi^L\mathrm{cmp}\,b^1_m)=
p^{s_0}\varphi^L\mathrm{cmp}\,b_{k_0}$, and $f$ in the fractional ideals determined by $S_i$ and the down letters, the arc
lifting property holds with exponent $s_0+\sigma(\mathcal A)$), applied with $S_i=\emptyset$ and $s_0=r_r$, gives the
claim; its conditions on $f$ are those of (Id), and the condition at up letters is vacuous.
\end{proof}

\begin{proposition}[Arc lifting for alternating types]\label{pf:linkN-adm:arcalt}
Let $C_b\ge2$ and let $\beta$ be the band indicator of the word with constant $C_b$. Assume that the type alternates and no
three consecutive letters are thick, and let $c$ be a cut with $\mathrm{len}(c)\le2$ or $40\,\mathrm{Sg}\le b/4$. Then $F$
has the arc lifting property for $(b,\mathrm{sol},k_0,f)$ at $c$ with some exponent $e_1\le40\,\mathrm{Sg}+\eta\beta$.
\formallabel{Formalism.MainNotions.N.FrobeniusPowerSeries.arc_proposition_alternative}
\end{proposition}

\begin{proof}
We have $\mathrm{Sg}\ge1$ (all terms of \eqref{pf:linkN-adm:eq-Sg} are nonnegative, with $\log_\ell0=0$) and
$\eta\beta\ge0$. By \cref{pf:linkN-adm:arcsetup}, $1\le\mathrm{len}(c)\le3$.
\begin{itemize}
\item If $\mathrm{len}(c)=1$, the arc consists of $c$ alone and \cref{pf:linkN-adm:single} gives $e_1\le6\,\mathrm{Sg}+1$.
\item If $\mathrm{len}(c)=2$, then $\mathrm{up}(c-1)\ne\mathrm{up}(c)$ because the type alternates, and the arc theorem for
  length two (\lref{Formalism.MainNotions.N.FrobeniusPowerSeries.arc_length_one_two_theorem}: under these hypotheses,
  for an arc of length $2$ with letters of different types, the arc lifting property holds with an exponent at most
  $15\,\mathrm{Sg}+1+\eta\beta$) applies.
\item If $\mathrm{len}(c)=3$, then $40\,\mathrm{Sg}\le b/4$. If $\mathrm{up}(c-2)$ fails, the arc theorem for length three
  (\lref{Formalism.MainNotions.N.FrobeniusPowerSeries.arc_length_three_theorem}: for an arc of length $3$ whose first
  letter has type down, with $40\,\mathrm{Sg}\le b/4$, the exponent can be taken at most $22\,\mathrm{Sg}+1+\eta\beta$)
  applies; if it holds, the companion result
  (\lref{Formalism.MainNotions.N.FrobeniusPowerSeries.exists_liftsAt_of_arc_len_three_of_up}: the same for a first
  letter of type up, with exponent at most $34\,\mathrm{Sg}+1+\eta\beta$) applies.
\end{itemize}
In all cases $e_1\le34\,\mathrm{Sg}+1+\eta\beta\le40\,\mathrm{Sg}+\eta\beta$.
\end{proof}

\paragraph{The arc length bound.}
For a cut $c$ of a valid arc datum let $J_c\subseteq\Z$ be the finite set attached to $c$ by the types
(\cref{pf:linkN-adm:def-arcs}), so that $\kappa_c=L+\sum_{i\in J_c\setminus\{0\}}(L+1+v_\ell(i))$ with $L=\mathrm{len}(c)$, and
let $P_c=P_{L,J_c}$ be the killer polynomial of (kill1). Thus (kill1) says $v(P_c(1))=\kappa_c$.

\begin{lemma}[Transfer of lifts]\label{pf:linkN-adm:transfer}
Let $W$ be in case~A with $p\ge5$, (relay) and (arc); let $\mathcal A$ be valid, $c$ a cut of length $L$, $F\colon T_1\to T_2$
a lattice map between lattices with the window condition, and $(b,\mathrm{sol})$ a horizontal family of $T_2$ at $c'$. If
$g\in E_W^d$ lifts with exponent $e\ge\sigma(\mathcal A)$, and $u\in I_c(T_2)$ with $\Lambda_cu\in P_{c',L}(T_2)$ has the
coefficients $(p^\sigma g_k)_k$ and $u\mapsto_cx_2$, then there is $x_1\in\linkNadmHf(T_1)$ with
$F_*x_1=p^{e-\sigma}x_2$.
\formallabel{Formalism.MainNotions.N.Ln17.PCutBoundGal.LV_galois}
\end{lemma}

\begin{proof}
Let $w_1$, $\Lambda'$ and $z$ witness the lift and put $X=F(w_1)-z-p^{e-\sigma}u$. Its coefficients are
$p^eg_k-p^{e-\sigma}p^\sigma g_k=0$ (coefficients are additive and compatible with multiplication by integers), and
$\Lambda_cX\in P_{c',L}(T_2)$, since $\Lambda_cF(w_1)$ and $\Lambda_cu$ lie there, and $z\in\pi_0N_c(T_2)\subseteq N_c(T_2)$
gives $\Lambda_cz\in P_{c',L}(T_2)$ by (arc). By the frame property, $X\in\pi_0\Lambda_cP_{c',L}(T_2)\subseteq\pi_0N_c(T_2)$, using
$\Lambda_cP_{c',L}\subseteq N_c$ from (arc). Hence $F(w_1)-p^{e-\sigma}u=z+X\in\pi_0N_c(T_2)$; it lies in $I_c(T_2)$
(as $F(w_1)$ and $u$ do), so its class is $0$ by (relay). By (relay) $w_1\mapsto_cx_1$ for some $x_1$, which is finite
since $\Lambda'w_1\in N_c(T_1)$; by naturality $F(w_1)\mapsto_cF_*x_1$. By additivity $p^{e-\sigma}u\mapsto_c p^{e-\sigma}x_2$,
and the uniqueness of classes gives $F_*x_1=p^{e-\sigma}x_2$.
\end{proof}

\begin{proposition}[Arc length bound]\label{pf:linkN-adm:arclen}
Let $W$ be in case~A with (C0), (relay), (arc), (nn), (univ), (kill1), (kill2), (P2), (P9) and $\gamma(\pi_0)\in\pi_0A^\Delta$.
Let $\mathcal A$ be a valid arc datum, $F\colon T_1\to T_2$ an onto map of crystalline lattices with the window condition,
$m\in\N$ such that the low counts of $T_2$ at the letters of type up or down are at most $m$, $c$ a cut of length $L$, and
$e_1\ge\sigma(\mathcal A)$. Assume that $F$ has the arc lifting property at $c$ with exponent $e_1$ for every horizontal
family $(b,\mathrm{sol})$ of $T_2$ at $c'$, every index $k_0$ satisfying the filtration condition and every $f$ satisfying
(Id). Then there is a set $Z_c\subseteq\linkNadmHf(T_2)$ of at most $m\,\mathrm{ro}(c)(L+1)$ elements such that for every
$w\in I_c(T_2)$ with $\Lambda_cw\in P_{c',L}(T_2)$ which satisfies the pole condition for every horizontal family at $c'$,
and every $x$ with $w\mapsto_cx$,
\[
p^{\sigma+\kappa_c+e_1}x\in F_*\linkNadmHf(T_1)+\OO_CZ_c .
\]
\formallabel{Formalism.MainNotions.N.Ln17.arc_length_bound_theorem}
\end{proposition}

\begin{proof}
Write $\Lambda=\Lambda_c$, $\sigma=\sigma(\mathcal A)$, $\kappa=\kappa_c$, $\Pi=P_{c',L}(T_2)$, $N_c=N_c(T_2)$ and
$\rho\colon I_c(T_2)\to H^1(T_2)$ for the map with graph $\mapsto_c$ given by (relay). It is additive, vanishes on
$I_c(T_2)\cap\pi_0N_c$, and is $\OO_C$-linear: apply the naturality in (relay) to the lattice map $T_2\to T_2$ given by
multiplication by $o\in\OO_C$, which induces multiplication by $o$ on $D(T_2)$ and on $H^1(T_2)$. Recall that $\gamma$
commutes with $e_c$ and with $\Delta$, and maps $N^\Delta(T_2)$ into itself by (P2). Hence $\pi_0N_c\subseteq I_c(T_2)$: for
$z=\pi_0y$ with $y\in N_c$, $z$ is fixed by $e_c$ and $\Delta$ (as $\pi_0\in A^\Delta$), and writing
$\gamma(\pi_0)=\pi_0u$ with $u\in A^\Delta$ we get $(\gamma-1)z=\pi_0(u\gamma y-y)\in\pi_0N_c$. Similarly $I_c(T_2)$ is stable
under $\gamma$ (if $(\gamma-1)w=\pi_0y$ then $(\gamma-1)\gamma w=\pi_0u\gamma y$), hence under all polynomials in $\gamma$
with coefficients in $\OO_C$.
By (arc) with $e=L$ we have $\Lambda\Pi\subseteq N_c$, $\Lambda N_c\subseteq\Pi$, and the frame property. Combining them:
\begin{itemize}
\item[(F)] \emph{if $X\in D(T_2)$ satisfies $\Lambda X\in\Pi$ and has coefficients $0$, then $X\in\pi_0N_c$}, and so
  $\rho(X)=0$ when $X\in I_c(T_2)$.
\end{itemize}
The hypothesis (nn) is not used. \emph{The universal data.} By (univ) for $T_2$, $m$ and $c$ there are a horizontal family $(b,\mathrm{sol})$ of $T_2$ at
$c'$, a set $R$ of at most $m\,\mathrm{ro}(c)$ indices, and vectors $J_1,\dots,J_g\in E_W^d$ with $g\le m\,\mathrm{ro}(c)(L+1)$,
such that:
\begin{itemize}
\item[(U0)] every index $k\notin R$ satisfies the filtration condition;
\item[(U1)] for every $w\in I_c(T_2)$ with $\Lambda w\in\Pi$ satisfying the pole condition there is $f\in E_W^d$ with
  entries in $\Lambda^{-1}e_cA^\Delta$ such that $p^\sigma w$ has the coefficients $f$, the elements $P_c(\gamma_E)f_k$
  ($k\notin R$) satisfy (Id), and the vectors $f''=(P_c(\gamma_E)f_k)_{k\notin R}$ and $f^R=(P_c(\gamma_E)f_k)_{k\in R}$
  (extended by $0$) satisfy the \emph{invariance condition} ($(\gamma_E-1)g=\pi_0g'$ with $g'$ satisfying the lattice
  condition of (arc));
\item[(U2)] the $J_i$ have entries in $\Lambda^{-1}e_cA^\Delta$, vanish outside $R$ and satisfy the invariance condition;
  and every vector with these three properties is $\sum_io_iJ_i+\pi_0g'$ with $o_i\in\OO_C$ and $g'$ satisfying the
  lattice condition;
\item[(U3)] every vector $g$ with entries in $\Lambda^{-1}e_cA^\Delta$ satisfying the invariance condition is realised:
  there is $u\in I_c(T_2)$ with $\Lambda u\in\Pi$ which has the coefficients $(p^\sigma g_k)_k$.
\end{itemize}
Moreover, by (arc), every $g'$ satisfying the lattice condition is realised by some $x_0\in N_c$ with coefficients
$(p^\sigma g'_k)_k$.

\emph{The set $Z_c$.} For each $i$ choose $u_i$ for $J_i$ by (U3) and put $x_i=\rho(u_i)$. Since $\Lambda u_i\in\Pi$, we have
$\Lambda^2u_i\in N_c$ with $\Lambda^2\in\mathcal P$, so $x_i\in\linkNadmHf(T_2)$ by (relay). Let $Z_c=\{x_1,\dots,x_g\}$.

\emph{The bound.} Let $w$, $x$ be as in the statement, and $f$ as in (U1). For $k\notin R$ the hypothesis, applied to the
index $k$ (which satisfies the filtration condition by (U0)) and to $P_c(\gamma_E)f_k$, shows that the vector with the
single entry $P_c(\gamma_E)f_k$ at $k$ lifts with exponent $e_1$; adding the witnesses, $f''$ lifts with exponent $e_1$. The
entries of $f''$ and $f^R$ lie in $\Lambda^{-1}e_cA^\Delta$
(\lref{Formalism.MainNotions.N.Ln17.PCutBoundGal.InFr_aeval}: this set is stable under polynomials in $\gamma_E$). By
(U3) choose $u''$ and $u^R$ for $f''$ and $f^R$, and put $x''=\rho(u'')$. By \cref{pf:linkN-adm:transfer} there is
$x_1\in\linkNadmHf(T_1)$ with $F_*x_1=p^{e_1-\sigma}x''$.

By (U2), $f^R=\sum_io_iJ_i+\pi_0g'$ with $g'$ satisfying the lattice condition; let $x_0\in N_c$ realise $g'$. The element
$X=u^R-\sum_io_iu_i-\pi_0x_0$ has coefficients $p^\sigma f^R-\sum_io_ip^\sigma J_i-p^\sigma\pi_0g'=0$, and $\Lambda X\in\Pi$
(note $\Lambda\pi_0x_0=\pi_0\Lambda x_0\in\Pi$). By (F), $\rho(X)=0$; as $\pi_0x_0\in\pi_0N_c$, also $\rho(\pi_0x_0)=0$, so
$\rho(u^R)=\sum_io_ix_i$.

Let $w'=p^\sigma w$, so that $\rho(w')=p^\sigma x$. The element $P_c(\gamma)w'$ lies in $I_c(T_2)$, satisfies
$\Lambda P_c(\gamma)w'\in\Pi$ (\lref{Formalism.MainNotions.N.Ln17.PCutBoundGal.InFPΛ_aeval}: the condition $\Lambda y\in\Pi$
is stable under polynomials in $\gamma$), and has the coefficients $(P_c(\gamma_E)f_k)_k$
(\lref{Formalism.MainNotions.N.Ln17.PCutBoundRep.RepVV_aeval}: if $y$ has coefficients $f$ then $Q(\gamma)y$ has
coefficients $(Q(\gamma_E)f_k)_k$ for every polynomial $Q$; its proof uses the first clause of (arc), the identity
$(\varphi^L\pi_0)^M\in\pi_0\Lambda^2A^\Delta$ of \lref{Formalism.MainNotions.N.Ln17.PCutBoundRho.φpow_π₀_JJ}, and the
$\gamma$-compatibility of $\Lambda$ of \lref{Formalism.MainNotions.N.Ln17.PCutBoundRho.γ_Lam}). By (kill2),
$P_c(\gamma)w'-P_c(1)w'\in\pi_0N_c$, so $\rho(P_c(\gamma)w')=P_c(1)p^\sigma x$ and $\rho(p^\sigma P_c(\gamma)w')=
p^\sigma P_c(1)p^\sigma x$. The element $X_2=p^\sigma P_c(\gamma)w'-u''-u^R$ has coefficients
$p^\sigma P_c(\gamma_E)f_k-p^\sigma f''_k-p^\sigma f^R_k=0$ and $\Lambda X_2\in\Pi$, so $\rho(X_2)=0$ by (F). Therefore
\[
p^\sigma P_c(1)\,p^\sigma x=\rho(u'')+\rho(u^R)=x''+\sum_io_ix_i .
\]
By (kill1), $v(P_c(1))=\kappa$; since $v(p)=1$, the element $p^\kappa/P_c(1)$ of $L_C$ has valuation $0$ and so, by (C0),
equals $y$ for some $y\in\OO_C$, so that $p^\kappa=y\,P_c(1)$. As
$\sigma+\kappa+e_1=\kappa+(e_1-\sigma)+2\sigma$,
\[
p^{\sigma+\kappa+e_1}x=y\,p^{e_1-\sigma}\bigl(p^\sigma P_c(1)p^\sigma x\bigr)=y\,F_*x_1+\sum_iy\,p^{e_1-\sigma}o_i\,x_i
\in F_*\linkNadmHf(T_1)+\OO_CZ_c .\qedhere
\]
\end{proof}

\begin{remark}[Two further arc inputs]\label{pf:linkN-adm:arcs}
The proof of \cref{pf:linkN-adm:sumalt} also uses two results proved in the formal development and not reproduced here.
We say that $x\in H^1(T_2)$ is \emph{concentrated at the cut} $c$ if $w\mapsto_cx$ for some $w\in I_c(T_2)$ with
$\Lambda_cw\in P_{c',L}(T_2)$ satisfying the pole condition for every horizontal family at $c'$.
\begin{itemize}
\item[(A4)] \emph{Splitting along cuts} (\lref{Formalism.MainNotions.N.Ln17.split_sum}): under the core local
  statements and the conclusions of \cref{pf:linkN-adm:arcsetup}, if for every cut $c$ the classes concentrated at $c$
  satisfy $\ell^Ex\in F_*\linkNadmHf(T_1)+\OO_C\bigcup_cZ_c$, then every $x\in\linkNadmHf(T_2)$ satisfies
  $\ell^{r_r+2\sigma(\mathcal A)+E}x\in F_*\linkNadmHf(T_1)+\OO_C\bigcup_cZ_c$. Its proof attaches to a finite class the
  crystalline extension of $\mathbf 1$ by $T_2$ that it defines, (ext), and splits it by (split).
\item[(A5)] \emph{Bounds for $\kappa$} (\lref{Formalism.MainNotions.N.Ln17.kappa_le_of_alternating},
  \lref{Formalism.MainNotions.N.Ln17.kappa_le_of_single_thick}): $\kappa_c\le13S+3\,\mathrm{lg}+3$ if the type alternates
  and $\mathrm{len}(c)\le3$; $\kappa_c\le3R^2(S+1+\mathrm{lg})$ if the type is constant and $\mathrm{len}(c)\le R$.
\end{itemize}
\end{remark}

\paragraph{The core statements at a place.}
Let the \emph{bounds package} $\mathcal B$ be the conjunction of the following seventeen statements, all contained in
$\mathcal A_B$ (\cref{pf:linkN-adm:def-AB}): the clauses (P2), (P3), (P4), (P5), (P9) of the Wach lattice theory (part of
(B1)), the local lifting theory (B2), (B3), (B4), the statements (kill1) and (kill2) of (B5), and (B6)--(B12). Recall that
the jump inputs $\mathcal J$ are Fontaine's inequality at every place, (cj) at every place, and the global letter condition.

\begin{proposition}[The core statements at a place]\label{pf:linkN-adm:cores}
Assume $\mathcal B$ (with constants $C_R,C_A,C_B,C_K$) and the jump inputs $\mathcal J$. Let $\ell\ge11$ be a prime and $(K,L_C)$ a place
in case~A (so $K$ is unramified), and let $W$ be its Wach structure of level $1$. Then the core local statements hold
at $W$ with the constant $C_R$ and the pole points of the period data.
\formallabel{Formalism.MainNotions.N.Ln17.coresAt}
\end{proposition}

\begin{proof}
We go through the list of \cref{pf:linkN-adm:def-cores}. (C0): $p=\ell\ge11$, and the valuation of the datum is
nonnegative exactly on $\OO_C$. (wr), (H) and (Tr) are components of the local lifting theory (B2) at $(\ell,K,1,L_C)$;
(dc) is (B12); (relay) is (B8), (arc) and (arc$'$) are (B9), (split) is (B10), (univ) and (nn) are (B11), at the place
$(\ell,K,L_C)$, which is in case~A with $\ell\ge5$; (sat) is (B3), (sat$'$) is (B4), (kill1) and (kill2) are parts of (B5),
(r2) is (B6), and (P) consists of the clauses of the Wach lattice theory at $W$. For (pt): the existence of the roots of
unity is \lref{Formalism.MainNotions.N.Ln17.ptEx_at} (in case~A every pole point condition is satisfied by some
primitive $p$-th root of unity), and $\gamma(\pi_0)\in\pi_0A^\Delta$ is
\lref{Formalism.MainNotions.UnusedSlot80.gamma_uniformizer_divisible_statement} (which derives it from (P9) in case~A with
$\ell\ge5$). (ext) is \lref{Formalism.MainNotions.LiftingTheory.lext_HExt_D} (a finite cocycle of a crystalline lattice
is attached to a crystalline extension of the unit lattice). (cj) is part of $\mathcal J$ and (c4) is the letter condition
at $W$. (sec), (gl), (hf), (fl) are \lref{Formalism.MainNotions.N.Ln17.arc_structures_hold_at}, which proves them from
$\mathcal B$ and $\mathcal J$. Finally (jq) and (jy) are \lref{Formalism.MainNotions.N.Ln17.jumps_of_surjective_lattice_map},
which derives them from the jump statements of $\mathcal J$.
\end{proof}

\begin{lemma}\label{pf:linkN-adm:cocycles}
Let $e\in\N$ and $Z_0\subseteq\linkNadmHf(T_2)$ finite with $\ell^ex\in F_*\linkNadmHf(T_1)+\OO_CZ_0$ for all
$x\in\linkNadmHf(T_2)$. Then there is a set $Z$ of at most $|Z_0|$ finite cocycles with
$\ell^e\linkNadmZf(T_2)\subseteq\linkNadmIm(F)+\OO_CZ$.
\formallabel{Formalism.MainNotions.N.Ln17.cocycle_classes_exponent}
\end{lemma}

\begin{proof}
Choose for each $z\in Z_0$ a finite cocycle $\tilde z$ representing it and let $Z=\{\tilde z\}$. For $c\in\linkNadmZf(T_2)$
with class $x$, write $\ell^ex=F_*y+\sum a_zz$ with $y\in\linkNadmHf(T_1)$ represented by a finite cocycle $\tilde y$. Then
$\ell^ec-F_*\tilde y-\sum a_z\tilde z$ is a coboundary, so $\ell^ec\in F_*\linkNadmZf(T_1)+B^1(T_2)+\OO_CZ$.
\end{proof}

\begin{proposition}[Exponents from arcs]\label{pf:linkN-adm:sumalt}
In a setting for theorem~B with $\ell\ge11$, $20w,20\eta,20S\le b$ and some letter not thick, under the core local
statements at $W$ (\cref{pf:linkN-adm:def-cores}), with $\beta$ the band indicator for a constant $C\ge2$:
\begin{enumerate}[label=(\alph*)]
\item if the type alternates with no three consecutive thick letters and either no letter is thick or
  $40\,\mathrm{Sg}\le b/4$, there are $e\le55\,\mathrm{Sg}+13S+\eta\beta$ and a set $Z$ of at most $4nm$ finite cocycles with
  $\ell^e\linkNadmZf(T_2)\subseteq\linkNadmIm(F)+\OO_CZ$;
\item if $R\ge1$ and the type is constant with no $R$ consecutive thick letters, the same holds with
  $e\le18\,\mathrm{Sg}+3R^2(S+1+\mathrm{lg})+1$ and $|Z|\le(R+1)nm$.
\end{enumerate}
\formallabel{Formalism.MainNotions.N.Ln17.sum_alt}
\formallabel{Formalism.MainNotions.N.Ln17.sumAltW}
\formallabel{Formalism.MainNotions.N.Ln17.sum_one}
\formallabel{Formalism.MainNotions.N.Ln17.sumOneW}
\end{proposition}

\begin{proof}
(a) Let $E_1=\lfloor40\,\mathrm{Sg}+\eta\beta\rfloor$; then $\sigma(\mathcal A)\le3\,\mathrm{Sg}\le E_1$ (note $\mathrm{Sg}\ge1$). If no
letter is thick, every arc has length $1$; so in all cases the hypothesis of \cref{pf:linkN-adm:arcalt} holds at every
cut, with the data $r_r$ of \cref{pf:linkN-adm:arcsetup}, and the arc lifting property holds with exponent $E_1$ for all
horizontal families of $T_2$ at $c'$, indices with the filtration condition and $f$ with (Id) (lifting with a smaller
exponent implies lifting with a larger one). By \cref{pf:linkN-adm:arclen} (the low counts of $T_2$ are at most $m$ at every
letter) there are sets $Z_c$; put $\kappa=\max_c\kappa_c$ and $E=\sigma(\mathcal A)+\kappa+E_1$, so that the hypothesis of (A4) holds with this $E$
(multiply by $\ell^{\kappa-\kappa_c}$). By (A4) and \cref{pf:linkN-adm:cocycles} we get $Z$ with $|Z|\le\sum_c|Z_c|$ and the
exponent $e=r_r+2\sigma(\mathcal A)+E=r_r+3\sigma(\mathcal A)+\kappa+E_1$. By \cref{pf:linkN-adm:arcsetup} and (A5), and $3\,\mathrm{lg}+3\le3\,\mathrm{Sg}$,
\[
e\le3\,\mathrm{Sg}+9\,\mathrm{Sg}+(13S+3\,\mathrm{lg}+3)+40\,\mathrm{Sg}+\eta\beta\le55\,\mathrm{Sg}+13S+\eta\beta .
\]
Since $\mathrm{len}(c)\le3$ and $\mathrm{ro}(c)\le\mathrm{len}(c)$, $\mathrm{ro}(c)(\mathrm{len}(c)+1)\le4\,\mathrm{len}(c)$, so
$|Z|\le\sum_c4m\,\mathrm{len}(c)=4nm$.

(b) The same argument, with \cref{pf:linkN-adm:single} in place of \cref{pf:linkN-adm:arcalt} and $E_1=\lfloor6\,\mathrm{Sg}+1\rfloor$, the bound $\kappa_c\le3R^2(S+1+\mathrm{lg})$ and
$\mathrm{len}(c)\le R$ gives $e\le3\,\mathrm{Sg}+9\,\mathrm{Sg}+3R^2(S+1+\mathrm{lg})+6\,\mathrm{Sg}+1$ and
$\mathrm{ro}(c)(\mathrm{len}(c)+1)\le(R+1)\mathrm{len}(c)$, so $|Z|\le(R+1)nm$.
\end{proof}

\begin{lemma}[Valuations in a window]\label{pf:linkN-adm:window}
Let $\ell$ be a prime, $a\in\Z$ and $S,N\in\N$, and assume that every integer $i$ with $a\le i<a+S$ satisfies $i\ne0$ and
$|i|\le N$. Then $\sum_{a\le i<a+S}v_\ell(i)\le N/(\ell-1)$.
\formallabel{Formalism.MainNotions.N.FrobeniusPowerSeries.sum_padicValInt_window_le}
\formallabel{Formalism.MainNotions.N.FrobeniusPowerSeries.exists_sum_padicValInt_window_eq}
\formallabel{Formalism.MainNotions.N.FrobeniusPowerSeries.padicValNat_factorial_add_sum}
\end{lemma}

\begin{proof}
We may assume $S\ge1$. The window does not contain $0$, so it consists either of negative or of positive integers; in
both cases the absolute values of its elements are $m+1,\dots,m+S$ for some $m\in\N$ with $m+S\le N$ (take $m=-a-S$,
resp.\ $m=a-1$, and use the bound for the first, resp.\ last, element). Hence
\[
\sum_{a\le i<a+S}v_\ell(i)=\sum_{j=1}^Sv_\ell(m+j)=v_\ell((m+S)!)-v_\ell(m!)\le v_\ell((m+S)!)\le\frac{m+S}{\ell-1}\le\frac N{\ell-1},
\]
by Legendre's formula $(\ell-1)v_\ell(k!)=k-s_\ell(k)\le k$, where $s_\ell(k)$ is the sum of the $\ell$-adic digits of $k$.
\end{proof}

\begin{lemma}[The lifting exponent for short words]\label{pf:linkN-adm:ePA}
Let $\ell\ge11$ be a prime and $h,\eta,H,M\in\N$ with $h\ge1$, $20\eta\le h$, $H\le h+\eta$ and $M\le h$, and let
$J\subseteq\Z$ be a finite set contained in $[-M,h]$. Then
\[
e_{PA}(\ell,M,H,J)\le\frac{7.1\,h}\ell+7.1+2\log_\ell(2h).
\]
\formallabel{Formalism.MainNotions.N.FrobeniusPowerSeries.lift_le_ePA}
\formallabel{Formalism.MainNotions.N.FrobeniusPowerSeries.lift_le}
\end{lemma}

\begin{proof}
Recall $e_{PA}=2e_W+\sum_{j\in J'}v_\ell(j)+a+v_\ell(a!)$ with $e_W=\lfloor H\ell/(\ell-1)^2+3+\log_\ell(H+1)\rfloor$,
$a=\lfloor M/(\ell-1)\rfloor+1$ and $J'$ as in \eqref{pf:linkN-adm:eq-ePA}. For $\ell\ge11$,
\[
\frac1{\ell-1}\le\frac1{10},\qquad\frac1{\ell-1}\le\frac{1.1}\ell,\qquad\frac\ell{(\ell-1)^2}\le\frac{1.21}\ell,\qquad
\frac1{(\ell-1)^2}\le\frac{0.11}\ell .
\]
Since $20\eta\le h$ and $h\ge1$ we have $\eta+1\le h$, so $H+1\le2h$ and $\log_\ell(H+1)\le\lambda:=\log_\ell(2h)$. Hence
$e_W\le1.21(h+\eta)/\ell+3+\lambda$. The elements of $J'$ are of the form $j+k$ with $-M\le j\le h$, $1\le k\le M$,
$j+k\ne0$, or lie in $[1,M]$; so $J'\subseteq[-M,-1]\cup[1,h+M]$. By \cref{pf:linkN-adm:window}, applied to the two windows
$[-M,-1]$ (with $N=M$) and $[1,h+M]$ (with $N=h+M$), $\sum_{j\in J'}v_\ell(j)\le M/(\ell-1)+(h+M)/(\ell-1)$. Next
$a\le M/(\ell-1)+1$, and Legendre's formula gives $v_\ell(a!)\le a/(\ell-1)\le M/(\ell-1)^2+1/(\ell-1)$. Adding up and using
$\eta\le h/20$, $M\le h$:
\[
e_{PA}\le\frac{2.42(h+\eta)+3\cdot1.1\,h+1.1\,h+0.11\,h}\ell+6+2\lambda+1+0.1
\le\frac{7.051\,h}\ell+7.1+2\lambda .\qedhere
\]
\end{proof}

\begin{lemma}[A thick letter and a short word]\label{pf:linkN-adm:thick}
In a setting for theorem~B with $\ell\ge11$, $20w,20\eta,20S\le b$, a non-thick letter and alternating type, if some letter is thick
and $b/4<40\,\mathrm{Sg}$, then $F$ lifts with an exponent $e_0\le160\,\mathrm{Sg}$.
\formallabel{Formalism.MainNotions.N.Ln17.exists_liftF_le_sg_of_thick}
\end{lemma}

\begin{proof}
Every letter has $\mathrm{sz}_a\le b+\eta$ (thick letters by definition, the others since $\mathrm{sz}_a\le w\le b$). At a
non-thick letter $a$ we have $\eta<\mathrm{sz}_a\le w$ and $20w\le b$, so $b\ge1$. The height is $h=b+\eta+S$. Define new
window bounds $M'^-_a=\eta$, $M'^+_a=\mathrm{sz}_a+S$ if $\mathrm{up}(a)$, and $M'^-_a=\mathrm{sz}_a+S$, $M'^+_a=\eta$
otherwise; the jumps of $T_1$ lie within $(M'^-,M'^+)$ because they lie in the windows and $\eta<\mathrm{sz}_a$. Then
$M'_{\max}\le h$, $H'_{\max}\le h+\eta$ and all jumps in the relevant range are in $[-M'_{\max},h]$. By (lift), $F$ lifts with
exponent $e_{PA}(\ell,M'_{\max},H'_{\max},J)$. As $h\ge b\ge1$ and $20\eta\le b\le h$, \cref{pf:linkN-adm:ePA} shows that
this exponent is at most $7.1\,h/\ell+7.1+2\log_\ell(2h)$; and \lref{Formalism.MainNotions.N.FrobeniusPowerSeries.step0} states that this
quantity is at most $160\,\mathrm{Sg}$ when $b\ge1$, $20w,20\eta,20S\le b$ and $b/4<40\,\mathrm{Sg}$.
\end{proof}

% ---------------------------------------------------------------------------
\subsection{The theorem on ramified coefficient chains (theorem C)}\label{pf:linkN-adm:thmCsec}

For a prime $\ell$ and $H\in\N$ put
\begin{equation}\label{pf:linkN-adm:eq-aR}
\mathrm{lam}(\ell,H)=\log_\ell(2H)+1,\qquad a_R(\ell,H)=\frac{H(2\ell-1)}{(\ell-1)^2}+3+\mathrm{lam}(\ell,H),\qquad
c_D(\ell,H)=\frac{H+1}{\ell-1}+1,
\end{equation}
with $\log_\ell0=0$, so that $\mathrm{lam}(\ell,H)\ge1$. Let $W$ be a Wach structure over $K$ with $n$ letters and
coefficients $L_C$.

\begin{definition}[The setting of theorem C]\label{pf:linkN-adm:def-setC}
Let $d\in\N$, $T_1,T_2$ lattices, $F\colon T_1\to T_2$ a lattice map, $b\colon\Z/n\to\N$ and $S,\eta,m\in\N$; put
$H_0=\max_ib_i+S$ and $H=H_0+\eta$. These form a \emph{setting of parity $d$ for theorem~C} if: $W$ is in case~A,
$\ell\ge11$ and $n$ is even; $\OO_C$ is a discrete valuation ring with fraction field $L_C$, and $L_C$ contains a primitive
$\ell$-th root of unity; $T_1,T_2$ are crystalline and $F$ is onto; for every letter $i$ and $\tau$ over $i$ the jumps of
$T_1$ at $\tau$ lie in
\[
\mathrm{W}^d_i=\begin{cases}[-(b_i+S),-b_i]\cup[-\eta,\eta]&\text{if }i+d\text{ is even},\\
[b_i,b_i+S]\cup[-\eta,\eta]&\text{if }i+d\text{ is odd;}\end{cases}
\]
$\sum_{|a|\le\eta}\mathrm{mult}_\tau(T_2,a)\le m$ for every $\tau$; and $2a_R(\ell,H)+\eta+1\le b_i$ and $H_0/2\le b_i$ for
all $i$. We then put $M^d_i=b_i+S$, $N^d_i=\eta$ if $i+d$ is even, $M^d_i=\eta$, $N^d_i=b_i+S$ otherwise, and
$e''=e_{PA}(\ell,\max_iM^d_i,\max_i(M^d_i+N^d_i),J)$ with $J$ the jumps of $T_1$ in $[-\max M^d,\max N^d]$.
\formal{Formalism.MainNotions.RamifiedCoefficientChains.SetCS}
\end{definition}

Since $\eta\le b_i$, the jumps of $T_1$ lie within $(M^d,N^d)$ in a setting of theorem~C, so the lifting theorem in
case~A shows that $F$ lifts with exponent $e''$.

\begin{theorem}[Theorem~C]\label{pf:linkN-adm:thmC}
Assume Bloch--Kato's local theory, local Tate duality and the local Euler characteristic formula. Let $D$ be a datum with
local Hodge--Tate, Fontaine and comparison data, $\ell$ a prime, $K/\Q_\ell$ finite with $n$ letters, $L_C$ admissible for
$K$ and $W$ the Wach structure. For every setting of parity $d$ for theorem~C there is $e\in\N$ with
\begin{gather*}
e\le e_{\mathrm{Bd}}=\frac{32H}{\ell^2}+4S+\frac{10\eta}{\ell-1}+25\,\mathrm{lam}(\ell,H),\\
\linkNadmcol(F)\preceq(n+1)re+4nm\,e''+\tfrac32n\bigl(rS+2m\,c_D(\ell,H)\bigr),
\end{gather*}
where $r=\rank T_2$, and such that $F$ lifts with exponent $e+3S$ if $m=0$.
\formallabel{GlobalTheory.T.prv_twenty_five_of_blochKato_tateDuality_eulerChar}
\formallabel{GlobalTheory.ramifiedCoefficientChains_prv_twenty_five}
\end{theorem}

\paragraph{The local lifting conditions.}
Let $\mathcal A_C$ denote the admissibility package of theorem~C with the constants $(10,23,32,32)$
(\cref{pf:linkN-adm:def-AC}); among its components are the generator bound (gen) and the lifting theorem in case~A (lift) of \cref{pf:linkN-adm:thmBsec}. At a place $(\ell,K,L_C)$
the \emph{local lifting conditions} are:

\begin{itemize}
\item[(LLC)] for every setting of parity $d$ for theorem~C there is $e_0\in\N$ such that: (a) there are $a,b'\in\N$ and
  reals $s_1,E,e_2$ with $b'\le H_0$, $a\le a_R(\ell,H)+1$,
  \[
  s_1\le\frac{2a+2\eta+S}{\ell-1}+2\mathrm{lam},\quad E\le\frac{3a+S+2\eta}{\ell-1}+3\mathrm{lam},\quad
  e_2\le\frac{b'}{\ell(\ell-1)}+E+2\sigma_2+2S+5+\mathrm{lam},
  \]
  and $e_0\le s_1+2\sigma_1+e_2+S+E+\sigma_2$, where $\mathrm{lam}=\mathrm{lam}(\ell,H)$,
  $\sigma_1=(\lfloor H_0/(\ell-1)\rfloor+1+\eta)/(\ell-1)+1$ and $\sigma_2=2H/(\ell-1)^2+1$; (b) for every free $\OO_C$-module
  $\Phi$ of finite rank and every $\OO_C$-linear map $g\colon\Phi\to\linkNadmZf(T_2)$ which is surjective modulo
  $B^1(T_2)$, there is a \emph{kernel chain}: submodules $\Phi=\mathrm{Fl}_0\supseteq\dots\supseteq\mathrm{Fl}_k$ and a set
  $\mathrm{gen}\subseteq\{0,\dots,k-1\}$ with $|\mathrm{gen}|\le4n$ and $k-|\mathrm{gen}|\le\lfloor3n/2\rfloor$, such that
  each $\mathrm{Fl}_i/\mathrm{Fl}_{i+1}$ embeds into a module which is generated by at most $m$ elements if
  $i\in\mathrm{gen}$ and has length at most $c\,e_C$ otherwise, where $c=rS+2m\,c_D(\ell,H)$, and
  $\ell^{e_0}g(\mathrm{Fl}_k)\subseteq\linkNadmIm(F)$; (c) if $m=0$, $F$ lifts with exponent $e_0+3S$.
\end{itemize}

\begin{definition}[Three local statements]\label{pf:linkN-adm:def-three}
Let $W$ be a Wach structure with $n$ letters. The following statements are quantified over all settings
$(d,T_1,T_2,F,b,S,\eta,m)$ of parity $d$ for theorem~C at $W$; the letters $c$ with $c+d$ even are called \emph{even}, and
$s_{CW}$, $e_{WV}$ are explicit functions of $(\ell,H_0,\eta,S)$ defined in the formal development. A \emph{chain with
$(g_1,g_2)$ steps} from $\Phi$ to $\Phi'$ is a chain $\Phi=\mathrm{Fl}_0\supseteq\dots\supseteq\mathrm{Fl}_k=\Phi'$ as in
(LLC)(b) (each quotient embeds into a module generated by at most $m$ elements, or of length at most $c\,e_C$ with
$c=rS+2m\,c_D(\ell,H)$) with at most $g_1$ steps of the first kind and at most $g_2$ of the second.
\begin{itemize}
\item[(CW)] \emph{Cocycles to Wach vectors}: there is $s\le s_{CW}$ such that for every finite free $\OO_C$-module $\Phi$ and
  every $\OO_C$-linear $g\colon\Phi\to\linkNadmZf(T_2)$ there are a submodule $\Phi_2\subseteq\Phi$ reached by a chain with
  $(n,0)$ steps, with $\Phi_2=\Phi$ if $m=0$, and $\OO_C$-linear maps $w_c\colon\Phi_2\to D(T_2)$ for the even letters $c$,
  such that each $w_c(f)$ satisfies the \emph{$q$-power condition} at $c$ ($q^\eta w_c(f)$ lies in the pull-back lattice at
  $c-1$, and $w_c(f)$ satisfies the integrality condition at $c$), and for every $f\in\Phi_2$ there are cocycles
  $\kappa_c$ attached to $w_c(f)$ (as classes of truncated $\Delta$-Herr cocycles) such that $\ell^sg(f)-\sum_c\kappa_c$ or
  $\ell^sg(f)+\sum_c\kappa_c$ is a coboundary.
\item[(WV)] \emph{Chains of Wach vectors}: there is $e\le e_{WV}$ such that for every even letter $c$, every submodule
  $\Phi_s$ of a finite free module and every $\OO_C$-linear $w\colon\Phi_s\to D(T_2)$ whose values satisfy the $q$-power
  condition at $c$, there is $\Phi'\subseteq\Phi_s$ reached by a chain with $(6,3)$ steps such that for all $f\in\Phi'$ and
  $j\in\N$ there is $u\in D(T_1)$ satisfying the integrality condition at $c$ with $\ell^{j+e}w(f)-F(u)\in N(T_2)$; and if
  $m=0$, the same holds for all $f\in\Phi_s$ with $\ell^{j+e+3S}$ in place of $\ell^{j+e}$.
\item[(RL)] \emph{Lifting of attached cocycles}: for every even letter $c$, $w\in D(T_2)$, $u\in D(T_1)$ satisfying the
  integrality condition at $c$, cocycle $\kappa$ attached to $w$ and $k\in\N$: if $\ell^kw-F(u)\in N(T_2)$, then
  $\ell^k\kappa\in\linkNadmIm(F)$.
\end{itemize}
We say that one of these statements holds \emph{at every good place} if it holds at the Wach structures of level $1$ of all
places $(\ell,K,L_C)$.
\formal{Formalism.MainNotions.RamifiedCoefficientChains.S9.CocyclesToWachVectorsBound}
\formal{Formalism.MainNotions.RamifiedCoefficientChains.S9.WachVectorSubmoduleChainBound}
\formal{Formalism.MainNotions.RamifiedCoefficientChains.S9.RhoLift}
\end{definition}

At an admissible place the three statements imply (LLC); this is
\lref{Formalism.MainNotions.RamifiedCoefficientChains.S9.lattice_bounds_of_admissible}, proved in the formal development
and not reproduced here: it composes the chains of (CW) and (WV), lifts the vectors $w_c(f)$ by (WV) and the attached
cocycles by (RL), and bounds the resulting exponent $e_0$ as in (LLC)(a).

\begin{lemma}\label{pf:linkN-adm:e0}
If $\ell\ge11$ and $e_0$ satisfies (a) of (LLC), then $e_0\le e_{\mathrm{Bd}}$.
\formallabel{Formalism.MainNotions.RamifiedCoefficientChains.e0_le_eBd_of_ExpC}
\end{lemma}

\begin{proof}
Adding the bounds of (a) (note that $e_2$ contains $E$ once more),
\[
e_0\le\frac{8a+6\eta+3S}{\ell-1}+9\,\mathrm{lam}+2\sigma_1+\frac{H_0}{\ell(\ell-1)}+3\sigma_2+3S+5 .
\]
Insert $a\le H(2\ell-1)/(\ell-1)^2+4+\mathrm{lam}$, $2\sigma_1\le2H_0/(\ell-1)^2+(2+2\eta)/(\ell-1)+2$, $3\sigma_2=
6H/(\ell-1)^2+3$ and $H_0\le H$. The coefficient of $H$ is then at most $\ell^{-2}f(\ell)$ with
$f(\ell)=8(2\ell-1)\ell^2/(\ell-1)^3+8\ell^2/(\ell-1)^2+\ell/(\ell-1)$, which is decreasing for $\ell\ge11$ with
$f(11)<31.2\le32$. The coefficient of $\eta$ is $8/(\ell-1)\le10/(\ell-1)$, that of $S$ is $3/(\ell-1)+3\le4$, that of
$\mathrm{lam}$ is $9+8/(\ell-1)\le9.8$, and the constant term is $34/(\ell-1)+10\le13.4\le13.4\,\mathrm{lam}$. Since
$9.8+13.4\le25$, $e_0\le e_{\mathrm{Bd}}$.
\end{proof}

\begin{lemma}[Length count]\label{pf:linkN-adm:lengthcount}
Let $\OO_C$ be a discrete valuation ring with $\ell\ne0$ in $\OO_C$, $\ell\ge2$, $F\colon T_1\to T_2$ a lattice map lifting with
exponent $e'$, $n,r,m,e\in\N$ and $c\ge0$. Let $\Phi$ be a free $\OO_C$-module of rank at most $(n+1)r$ and
$g\colon\Phi\to\linkNadmZf(T_2)$ surjective modulo $B^1(T_2)$, and let $\mathrm{Fl},\mathrm{gen}$ be a kernel chain for
$(\Phi,g)$ with the constants $m,c$ and exponent $e$ as in (LLC)(b). Then
$\linkNadmcol(F)\preceq(n+1)re+4nm\,e'+\frac32n\,c$.
\formallabel{Formalism.MainNotions.RamifiedCoefficientChains.colLeR_of_lengthCount}
\formallabel{Formalism.MainNotions.RamifiedCoefficientChains.lengthCount_of_kerChain}
\end{lemma}

\begin{proof}
Let $P=g^{-1}(\linkNadmIm(F))\subseteq\Phi$ and $e_C=\linkNadmlen(\OO_C/\ell)$, finite as $\ell\ne0$. The map
$\Phi/P\to\linkNadmZf(T_2)/(\linkNadmIm(F)\cap\linkNadmZf(T_2))$ induced by $g$ is injective by definition of $P$ and surjective
because $g$ is surjective modulo $B^1(T_2)\subseteq\linkNadmIm(F)$; so $\linkNadmcol(F)\le\linkNadmlen(\Phi/P)$. We apply the
filtration count (\cref{pf:linkN-adm:filtration}) with $a=\ell$ and $L=\lfloor c\,e_C\rfloor$. Condition (ii) holds since $F$ lifts
with exponent $e'$, so $\ell^{e'}g(\Phi)\subseteq\linkNadmIm(F)$; (iv) is the last property of the kernel chain. For
$i\in\mathrm{gen}$, $\mathrm{Fl}_i/\mathrm{Fl}_{i+1}$ is a submodule of a module generated by $m$ elements over a principal ideal
domain, hence generated by at most $m$ elements, which gives (i). For $i\notin\mathrm{gen}$,
$\mathrm{Fl}_i/((\mathrm{Fl}_{i+1}+P)\cap\mathrm{Fl}_i)$ is a quotient of $\mathrm{Fl}_i/\mathrm{Fl}_{i+1}$, whose length is at most
$c\,e_C$, which gives (iii) with $L$. Hence
\[
\linkNadmlen(\Phi/P)\le\rank(\Phi)\,e\,e_C+|\mathrm{gen}|\,m\,e'e_C+(k-|\mathrm{gen}|)\,c\,e_C\le
\Bigl((n+1)re+4nme'+\tfrac32nc\Bigr)e_C .\qedhere
\]
\end{proof}

\begin{proposition}[Theorem~C from the local lifting conditions]\label{pf:linkN-adm:thmCLLC}
If the generator bound (gen), the lifting theorem in case~A and (LLC) hold at every admissible place, then the conclusion
of \cref{pf:linkN-adm:thmC} holds with the constant $25$.
\formallabel{Formalism.MainNotions.RamifiedCoefficientChains.theorem_of_lifting_steps}
\formallabel{Formalism.MainNotions.RamifiedCoefficientChains.Theorem_C_of_open}
\formallabel{Formalism.MainNotions.RamifiedCoefficientChains.Theorem_C_range_of_open}
\formallabel{Formalism.MainNotions.RamifiedCoefficientChains.thmCKR_of_open}
\formallabel{Formalism.MainNotions.RamifiedCoefficientChains.openKR_of_steps}
\formallabel{Formalism.MainNotions.RamifiedCoefficientChains.openC_of_steps}
\end{proposition}

\begin{proof}
Let a setting of parity $d$ for theorem~C be given. By (gen), $\linkNadmHf(T_2)$ is generated by at most $(n+1)r$ classes;
choose finite cocycles $z_1,\dots,z_s$ representing them ($s\le(n+1)r$), let $\Phi=\OO_C^s$ and $g(x)=\sum x_jz_j$. Then every
finite cocycle differs from an element of $g(\Phi)$ by a coboundary. Let $e_0$ be given by (LLC) and take the kernel chain
for $(\Phi,g)$. By \cref{pf:linkN-adm:e0}, $e_0\le e_{\mathrm{Bd}}$. By the remark after \cref{pf:linkN-adm:def-setC}, $F$
lifts with exponent $e''$, and \cref{pf:linkN-adm:lengthcount} with $e'=e''$ and $c=rS+2mc_D(\ell,H)$ gives the bound for
$\linkNadmcol(F)$. The last assertion is (LLC)(c). Take $e=e_0$.
\end{proof}

\begin{proposition}[Theorem~C from the local statements]\label{pf:linkN-adm:thmC3}
Let $D$ be a datum with local Hodge--Tate, Fontaine and comparison data, and $C_R,C_A,C_B,C_K\in\R$.
\begin{enumerate}[label=(\alph*)]
\item If $\mathcal A_C$ implies (LLC) at every admissible place, then $\mathcal A_C$ implies the conclusion of
  \cref{pf:linkN-adm:thmC} with the constant $25$.
\item If $\mathcal A_C$ implies (CW), (WV) and (RL) at every good place, the same holds.
\end{enumerate}
\formallabel{Formalism.MainNotions.RamifiedCoefficientChains.stmt_of_stepsAdm}
\formallabel{Formalism.MainNotions.RamifiedCoefficientChains.S9.stmt_of_lemmas}
\end{proposition}

\begin{proof}
(a) Assume $\mathcal A_C$. Its component (C10) contains the lifting theorem in case~A, and its component (C19) is the
generator bound (gen) (in an equivalent formulation). Apply \cref{pf:linkN-adm:thmCLLC}. (b) By
\lref{Formalism.MainNotions.RamifiedCoefficientChains.S9.lattice_bounds_of_admissible}, the hypothesis of (b) implies that
of (a).
\end{proof}

\begin{proposition}[The chain bound and the lifting of cocycles]\label{pf:linkN-adm:TLL}
Let $D$ be a datum with local Hodge--Tate, Fontaine and comparison data and $C_R,C_A,C_B,C_K\in\R$. Assume that $\mathcal
A_C$ implies (CW) at every good place, that the jumps behave as in (cj) at every place, and the global letter condition.
Then $\mathcal A_C$ implies the conclusion of \cref{pf:linkN-adm:thmC} with the constant $25$.
\formallabel{Formalism.MainNotions.RamifiedCoefficientChains.Asm.theorem_of_local_lemmas}
\end{proposition}

\begin{proof}
By \cref{pf:linkN-adm:thmC3}(b) it suffices to show that $\mathcal A_C$ implies (WV) and (RL) at every good place. Assume
$\mathcal A_C$. (RL) holds at every good place by
\lref{Formalism.MainNotions.RamifiedCoefficientChains.Asm.rhoLift_of_Adm}, which derives it from $\mathcal A_C$ (through a
truncation clause and a local clause of the crystalline lifting inputs). For (WV) we use four local lemmas, each stated
for all settings of theorem~C at a Wach structure $W$ of a good place:
\begin{itemize}
\item[(Gs)] at the letters $i$ with $i+d$ odd there is $E\le(3a_R(\ell,H)+S+2\eta)/(\ell-1)+\mathrm{lam}(\ell,H)$ such that
  every $x\in\varphi^2(D_{i-2}(T_2))\cap\mathrm{Fil}^{b_i}D_{\mathrm{cris}}$ satisfies
  $p^{\lceil E+\sigma_2+b_{i-1}+S-b_i\rceil}x\in F(D_i(T_1))$, where $D_i(T)$ denotes the lattice in $D_{\mathrm{cris}}$
  attached to the letter $i$ and $\sigma_2$ is as in (LLC);
\item[(DN)] at every even letter $i$ there is a set of at most $m$ elements of $N_i(T_2)$ such that the quotient of $N_i(T_2)$
  by $\pi_0N_i(T_2)$, by the $q^{b_i+S}$-multiples of its Frobenius pull-back and by these elements has length at most
  $\mathrm{len}_C\cdot\linkNadmlen(\OO_C/\ell)$, for an explicit $\mathrm{len}_C$;
\item[(DN0)] if $m=0$, then at every even letter $i$, $p^S N_i(T_2)$ lies in the sum of $\pi_0N_i(T_2)$ and the
  $q^{b_i+S}$-multiples of the Frobenius pull-back;
\item[(H2)] an $H^2$-type bound for the lattices of the setting.
\end{itemize}
They hold at every good place: (Gs) by \lref{Formalism.MainNotions.RamifiedCoefficientChains.Asm.lemmaGs_rec} (from
$\mathcal A_C$, (cj) and the letter condition); (DN) by
\lref{Formalism.MainNotions.RamifiedCoefficientChains.S9.bigCellDetWeightedEnergyIntegral.dual_lemma_of_admitted} (a
duality lemma for determinants, from $\mathcal A_C$); (DN0) by
\lref{Formalism.MainNotions.RamifiedCoefficientChains.S9.DLev.lemmaDN0_L} (from the components (C1) and (C11) of
$\mathcal A_C$); (H2) by \lref{Formalism.MainNotions.RamifiedCoefficientChains.Asm.h2_bound_of_admitted} (from $\mathcal
A_C$). From them, \lref{Formalism.MainNotions.RamifiedCoefficientChains.Asm.frobenius_lift_lemmas_of_admitted} derives the
\emph{Wach-vector lifting bound} at every good place, and \lref{Formalism.MainNotions.N.Pv.LocalLengthBound.steps6_of}
derives, from $\mathcal A_C$, (Gs), (DN) and (DN0), a six-step chain statement at every good place. At a good place,
\lref{Formalism.MainNotions.RamifiedCoefficientChains.S9.lattice_bound_three_six_of} combines these two statements into
(WV).
\end{proof}

\begin{proposition}[From cocycles to Wach vectors]\label{pf:linkN-adm:CW}
Let $D$ be a datum with local Hodge--Tate, Fontaine and comparison data and $C_R,C_A,C_B,C_K\in\R$. Assume $\mathcal A_C$,
the statement (cj) at every place, the natural extension of Herr classes, Fontaine's inequality
$\dim_{L_C}D_{\mathrm{cris}}(V)\le[K_0:\Q_\ell]\dim_{L_C}V$ at every place, and Fontaine's full faithfulness statement. Then
(CW) holds at every good place.
\formallabel{Formalism.MainNotions.RamifiedCoefficientChains.S9.LocalLemmasFromLiterature.local_lemmas_of_literature}
\end{proposition}

\begin{proof}
Let $W$ be the Wach structure of a good place $(\ell,K,L_C)$ and fix a setting of theorem~C at $W$; then $W$ is in case~A,
$\ell\ge11$, and $K$ is unramified. By \lref{Formalism.MainNotions.RamifiedCoefficientChains.S9.ExtensionLemmaAlt.cocyclesToWachVectorsBound_of_extTool},
(CW) at $W$ (for this setting) follows from three local statements at $W$: a \emph{first lifting step}, a \emph{second lifting
step} and an \emph{extension tool}. They are obtained as follows.
\begin{itemize}
\item The first lifting step is \lref{Formalism.MainNotions.RamifiedCoefficientChains.S9.ExtensionLemmaAlt.first_lift_of_uniformizer_clause}
  (a lifting lemma for killing polynomials, from the valuation property of $\OO_C$, the clause of (P9) on $\pi_0$, a
  numerical tool of the crystalline lifting inputs, (P9), the statement (C18) on $(\gamma-1)/\pi_0$ at the letters, the
  exponent rules and exponent ring of $W$, the lemma K, the filtration dimension statement, the basis
  tool, and the three clauses of Fontaine's statement on $D_{\mathrm{cris}}$ (C11)).
\item The second lifting step is \lref{Formalism.MainNotions.RamifiedCoefficientChains.S9.ExtensionLemmaAlt.second_lift_of_hypotheses}
  (from a truncated-series tool, a Frobenius tool, the basis tool, the lemma K, the filtration dimension statement, the
  exponent ring and rules, a clause of (C3), the clauses (P5) and (P9) contained in (C1), and (C11)).
\item The extension tool is \lref{Formalism.MainNotions.RamifiedCoefficientChains.Sh.extData_extTool} (from the natural
  extension of Herr classes, the extension and unit statements for the lifting of classes
  \lref{Formalism.MainNotions.LiftingTheory.liftb_HExt_holds}, \lref{Formalism.MainNotions.LiftingTheory.liftb_HOne_holds},
  a kernel statement \lref{Formalism.MainNotions.LiftingTheory.lker_HKer_D} derived from Fontaine's inequality, (cj),
  Fontaine's full faithfulness, the invertibility of $|\Delta|$ in $\OO_C$, the clauses (P5), (P6), (P9) contained in (C1),
  and a level tool).
\end{itemize}
The auxiliary statements come from $\mathcal A_C$ through the following results: the exponent rules from the level map
statement of the level tools of $\mathcal A_C$ and the $\gamma$-structure attached to the comparison data
(\lref{Formalism.MainNotions.N.Pv.LocalLengthBound.expRules_of}, \lref{Formalism.MainNotions.N.Pv.LocalLengthBound.gam_of_tie});
the exponent ring from the ring clause of (P9) (\lref{Formalism.MainNotions.N.Pv.LocalLengthBound.expRing_of}); the basis
tool from (P9) (\lref{Formalism.MainNotions.RamifiedCoefficientChains.Sh.basTool_of}); the lemma K
(\lref{Formalism.MainNotions.N.Pv.LocalLengthBound.lemK_D}, from $\mathcal A_C$ and the level coefficients at the points
given by the uniformiser clause of $\mathcal A_C$, \lref{Formalism.MainNotions.UnusedSlot80.points_exist_of_ring_condition},
\lref{Formalism.MainNotions.UnusedSlot80.levCoeff_of_pt}); the filtration dimension statement from the third clause of (C11)
(\lref{Formalism.MainNotions.N.Pv.LocalLengthBound.filDim_of}); the truncated-series tool
(\lref{Formalism.MainNotions.CrystallineLiftingInputs.truncatedSeriesData.tTool2Z_DU}); the Frobenius and level tools
(\lref{Formalism.MainNotions.CrystallineLiftingInputs.Tools.local_clause_of_admitted_statements},
\lref{Formalism.MainNotions.CrystallineLiftingInputs.Tools.level_clause_of_admitted_statements}); and the numerical tool
(\lref{Formalism.MainNotions.CrystallineLiftingInputs.RToolNum_holds}). Finally $|\Delta|=\ell-1$ in case~A
(\lref{Formalism.MainNotions.UnusedSlot80.LevR.card_delF}), and $\ell-1$ is a unit of $\OO_C$: otherwise it would lie in
the maximal ideal, which contains $\ell$, and so would $1=\ell-(\ell-1)$.
\end{proof}

\begin{proposition}[Reduction of theorem~C]\label{pf:linkN-adm:thmCred}
Assume the exactness of jumps (the statement (cj) at every place), the global letter condition, the natural extension of
Herr classes, Fontaine's inequality $\dim_{L_C}D_{\mathrm{cris}}(V)\le[K_0:\Q_\ell]\dim_{L_C}V$ at every place, and the full
faithfulness and exactness statement for Fontaine's functor. Then $\mathcal A_C$ implies the conclusion of
\cref{pf:linkN-adm:thmC} with the constant $25$.
\formallabel{GlobalTheory.T.ramified_coefficient_chains_of_hypotheses}
\formallabel{GlobalTheory.theoremC_of_admitted_inputs}
\formallabel{Formalism.MainNotions.RamifiedCoefficientChains.Asm.theorem_of_literature_hypotheses}
\end{proposition}

\begin{proof}
By \cref{pf:linkN-adm:CW}, whose hypotheses besides $\mathcal A_C$ are among the assumptions, $\mathcal A_C$ implies (CW)
at every good place. \Cref{pf:linkN-adm:TLL}, applied with (cj) and the global letter condition, concludes.
\end{proof}

\begin{proof}[Proof of \cref{pf:linkN-adm:thmC}]
The package $\mathcal A_C$ is proved from Bloch--Kato's local theory, local duality and the Euler characteristic formula in
\cref{pf:linkN-adm:admC}. The other inputs of \cref{pf:linkN-adm:thmCred} are:
the exactness of jumps \lref{GlobalTheory.jumps_exact_adm}; the global letter condition, from
\cref{pf:linkN-adm:letter} and \lref{GlobalTheory.T.wachLattice_laterClauses_hold} as in the proof of theorem~B; the natural extension of Herr classes \lref{GlobalTheory.AdmHerrLocal.main}, which provides at
every place the formula for the Herr class of an extension; the one-place Fontaine inequality
\lref{GlobalTheory.FontaineIneq.one_place_admissible_of_gap_constant_le}; and
\lref{GlobalTheory.T.fontaine_fullFaithful_exact_psi}, which proves the full faithfulness and exactness properties of
Fontaine's functor used by the local lemmas. Apply \cref{pf:linkN-adm:thmCred}.
\end{proof}

% ---------------------------------------------------------------------------
\subsection{Admissibility of the local inputs}\label{pf:linkN-adm:admsec}

This subsection proves the admissibility packages $\mathcal A_B$ and $\mathcal A_C$ of theorems~B and~C from
Bloch--Kato's local theory, local Tate duality and the local Euler characteristic formula, and the global letter condition
from the norm bound for $\Delta$-invariant lattices. The proofs assemble the components from results of the formal
development; the results with long proofs are cited with a statement of what they give. Three results proved elsewhere are
used throughout: the \emph{later clauses of the Wach lattice theory}
\lref{GlobalTheory.T.wachLattice_laterClauses_hold}, which prove the norm bound for $\Delta$-invariant lattices
together with its clauses, among them the freeness (P1) of the Wach lattices and the clauses (P2)--(P9); the \emph{local
constants} \lref{GlobalTheory.T.localConstants_from_admitted_local_Galois_statements}, which prove from the three
hypotheses the local lifting theory with constants $(10,23,32,32)$, containing (H), (Tr), an annihilator statement (W1) for
Wach modules and the statement (WX) that Wach modules kill the relevant quotients; and the \emph{arc tools}
\lref{GlobalTheory.T.arcTheorem_tools_of_wachModuleConclusions}, which prove (sat) and (sat$'$), the statement
$g(\pi_0)\in\pi_0A^\Delta$ for all $g\in\Gamma$, (kill1), the killer statements with constant $8$, and two filtration statements (F1), (F2),
at every place in case~A with $\ell\ge5$.

\begin{definition}[The admissibility package of theorem~B]\label{pf:linkN-adm:def-AB}
Let $D$ be a datum with local Hodge--Tate, Fontaine and comparison data and $C_R,C_A,C_B,C_K,C_\kappa\in\R$. The package
$\mathcal A_B$ consists of: (B1) the used local conclusion of the Wach lattice theory; (B2) the local lifting theory with
constants $(C_R,C_A,C_B,C_K)$, which contains (wr), (H), (Tr) and the lifting theorem in case~A; (B3) (sat) and (B4) (sat$'$)
at every place in case~A with $\ell\ge5$; (B5) the killer statements: two auxiliary statements on window sums, (kill1) and
(kill2) at these places, and a bound for $\kappa$ with constant $C_\kappa$; (B6) the content divisibility statement (r2);
(B7) the generator bound (gen); (B8) (relay), (B9) (arc) and (arc$'$), (B10) (split), (B11) (nn) and (univ), each at every
place in case~A with $\ell\ge5$; (B12) (dc) at every place; and (B13) the Bloch--Kato crystalline extension criterion.
\formal{Formalism.MainNotions.ArcLiftingSetup.admitted}
\end{definition}

\begin{lemma}[The relay statement]\label{pf:linkN-adm:relaylem}
Let $D$ be a datum with local Hodge--Tate and Fontaine data, and assume (P5), (P9), (H) and (Tr) at every place, and (sat),
(F1), (F2) at every place in case~A with $\ell\ge5$. Then (relay) holds at every place in case~A with $\ell\ge5$.
\formallabel{Formalism.MainNotions.UnusedSlot80.relay_general_statement}
\formallabel{Formalism.MainNotions.UnusedSlot80.relay_general_case_of_admissible}
\end{lemma}

\begin{proof}
Let $W$ be the Wach structure of such a place. \Lref{Formalism.MainNotions.UnusedSlot80.Relay2.relayCore_of} proves (relay)
at $W$ from: the injectivity of $\OO_C\to L_C$ and the fact that every element of $L_C$ of nonnegative valuation comes from
$\OO_C$ (true for the ring of integers of $L_C$), (P9), (P5), (H), (Tr), (sat), (sat$'$), (F1) and (F2) at $W$. Here (sat$'$) at
$W$ holds by \lref{Formalism.MainNotions.UnusedSlot80.SatPW.satPCore_of}, which derives it from (P9) and the injectivity
of $\OO_C\to L_C$.
\end{proof}

\begin{proposition}[Relay statement and arc package]\label{pf:linkN-adm:relay}
Assume Bloch--Kato's local theory, local Tate duality and the local Euler characteristic formula, and let $D$ be a datum with
local Hodge--Tate, Fontaine and comparison data. Then (relay) holds at every place in case~A with $\ell\ge5$, and the
\emph{arc package} $\mathcal A_{\mathrm{arc}}$ holds: (P3), (P4), (P5), (P9), (H), (Tr), (W1), and (sat) at every place in
case~A with $\ell\ge5$.
\formallabel{GlobalTheory.T.wachRelayLemma_arcCoordinateHypotheses}
\formallabel{GlobalTheory.NQ.relay_and_arc_inputs}
\end{proposition}

\begin{proof}
(P3), (P4), (P5) and (P9) are among the later clauses of the Wach lattice theory; (H), (Tr) and (W1) are part of the local
constants; (sat), (F1) and (F2) are among the arc tools. This gives $\mathcal A_{\mathrm{arc}}$ and the hypotheses of
\cref{pf:linkN-adm:relaylem}, which gives (relay).
\end{proof}

\begin{proposition}[Class splitting along arcs]\label{pf:linkN-adm:splitting}
Let $D$ be a datum with local Hodge--Tate, Fontaine and comparison data. Assume the \emph{splitting package} ((P8), (P9),
(W1), (cj) at every place, the Herr extension class formula, (F1), (relay) and the arc statements (B9)), the arc package
$\mathcal A_{\mathrm{arc}}$, (P1), (P6), (WX), Fontaine's full faithfulness statement, (cj) at every place, (dc) at every
place, and Fontaine's statement on $D_{\mathrm{cris}}$ (the component (C11) of $\mathcal A_C$). Then (split) holds at every
place in case~A with $\ell\ge5$.
\formallabel{GlobalTheory.T.classSplitting_along_arcs_of_admissible}
\formallabel{GlobalTheory.classSplitAlongArcs}
\formallabel{Formalism.MainNotions.UnusedSlot80.split_case_A_final}
\end{proposition}

\begin{proof}
The proof passes through statements on the \emph{level-zero parts} of the coefficients relative to horizontal families.
First, \lref{Formalism.MainNotions.N.Pv.LocalLengthBound.memCritL_of} gives a membership criterion at level zero from
(P1), the clauses (P3), (P4), (P5), (P9) of $\mathcal A_{\mathrm{arc}}$ and (C11), and
\lref{Formalism.MainNotions.UnusedSlot80.lev0MprL_of} turns it, with $\mathcal A_{\mathrm{arc}}$, into a refined
level-zero statement modulo $p$. Second, \lref{Formalism.MainNotions.UnusedSlot80.lev0MpL_of} gives a level-zero
membership statement from $\mathcal A_{\mathrm{arc}}$, (P4) and (P9). Third,
\lref{Formalism.MainNotions.UnusedSlot80.splitShapeL_rec} derives from $\mathcal A_{\mathrm{arc}}$, these two
statements, (cj) and (dc) the \emph{splitting shape}: a decomposition of finite classes along the cuts in the coordinates of
the arcs. Finally, \lref{Formalism.MainNotions.UnusedSlot80.splitting_of_literature_exactness} derives (split) from the
splitting package, (P6), (H), (Tr) (from $\mathcal A_{\mathrm{arc}}$), (WX), Fontaine's full faithfulness statement and
the splitting shape.
\end{proof}

\begin{proposition}[The universal statements]\label{pf:linkN-adm:univ}
Let $D$ be a datum with local Hodge--Tate, Fontaine and comparison data. Assume the \emph{universal package} ((sat) and
(kill1) at every place in case~A with $\ell\ge5$, the statement $g(\pi_0)\in\pi_0A^\Delta$ for all $g\in\Gamma$ at these
places, and the arc statements (B9)), (P9),
(H) at every place, and the equality of $D_{\mathrm{dR}}$ and $D_{\mathrm{cris}}$ at every place. Then (nn) and (univ)
hold at every place in case~A with $\ell\ge5$.
\formallabel{GlobalTheory.T.crystallineUniversalPart_caseA}
\formallabel{GlobalTheory.universalPart_caseA}
\formallabel{Formalism.MainNotions.UnusedSlot80.universal_general_statement_holds}
\end{proposition}

\begin{proof}
By \lref{Formalism.MainNotions.UnusedSlot80.invLevCore_fin}, (P9) implies the following \emph{inverse level statement}
at every place in case~A with $\ell\ge5$: for a letter $a$, an integer $s$, a finite set $T\not\ni s$ and exponents
$e(t)$, if $f\in E_W$ lies in the fractional ideal $e_a\bigl(\prod_{t\in T}q_t^{e(t)}\bigr)^{-1}A^\Delta$ and
$p^\kappa(\gamma_E-1)f$ lies in $q_s$ times this ideal for some $\kappa$, then $p^\kappa f-e_ax$ lies in $q_s$ times this
ideal for some $\kappa$ and some $x\in\OO_C$. \Lref{Formalism.MainNotions.UnusedSlot80.universal_general_of_level_inverse}
derives (nn) and (univ) from the universal package, (P9), (H), the equality of $D_{\mathrm{dR}}$ and $D_{\mathrm{cris}}$ and
the inverse level statement.
\end{proof}

\begin{theorem}[Admissibility for theorem~B]\label{pf:linkN-adm:admB}
Assume Bloch--Kato's local theory, local Tate duality and the local Euler characteristic formula. Then for every datum $D$
with local Hodge--Tate, Fontaine and comparison data, $\mathcal A_B$ holds with the constants $(10,23,32,32,8)$.
\formallabel{GlobalTheory.T.localLifting_theoremB_hypotheses_admitted}
\formallabel{GlobalTheory.TheoremBAdm.admissible_ranks_of_CMDatum}
\end{theorem}

\begin{proof}
(B1) is part of \lref{GlobalTheory.T.wachModuleLocalConclusions_usedForms} (the used forms of the local conclusions of the
Wach lattice theory); (B2) is the local constants; (B3), (B4), (B5) are among the arc tools; (B6) is
\lref{Formalism.MainNotions.ArcLiftingSetup.content_divisibility_clause}; (B7) is
\lref{GlobalTheory.T.h1fGeneratorBoundAll_loc} (the generator bound at every place); (B12) is
\lref{GlobalTheory.T.dcris_exact_strict_elementwise_crystalline}; and (B13) is
\lref{GlobalTheory.T.blochKato_crystalline_extension_criterion}. By \cref{pf:linkN-adm:relay} we have (B8) and
$\mathcal A_{\mathrm{arc}}$. (B9) follows from $\mathcal A_{\mathrm{arc}}$, (P1) and Fontaine's statement on
$D_{\mathrm{cris}}$ (\lref{GlobalTheory.T.ramifiedCoefficientChains_adm11}) by
\lref{GlobalTheory.T.arcCoordinatesCaseA_of_admArcA_of_P1free_of_adm11}. (B10) follows from
\cref{pf:linkN-adm:splitting}: the splitting package consists of (P8), (P9), (W1), (cj) (proved at every place by
\lref{GlobalTheory.jumps_exact_adm}), the Herr extension class formula
(\lref{GlobalTheory.HerrExtClass.herr_extension_class_formula}), (F1), (B8) and (B9); the other hypotheses are
$\mathcal A_{\mathrm{arc}}$, (P1), (P6), (WX), Fontaine's full faithfulness statement
(\lref{GlobalTheory.T.fontaine_fullFaithful_exact_psi}), (cj), (B12) and Fontaine's statement on $D_{\mathrm{cris}}$.
(B11) follows from \cref{pf:linkN-adm:univ}, the universal package consisting of (sat), (kill1), the statement $g(\pi_0)\in\pi_0A^\Delta$ ($g\in\Gamma$)
and (B9), together with (P9), (H) and the equality of $D_{\mathrm{dR}}$ and $D_{\mathrm{cris}}$
(\lref{GlobalTheory.T.dDeRham_eq_dCris_of_unramified}).
\end{proof}

\begin{definition}[The admissibility package of theorem~C]\label{pf:linkN-adm:def-AC}
Let $D$ be a datum with local Hodge--Tate, Fontaine and comparison data and $C_R,C_A,C_B,C_K\in\R$. The package $\mathcal A_C$
consists of: (C1)--(C4) four clauses of the used local conclusion of the Wach lattice theory; (C5)--(C10) six components of
the local lifting theory with constants $(C_R,C_A,C_B,C_K)$ (a bound for sizes, a restriction dimension statement, the level
tools, the hypothesis of a lemma on $M'$, a crystalline lifting statement, and the lifting bounds, which contain the lifting
theorem in case~A); (C11) Fontaine's statement on $D_{\mathrm{cris}}$ (three clauses on exactness and comparison); (C12) the
Bloch--Kato dimension formula at all local situations; (C13) generators of $\linkNadmHf$; (C14) Herr's description of $H^1$;
(C15) the Bloch--Kato crystalline extension criterion; (C16) the Herr extension class formula; (C17) the naturality of the
Herr complex; (C18) a statement on $(\gamma-1)/\pi_0$ at the letters; and (C19) the generator bound (gen).
\formal{Formalism.MainNotions.RamifiedCoefficientChains.admitted}
\end{definition}

\begin{theorem}[Admissibility for theorem~C]\label{pf:linkN-adm:admC}
Assume Bloch--Kato's local theory, local Tate duality and the local Euler characteristic formula. Then for every datum $D$
with local Hodge--Tate, Fontaine and comparison data, $\mathcal A_C$ holds with the constants $(10,23,32,32)$.
\formallabel{GlobalTheory.T.localLifting_theoremC_hypotheses_admitted}
\formallabel{GlobalTheory.TheoremCAdm.admissible_ranks_of_CMDatum}
\end{theorem}

\begin{proof}
(C1)--(C4) are \lref{GlobalTheory.T.wachModuleLocalConclusions_usedForms}; (C5)--(C10) are components of the local
lifting theory given by the local constants; (C11) is \lref{GlobalTheory.T.ramifiedCoefficientChains_adm11}; (C12) is
\lref{GlobalTheory.T.blochKato_dimFormula_all_localSituations}; (C13) is
\lref{GlobalTheory.T.ramifiedCoefficientChains_adm13}; (C14) consists of two clauses of Herr's theorem for every record
(\lref{GlobalTheory.HerrFull.herr_part_a_every_record}, which states the injectivity and bijectivity properties of the
class maps of the Wach structures); (C15) is \lref{GlobalTheory.T.blochKato_crystalline_extension_criterion}; (C16) is
\lref{GlobalTheory.HerrExtClass.herr_extension_class_formula}; (C17) is
\lref{Formalism.MainNotions.HerrComplexNaturality.herr_naturality_theorem_holds} (the naturality of the Herr complex,
from Herr's theorem for every record); (C18) is
\lref{GlobalTheory.T.gamma_sub_one_div_uniformiser_at_letter_of_conclusions}, applied to the norm bound for
$\Delta$-invariant lattices; and (C19) is \lref{GlobalTheory.T.h1fGeneratorBoundAll_loc}.
\end{proof}

\begin{definition}[The letter condition]\label{pf:linkN-adm:def-letter}
Let $W$ be a Wach structure of residue characteristic $p$. The \emph{letter condition} at $W$ is: if $W$ is in case~A,
then for every crystalline lattice $T$, letter $i$, integer $y\ge1$, $\OO_C$-linear map $\delta\colon D(T)\to D(T)$ with
$\pi_0\delta(z)=\gamma z-z$ for all $z$, map $c\colon\Z\to\OO_C$ with $-p\,c(x)=\chi(\gamma_0)^{-x}-1$ for all $x$, and finite
set $X\subseteq\Z$ containing $a+j$ for every jump $a$ of $T$ at an embedding over $i$ and every $1\le j\le y$: if $z\in D(T)$
satisfies $q^yz\in N_i(T)$, then $\prod_{x\in X}(\delta-c(x))\,z\in N_i(T)$. The \emph{global letter condition} is the
letter condition at the Wach structures of all places $(\ell,K,s,L_C)$.
\formal{Formalism.MainNotions.N.LetterConditionLocal}
\formal{Formalism.MainNotions.N.LetterConditionGlobal}
\end{definition}

\begin{lemma}[The letter condition at one place]\label{pf:linkN-adm:letterloc}
Let $W$ be a Wach structure of residue characteristic $p=\ell$. Assume that in case~A every crystalline lattice $T$
satisfies the ring hypotheses of (P9), and the following \emph{shift statement}: in case~A, for every crystalline $T$, letter
$i$, integers $y'\le Y$, map $c$ with $\ell\,c(x)=1-\chi(\gamma_0)^{-x}$, list $x_0,\dots,x_{t-1}$ of integers containing
$a+j$ for every jump $a$ of $T$ over $i$ and every $y'<j\le Y$, and sequence $(z_k)$ in $D(T)$ such that $z_0$ is fixed by
$e_i$ and by $\Delta$, $q^Yz_0\in N(T)$ and $\pi_0(z_{k+1}+c(x_k)z_k)=\gamma z_k-z_k$ for $k<t$, the element $z_t$ is fixed by
$e_i$ and by $\Delta$ and satisfies $q^{y'}z_t\in N(T)$. Then the letter condition holds at $W$.
\formallabel{GlobalTheory.PrimeSequence.eigenvalue_product_maps_letter_lattice_local}
\end{lemma}

\begin{proof}
Let $T,i,y,\delta,c,X,z$ be as in the letter condition, with $q^yz\in N_i(T)$. By the ring hypotheses, this implies that $z$
is fixed by $e_i$ and by $\Delta$ and that $q^yz\in N(T)$
(\lref{GlobalTheory.PrimeSequence.dl_of_qp}). Since $p=\ell$, the map $c$ satisfies $\ell\,c(x)=1-\chi(\gamma_0)^{-x}$.
List $X$ as $x_0,\dots,x_{t-1}$ and put $z_k=\prod_{s<k}(\delta-c(x_s))\,z$. Then $z_0=z$ and $z_{k+1}=\delta z_k-c(x_k)z_k$,
so $\pi_0(z_{k+1}+c(x_k)z_k)=\pi_0\delta z_k=\gamma z_k-z_k$. The shift statement with $y'=0$ and $Y=y$ shows that
$z_t=\prod_{x\in X}(\delta-c(x))\,z$ is fixed by $e_i$ and by $\Delta$ and lies in $N(T)$; hence it lies in $N_i(T)$.
\end{proof}

\begin{proposition}[The global letter condition]\label{pf:linkN-adm:letter}
Let $D$ be a datum with local Hodge--Tate, Fontaine and comparison data. The norm bound for $\Delta$-invariant lattices
implies the global letter condition.
\formallabel{GlobalTheory.T.letterConditionGlobal_of_N1used}
\formallabel{GlobalTheory.PrimeSequence.eigenvalue_product_maps_letter_lattice_of_conclusions}
\end{proposition}

\begin{proof}
Let $(\ell,K,s,L_C)$ be a place with Wach structure $W$ in case~A (otherwise there is nothing to prove). Then the level is
$s=1$, $K$ is unramified, $\ell\ne2$, and the residue characteristic of $W$ is $\ell$. The clause (P9) of the norm bound at
$W$ gives the ring hypotheses for every crystalline $T$
(\lref{Formalism.MainNotions.RamifiedCoefficientChains.Sh.ZT.ring_hypotheses_of_crystalline}), and the shift statement is
\lref{GlobalTheory.T.wachOperatorProduct_shift_mem_lattice}, which proves it from the norm bound at every place in
case~A. Apply \cref{pf:linkN-adm:letterloc}.
\end{proof}

% ---------------------------------------------------------------------------
\subsection{Remarks on the formal statements}\label{pf:linkN-adm:remarks}

\begin{remark}\label{pf:linkN-adm:rem}
(1) The main theorem is stated for arbitrary local Hodge--Tate, Fontaine and comparison data and an arbitrary parameter
choice; the Wach structures of \cref{pf:linkN-adm:wach} and the period data enter only through the axioms of these
structures. The parameter choice is not used in the bound for the error term (\cref{pf:linkN-adm:bbmain}), and the row~B
bound (\cref{pf:linkN-adm:rowB}) does not use its hypothesis $S_c+\eta_{\mathrm{sz}}+t+2\le b_L$.

(2) The constant $7200$ of \eqref{pf:linkN-adm:eq-bb} is $60\cdot120$, and \eqref{pf:linkN-adm:eq-bb} is proved with $\Lambda$
in place of $\Lambda^2$ and then weakened using $\Lambda\ge1$.

(3) The row lengths $\linkNadmlen_B$ and $\linkNadmlen_{CD}$ are rational numbers chosen strictly between two real
bounds; if the interval is empty the value $0$ is used. For $\linkNadmlen_{CD}$ the interval is never empty (since $r\ge5$);
for $\linkNadmlen_B$ it is empty exactly when $\mathrm{core}_B\le0$, and the proof of \cref{pf:linkN-adm:rowB} covers this
case.

(4) In theorem~B the band indicator is formed with the same constant $C$ as the bound; it is used with $C=320$, which is
the constant $C_B$ of the band predicate of the hub data.

(5) Theorem~C has no hypothesis that $K$ is unramified, but its setting contains case~A ($K=K_0$, level $1$,
$\ell\ne2$); its conclusion is used only in row~D, at unramified places.

(6) At an inert prime the declarations of the hub data have no columns, so the lift condition (Lift) of the inert
admissibility is vacuous; all places above $\ell$ are rows.

(7) The proof shows that the inert primes of the data considered here are at least $11$, because the split bound is at
least $60$ (\cref{pf:linkN-adm:eleven}).

(8) Some formal statements carry hypotheses that their proofs do not use: the arc setup (\cref{pf:linkN-adm:arcsetup})
assumes the exactness of $D_{\mathrm{cris}}$ (dc), and the arc length bound (\cref{pf:linkN-adm:arclen}) assumes the
universal statement (nn) and that $F$ is onto. In the arc datum of a word every letter has type up or down, so the
``thick'' types of the arc data (types other than thin) are not the thick letters of the word; for instance
$\mathrm{ro}(c)=\mathrm{len}(c)$ for these arc data.
\end{remark}

